\documentclass[12pt]{article}
\usepackage[T5]{fontenc}
\usepackage[utf8]{inputenc}
\usepackage[vietnamese]{babel}
\shorthandoff{"}

% Cấu hình Font chữ và Cỡ chữ 13pt tiêu chuẩn
\usepackage[fontsize=13pt]{scrextend}
\usepackage{mathptmx}

% Cấu hình lề trang theo quy chuẩn kỹ thuật
\usepackage[top=2.5cm,bottom=2.5cm,left=3.5cm,right=2.0cm]{geometry}

% Cấu hình giãn dòng và thụt lề đoạn văn
\usepackage{setspace}
\setstretch{1.5}
\usepackage{indentfirst}
\setlength{\parindent}{1.0cm}
\setlength{\parskip}{4pt}

% Các gói đồ họa, bảng biểu và toán học
\usepackage{graphicx}
\usepackage{xcolor}
\usepackage{tikz}
\usetikzlibrary{calc,shapes.geometric,arrows,arrows.meta,positioning,fit,trees}
\usepackage{booktabs}
\usepackage{array}
\usepackage{tabularx}
\usepackage{longtable}
\usepackage{multirow}
\usepackage{amsmath}
\usepackage{amssymb}
\usepackage{caption}
\captionsetup[table]{labelfont=bf,textfont=bf,position=above,skip=6pt}
\captionsetup[figure]{labelfont=bf,textfont=bf,position=below,skip=6pt}
\usepackage{url}
\usepackage[hidelinks]{hyperref}

% Cấu hình Header, Footer và Đánh số trang góc dưới bên phải
\usepackage{fancyhdr}
\pagestyle{fancy}
\fancyhf{}
\fancyfoot[R]{\thepage}
\renewcommand{\headrulewidth}{0pt}
\renewcommand{\footrulewidth}{0pt}

% Cấu hình trích dẫn chuẩn APA bằng natbib
\usepackage[round]{natbib}
\newcommand{\parencite}{\citep}

% Định dạng cột tự canh lề và xuống dòng cho tabularx
\newcolumntype{L}[1]{>{\raggedright\arraybackslash}p{#1}}
\newcolumntype{C}[1]{>{\centering\arraybackslash}p{#1}}
\newcolumntype{Y}{>{\raggedright\arraybackslash}X}
\newcolumntype{Z}{>{\centering\arraybackslash}X}

% Cấu hình định dạng tiêu đề
\usepackage{titlesec}
\titleformat{\section}{\fontsize{14}{17}\selectfont\bfseries}{}{0em}{}
\titleformat{\subsection}{\fontsize{13}{16}\selectfont\bfseries}{}{0em}{}
\titleformat{\subsubsection}{\fontsize{13}{16}\selectfont\bfseries\itshape}{}{0em}{}
\titlespacing*{\section}{0pt}{16pt}{7pt}
\titlespacing*{\subsection}{0pt}{12pt}{5pt}
\titlespacing*{\subsubsection}{0pt}{8pt}{3pt}

% Cấu hình mục lục có dấu chấm nối cho tất cả các cấp
\usepackage{tocloft}
\renewcommand{\cftsecleader}{\cftdotfill{\cftdotsep}}
\setlength{\cftsecindent}{0pt}
\setlength{\cftsubsecindent}{15pt}
\setlength{\cftsubsubsecindent}{30pt}

\begin{document}

% 1. Trang bìa chuẩn của Đại học Bách khoa Hà Nội
\begin{titlepage}
    % Khung viền đơn mảnh bám theo lề giấy chuẩn của Đại học Bách khoa Hà Nội
    \begin{tikzpicture}[overlay,remember picture]
        \draw [line width=1.5pt] ($ (current page.north west) + (3cm,-2cm) $) rectangle ($ (current page.south east) + (-2cm,1.5cm) $);
    \end{tikzpicture}
    
    \centering
    
    % Tên Trường & Khoa chủ quản
    {\fontsize{14}{16}\selectfont\textbf{ĐẠI HỌC BÁCH KHOA HÀ NỘI}\\} 
    {\fontsize{12}{14}\selectfont\textbf{KHOA KHOA HỌC VÀ CÔNG NGHỆ GIÁO DỤC}\\} 
    
    \vspace{1.5cm}
    
    % Khối định vị LOGO Trường (tương thích khi có hoặc chưa có tệp ảnh)
    \IfFileExists{hust_logo.png}{%
        \includegraphics[height=4.5cm]{hust_logo.png}%
    }{%
        \fbox{\begin{minipage}[c][4.5cm][c]{4.5cm}\centering\textbf{\color{red!80!black}HUST LOGO}\\\vspace{4pt}\fontsize{9}{11}\selectfont(Đặt file hust\_logo.png tại đây)\end{minipage}}%
    }
    
    \vspace{1.5cm}
    
    % Loại hình tài liệu (Màu đỏ, in hoa đậm)
    {\fontsize{16}{18}\selectfont\textbf{\color{red}BÁO CÁO HỌC PHẦN}\\} 
    
    \vspace{0.8cm}
    
    % Tên đề tài báo cáo
    {\fontsize{14}{18}\selectfont\textbf{DỰ ÁN XÂY DỰNG BỘ TÀI LIỆU TƯ VẤN PHÂN LUỒNG\\ VÀ ĐỊNH HƯỚNG NGHỀ NGHIỆP CHO HỌC SINH\\ CUỐI CẤP TRUNG HỌC CƠ SỞ}\\} 
    
    \vspace{0.4cm}
    \noindent\rule{0.65\textwidth}{0.6pt} % Đường kẻ ngang phân cách tiêu đề
    \vspace{0.6cm}
    
    % Tên Học phần tương ứng
    {\fontsize{12}{14}\selectfont\textit{\textbf{Học phần:} Quản lý dự án giáo dục}\\} 
    
    \vspace{2.0cm}
    
    % Khối thông tin Giảng viên và Học viên thực hiện
    \begin{flushleft}
        \hspace{2.5cm} 
        \begin{tabular}{ll}
            \fontsize{12}{14}\selectfont\textbf{Giảng viên giảng dạy:} & \fontsize{12}{14}\selectfont PGS. TS. Bùi Thị Thúy Hằng \\
            \fontsize{12}{14}\selectfont\textbf{Học viên thực hiện:} & \fontsize{12}{14}\selectfont Lê Hạnh Nguyên -- 20252233M
        \end{tabular}
    \end{flushleft}
    
    \vfill
    
    % Địa điểm và thời gian xuất bản ở cuối trang
    {\fontsize{12}{14}\selectfont\textbf{Hà Nội -- 2026}\\}
    \vspace*{0.4cm}
    
\end{titlepage}



\newpage
% 2. Trang mục lục (không đánh số trang)
\pagenumbering{gobble}
\renewcommand{\contentsname}{\centering\fontsize{14}{17}\selectfont\textbf{MỤC LỤC}}
{\setstretch{1.2}
\tableofcontents
}
\newpage

% Bắt đầu đánh số trang Ả Rập từ Phần Mở đầu
\pagenumbering{arabic}
\setcounter{page}{1}

% 3. Phần Mở đầu
\section*{PHẦN MỞ ĐẦU}
\addcontentsline{toc}{section}{PHẦN MỞ ĐẦU}

\subsection*{1. Tính cấp thiết của dự án}
\addcontentsline{toc}{subsection}{1. Tính cấp thiết của dự án}

\subsubsection*{1.1. Căn cứ chính sách và yêu cầu phân luồng giáo dục quốc gia}

Phân luồng và giáo dục hướng nghiệp là chủ trương lớn nhằm điều chỉnh cơ cấu nguồn nhân lực tại Việt Nam. Tại Điều 9 Luật Giáo dục số 43/2019/QH14 (được sửa đổi, bổ sung bởi Luật số 123/2025/QH15), hướng nghiệp trong giáo dục được xác định là hệ thống các biện pháp tiến hành trong và ngoài cơ sở giáo dục nhằm giúp người học có tri thức về nghề nghiệp, hình thành khả năng lựa chọn nghề nghiệp trên cơ sở kết hợp nguyện vọng, sở trường cá nhân với nhu cầu lao động của xã hội. Gắn liền với đó, phân luồng trong giáo dục tạo điều kiện để người học sau khi hoàn thành chương trình trung học cơ sở tiếp tục học ở cấp học cao hơn hoặc tham gia học nghề phù hợp với năng lực và điều kiện cụ thể.

Cụ thể hóa định hướng này, Quyết định số 522/QĐ-TTg ngày 14 tháng 5 năm 2018 của Thủ tướng Chính phủ phê duyệt Đề án Giáo dục hướng nghiệp và định hướng phân luồng học sinh trong giáo dục phổ thông giai đoạn 2018 -- 2025 đã đặt mục tiêu có ít nhất 40\% học sinh tốt nghiệp trung học cơ sở tiếp tục học tập tại các cơ sở giáo dục nghề nghiệp đào tạo trình độ sơ cấp, trung cấp; 100\% trường trung học cơ sở thực hiện chương trình giáo dục hướng nghiệp và ít nhất 70\% học sinh được tiếp cận thông tin thị trường lao động. Nhiệm vụ trọng tâm thứ hai của Đề án yêu cầu biên soạn và cung cấp tài liệu hướng nghiệp chuẩn hóa, phù hợp với tâm lý lứa tuổi và điều kiện thực tiễn.

Chủ trương đẩy mạnh phân luồng tiếp tục được nhấn mạnh tại Kết luận số 91-KL/TW ngày 12 tháng 8 năm 2024 của Bộ Chính trị về tiếp tục thực hiện Nghị quyết số 29-NQ/TW. Về cơ chế tài chính, Nghị định số 81/2021/NĐ-CP và Nghị định số 97/2023/NĐ-CP của Chính phủ đã quy định miễn 100\% học phí đối với người tốt nghiệp trung học cơ sở đi học trung cấp nghề, tạo hành lang thuận lợi cho học sinh theo học mô hình đào tạo nghề tích hợp văn hóa hệ chín cộng.

\subsubsection*{1.2. Khoảng cách thực tiễn tại các trường trung học cơ sở}

Mặc dù khung pháp lý và chính sách miễn học phí học nghề đã được ban hành, kết quả thực hiện phân luồng sau trung học cơ sở trên thực tế vẫn còn khoảng cách lớn so với mục tiêu đề ra. Số liệu tổng kết của Bộ Giáo dục và Đào tạo cùng Tổng cục Giáo dục nghề nghiệp cho thấy tỷ lệ học sinh tốt nghiệp trung học cơ sở chuyển sang học nghề trên toàn quốc mới đạt khoảng 15\% đến 18\%, thấp hơn nhiều so với chỉ tiêu 40\% của Quyết định số 522/QĐ-TTg.

Nguyên nhân trước hết đến từ áp lực thi tuyển vào lớp 10 trung học phổ thông công lập. Phần lớn học sinh và phụ huynh vẫn coi trung học phổ thông công lập là lựa chọn duy nhất, dẫn đến sự cạnh tranh gay gắt trong kỳ thi tuyển sinh đầu cấp. Nghiên cứu của Nguyễn Thị Mai Hương (2022) chỉ ra rằng việc thiếu thông tin về các lộ trình học tập thay thế có mối liên hệ trực tiếp với mức độ căng thẳng, lo âu thi cử ở học sinh độ tuổi 14 đến 15. Bên cạnh đó, tổng kết của Đặng Bá Lãm (2019) cho thấy tâm lý chuộng bằng cấp khiến nhiều gia đình nhìn nhận việc học nghề sau lớp 9 như một giải pháp bất đắc dĩ, đồng thời chưa nắm rõ quyền lợi miễn 100\% học phí cũng như cơ hội học liên thông lên cao đẳng, đại học.

Về phía nhà trường, nguồn học liệu hướng nghiệp chuyên biệt phục vụ học sinh cuối cấp còn thiếu hụt. Khảo sát của Nguyễn Đức Sơn và cộng sự (2026) tại 24 trường trung học cơ sở ghi nhận 78,4\% học sinh lớp 9 có nhu cầu tiếp cận một bộ tài liệu hướng dẫn cụ thể về các hướng đi sau tốt nghiệp. Trong khi đó, giáo viên phụ trách tư vấn hướng nghiệp tại trường phổ thông chủ yếu làm việc kiêm nhiệm, thiếu công cụ trắc nghiệm sở thích nghề nghiệp đã được chuẩn hóa để tư vấn cho học sinh và phụ huynh.

Từ yêu cầu của chính sách và khoảng trống học liệu tại cơ sở, Dự án xây dựng bộ tài liệu tư vấn phân luồng và định hướng nghề nghiệp cho học sinh cuối cấp trung học cơ sở được thiết lập nhằm cung cấp một ấn phẩm hướng dẫn khoa học, hỗ trợ trực tiếp cho học sinh lớp 9 và gia đình trước kỳ chuyển cấp.

\subsection*{2. Phạm vi triển khai dự án}
\addcontentsline{toc}{subsection}{2. Phạm vi triển khai dự án}

Dự án tập trung thực hiện trọn vẹn chu trình phát triển và xuất bản học liệu hướng nghiệp qua năm giai đoạn kế tiếp nhau: khảo sát nhu cầu thực tế trên 350 bộ phiếu và thu thập dữ liệu hướng nghiệp tại năm trường trung học cơ sở đại diện; xây dựng khung cấu trúc bốn phần và biên soạn bản thảo 80 trang ruột; tổ chức phản biện độc lập, lấy ý kiến đọc thử của học sinh, giáo viên, họp Hội đồng thẩm định chuyên môn và hiệu đính bản thảo sạch; thiết kế đồ họa, dàn trang chế bản điện tử, lựa chọn nhà thầu và in ấn công nghiệp 2.000 cuốn tài liệu; cuối cùng là kiểm tra chất lượng theo mẫu thống kê 384 cuốn, nghiệm thu cấp cơ sở và bàn giao toàn bộ ấn phẩm cùng bộ tệp tin chế bản gốc cho năm trường trung học cơ sở thụ hưởng để đưa vào khai thác trước kỳ đăng ký nguyện vọng tuyển sinh lớp 10.

\subsection*{3. Các thông số tổng quan xác lập dự án}
\addcontentsline{toc}{subsection}{3. Các thông số tổng quan xác lập dự án}

\begin{table}[htbp]
\centering
\caption{Các thông số tổng quan xác lập dự án}
\label{tab:thong_so_tong_quan}
\begin{tabularx}{\textwidth}{L{4.2cm}Y}
\toprule
\textbf{Thông số quản lý} & \textbf{Nội dung xác lập} \\
\midrule
Tên dự án & Dự án xây dựng bộ tài liệu tư vấn phân luồng và định hướng nghề nghiệp cho học sinh cuối cấp trung học cơ sở \\
Cơ quan chủ quản & Phòng Giáo dục và Đào tạo phối hợp cụm các trường trung học cơ sở \\
Đơn vị thực hiện & Ban Quản lý dự án phát triển học liệu hướng nghiệp trung học cơ sở \\
Đối tượng thụ hưởng trực tiếp & Học sinh lớp 9 cuối cấp trung học cơ sở và cha mẹ học sinh \\
Sản phẩm bàn giao cốt lõi & 2.000 cuốn tài liệu khổ 17 $\times$ 24 cm (gồm 80 trang ruột giấy Fort 100 gsm và 4 trang bìa ngoài giấy Couche 250 gsm cán màng mờ, in offset 4 màu, đóng keo nhiệt) kèm 01 bộ tệp tin chế bản điện tử gốc \\
Tổng thời gian thực hiện & 90 ngày lịch (13 tuần, từ ngày 02/03/2026 đến ngày 30/05/2026) \\
Tổng mức ngân sách & 185.000.000 đồng (một trăm tám mươi lăm triệu đồng chẵn) \\
Địa bàn triển khai & Cụm 5 trường trung học cơ sở công lập đại diện trên địa bàn \\
Phương thức quản lý & Quản lý theo mục tiêu, cấu trúc phân rã công việc và kiểm soát tiến độ theo đường găng \\
\bottomrule
\end{tabularx}
\end{table}


\newpage

% 4. Chương 1: Bối cảnh thực tiễn và cơ sở xác định dự án giáo dục
\section*{CHƯƠNG 1. BỐI CẢNH THỰC TIỄN VÀ CƠ SỞ XÁC ĐỊNH DỰ ÁN GIÁO DỤC}
\addcontentsline{toc}{section}{CHƯƠNG 1. BỐI CẢNH THỰC TIỄN VÀ CƠ SỞ XÁC ĐỊNH DỰ ÁN GIÁO DỤC}

Chương này vận dụng \textit{Logical Framework Approach} (Phương pháp tiếp cận khung logic) và lý thuyết chu trình dự án giáo dục để phân tích quan hệ nhân quả của bài toán phân luồng tại trường trung học cơ sở, xác lập hệ thống mục tiêu, luận chứng lựa chọn chiến lược can thiệp và định hình bốn giai đoạn vòng đời dự án.

\subsection*{1.1. Cấu trúc vấn đề thực tiễn theo mô hình Cây vấn đề}
\addcontentsline{toc}{subsection}{1.1. Cấu trúc vấn đề thực tiễn theo mô hình Cây vấn đề}

Mô hình \textit{Problem Tree} (Cây vấn đề) được sử dụng để bóc tách các nguyên nhân trực tiếp, nguyên nhân gốc rễ và hệ quả của tình trạng lúng túng trong định hướng phân luồng cho học sinh cuối cấp trung học cơ sở.

\subsubsection*{1.1.1. Vấn đề trung tâm}

Vấn đề trung tâm cần giải quyết là: Học sinh lớp 9 thiếu thông tin hệ thống và công cụ tự đánh giá năng lực, dẫn đến lúng túng và định hướng chưa sát thực tế khi lựa chọn con đường học tập sau tốt nghiệp trung học cơ sở.

\subsubsection*{1.1.2. Hệ thống nguyên nhân và hệ quả}

Vấn đề trung tâm xuất phát từ ba nhánh nguyên nhân chính tại cơ sở giáo dục phổ thông. Nhánh nguyên nhân thứ nhất là sự thiếu hụt một bộ tài liệu cẩm nang phân luồng tổng hợp: thông tin về các hướng đi sau lớp 9, điều kiện dự tuyển và chính sách miễn học phí học nghề hiện nằm rải rác ở các văn bản hành chính hoặc tờ rơi tuyển sinh riêng lẻ, trong khi thời lượng nội dung hướng nghiệp trong sách giáo khoa còn khá khái quát. Nhánh nguyên nhân thứ hai là việc học sinh thiếu công cụ khoa học để tự nhận diện sở thích và năng lực nghề nghiệp phù hợp với lứa tuổi 14 đến 15, dẫn đến xu hướng chọn trường theo cảm tính hoặc theo tâm lý đám đông. Nhánh nguyên nhân thứ ba đến từ rào cản nhận thức của phụ huynh và sự thiếu hụt học liệu hỗ trợ đối thoại giữa nhà trường với gia đình, khiến giáo viên kiêm nhiệm gặp khó khăn khi tư vấn phân luồng.

Nếu không được tháo gỡ kịp thời, vấn đề này kéo theo hai hệ quả trực tiếp là áp lực thi tuyển vào lớp 10 trung học phổ thông công lập tiếp tục quá tải và tỷ lệ phân luồng vào giáo dục nghề nghiệp duy trì ở mức thấp (15\% đến 18\%). Về dài hạn, sự lệch pha trong lựa chọn đường học dẫn đến tình trạng học sinh đuối sức ở bậc trung học phổ thông và gây mất cân đối cơ cấu nhân lực qua đào tạo.

\begin{figure}[htbp]
\centering
\caption{Sơ đồ Cây vấn đề trong công tác phân luồng học sinh trung học cơ sở}
\label{fig:cay_van_de}
\resizebox{0.95\textwidth}{!}{%
\begin{tikzpicture}[
    node distance = 0.7cm and 0.5cm,
    box_hq/.style = {rectangle, rounded corners=3pt, draw=red!70!black, fill=red!6, line width=0.8pt, text width=4.5cm, align=center, inner sep=5pt, font=\footnotesize},
    box_vd/.style = {rectangle, rounded corners=3pt, draw=red!90!black, fill=red!15, line width=1.2pt, text width=12cm, align=center, inner sep=6pt, font=\footnotesize\bfseries},
    box_nn/.style = {rectangle, rounded corners=3pt, draw=blue!70!black, fill=blue!6, line width=0.8pt, text width=4.5cm, align=center, inner sep=5pt, font=\footnotesize},
    box_sub/.style = {rectangle, rounded corners=3pt, draw=gray!70, fill=gray!8, line width=0.6pt, text width=4.2cm, align=center, inner sep=4pt, font=\scriptsize},
    arr/.style = {-{Stealth[scale=0.9]}, line width=0.8pt, draw=gray!70!black}
]

% Vấn đề trung tâm
\node (vd) [box_vd] {HỌC SINH LỚP 9 LÚNG TÚNG, THIẾU THÔNG TIN VÀ ĐỊNH HƯỚNG CHƯA SÁT NĂNG LỰC KHI PHÂN LUỒNG SAU TRUNG HỌC CƠ SỞ};

% Hệ quả phía trên
\node (hq1) [box_hq, above left=1.2cm and -0.5cm of vd] {Áp lực thi vào lớp 10 quá tải, căng thẳng tâm lý ở học sinh lớp 9};
\node (hq2) [box_hq, above right=1.2cm and -0.5cm of vd] {Tỷ lệ phân luồng sau THCS vào học nghề thấp (15--18\% so với chỉ tiêu 40\%)};
\node (hq3) [box_hq, text width=10cm, above=0.7cm of $(hq1.north)!0.5!(hq2.north)$] {Cơ cấu lao động mất cân đối, lãng phí nguồn lực xã hội và gia đình};

\draw [arr] (vd.north) -- ++(0,0.3) -| (hq1.south);
\draw [arr] (vd.north) -- ++(0,0.3) -| (hq2.south);
\draw [arr] (hq1.north) -- ++(0,0.2) -| (hq3.south);
\draw [arr] (hq2.north) -- ++(0,0.2) -| (hq3.south);

% Nguyên nhân phía dưới
\node (nn2) [box_nn, below=1.0cm of vd] {Thiếu công cụ trắc nghiệm sở thích, năng lực phù hợp lứa tuổi 14--15};
\node (nn1) [box_nn, left=0.6cm of nn2] {Thiếu bộ tài liệu cẩm nang phân luồng chuẩn hóa, dễ tra cứu};
\node (nn3) [box_nn, right=0.6cm of nn2] {Tâm lý chuộng bằng cấp của phụ huynh; giáo viên thiếu học liệu tư vấn};

\draw [arr] (nn1.north) -- ++(0,0.3) -| (vd.south);
\draw [arr] (nn2.north) -- (vd.south);
\draw [arr] (nn3.north) -- ++(0,0.3) -| (vd.south);

% Nguyên nhân gốc
\node (nn1_1) [box_sub, below=0.6cm of nn1] {Thông tin các hướng đi sau lớp 9 và chính sách học phí bị phân mảnh};
\node (nn2_1) [box_sub, below=0.6cm of nn2] {Học sinh chọn trường theo cảm tính, trào lưu đám đông};
\node (nn3_1) [box_sub, below=0.6cm of nn3] {Chưa có tài liệu làm cầu nối đối thoại giữa học sinh và cha mẹ};

\draw [arr] (nn1_1.north) -- (nn1.south);
\draw [arr] (nn2_1.north) -- (nn2.south);
\draw [arr] (nn3_1.north) -- (nn3.south);

\end{tikzpicture}%
}
\end{figure}

\subsection*{1.2. Hệ thống mục tiêu dự án theo mô hình Cây mục tiêu}
\addcontentsline{toc}{subsection}{1.2. Hệ thống mục tiêu dự án theo mô hình Cây mục tiêu}

Từ cấu trúc Cây vấn đề, mô hình \textit{Objective Tree} (Cây mục tiêu) chuyển hóa các trạng thái hạn chế hiện tại thành hệ thống mục tiêu kỳ vọng theo ba cấp độ: mục tiêu phát triển cấp cao, mục tiêu trực tiếp của dự án và các kết quả đầu ra cụ thể.

Ở cấp độ phát triển vĩ mô, dự án góp phần nâng cao hiệu quả phân luồng học sinh sau trung học cơ sở theo định hướng của Quyết định số 522/QĐ-TTg và giảm tải áp lực tâm lý trong kỳ thi tuyển sinh vào lớp 10. Ở cấp độ mục tiêu trực tiếp, dự án trang bị cho học sinh lớp 9 công cụ tự đánh giá bản thân và hệ thống thông tin minh bạch về các lộ trình học tập để chủ động xây dựng kế hoạch chọn trường cùng gia đình.

Để đạt được mục tiêu trực tiếp đó, dự án xác lập bốn kết quả chuyển giao cụ thể: hoàn thành bản thảo bộ tài liệu tư vấn phân luồng có tích hợp trắc nghiệm sở thích nghề nghiệp theo mô hình mật mã Holland rút gọn và bản đồ các hướng đi sau lớp 9; thẩm định chuyên môn bản thảo đáp ứng yêu cầu cần đạt của Chương trình Giáo dục phổ thông 2018 (Thông tư số 32/2018/TT-BGDĐT); xuất bản và bàn giao 2.000 cuốn ấn phẩm đạt chuẩn kỹ thuật cho năm trường trung học cơ sở đại diện; đồng thời cung cấp bộ phiếu đối thoại hướng nghiệp giúp gắn kết học sinh, giáo viên tư vấn và phụ huynh.

\begin{figure}[htbp]
\centering
\caption{Sơ đồ Cây mục tiêu của dự án}
\label{fig:cay_muc_tieu}
\resizebox{0.95\textwidth}{!}{%
\begin{tikzpicture}[
    node distance = 0.7cm and 0.5cm,
    box_mc/.style = {rectangle, rounded corners=3pt, draw=green!60!black, fill=green!6, line width=0.8pt, text width=4.5cm, align=center, inner sep=5pt, font=\footnotesize},
    box_ct/.style = {rectangle, rounded corners=3pt, draw=blue!80!black, fill=blue!10, line width=1.2pt, text width=12cm, align=center, inner sep=6pt, font=\footnotesize\bfseries},
    box_kq/.style = {rectangle, rounded corners=3pt, draw=blue!70!black, fill=blue!5, line width=0.8pt, text width=4.5cm, align=center, inner sep=5pt, font=\footnotesize},
    box_hd/.style = {rectangle, rounded corners=3pt, draw=gray!70, fill=gray!8, line width=0.6pt, text width=4.2cm, align=center, inner sep=4pt, font=\scriptsize},
    arr/.style = {-{Stealth[scale=0.9]}, line width=0.8pt, draw=gray!70!black}
]

% Mục tiêu trực tiếp
\node (ct) [box_ct] {HỌC SINH LỚP 9 ĐƯỢC TRANG BỊ CÔNG CỤ TỰ ĐÁNH GIÁ VÀ ĐẦY ĐỦ THÔNG TIN ĐỂ CHỦ ĐỘNG ĐỊNH HƯỚNG PHÂN LUỒNG};

% Mục tiêu cấp cao
\node (mc1) [box_mc, above left=1.2cm and -0.5cm of ct] {Giảm áp lực thi cử lớp 10, học sinh tự tin bước vào giai đoạn chuyển cấp};
\node (mc2) [box_mc, above right=1.2cm and -0.5cm of ct] {Nâng cao tỷ lệ phân luồng học sinh sau THCS vào các hướng đi phù hợp};
\node (mc3) [box_mc, text width=10cm, above=0.7cm of $(mc1.north)!0.5!(mc2.north)$] {Cơ cấu nhân lực cân bằng, sử dụng hiệu quả nguồn lực đào tạo xã hội và gia đình};

\draw [arr] (ct.north) -- ++(0,0.3) -| (mc1.south);
\draw [arr] (ct.north) -- ++(0,0.3) -| (mc2.south);
\draw [arr] (mc1.north) -- ++(0,0.2) -| (mc3.south);
\draw [arr] (mc2.north) -- ++(0,0.2) -| (mc3.south);

% Kết quả đầu ra
\node (kq2) [box_kq, below=1.0cm of ct] {Chuẩn hóa bài trắc nghiệm sở thích nghề nghiệp dành cho lứa tuổi 14--15};
\node (kq1) [box_kq, left=0.6cm of kq2] {Biên soạn hoàn chỉnh Bộ tài liệu tư vấn phân luồng trực quan, dễ tra cứu};
\node (kq3) [box_kq, right=0.6cm of kq2] {Hệ thống hóa bản đồ 4 hướng đi sau lớp 9, chính sách học phí và trường nghề};

\draw [arr] (kq1.north) -- ++(0,0.3) -| (ct.south);
\draw [arr] (kq2.north) -- (ct.south);
\draw [arr] (kq3.north) -- ++(0,0.3) -| (ct.south);

% Hoạt động
\node (hd1) [box_hd, below=0.6cm of kq1] {Khảo sát nhu cầu, biên soạn và thẩm định chuyên môn bản thảo};
\node (hd2) [box_hd, below=0.6cm of kq2] {Thiết kế đồ họa, chế bản điện tử và in ấn offset 2.000 cuốn tài liệu};
\node (hd3) [box_hd, below=0.6cm of kq3] {Kiểm định chất lượng mẫu thống kê và nghiệm thu bàn giao 5 trường};

\draw [arr] (hd1.north) -- (kq1.south);
\draw [arr] (hd2.north) -- (kq2.south);
\draw [arr] (hd3.north) -- (kq3.south);

\end{tikzpicture}%
}
\end{figure}

\subsection*{1.3. Cơ sở luận chứng và lựa chọn chiến lược can thiệp}
\addcontentsline{toc}{subsection}{1.3. Cơ sở luận chứng và lựa chọn chiến lược can thiệp}

Nhằm lựa chọn phương thức can thiệp phù hợp nhất với nguồn lực 185.000.000 đồng và thời gian 90 ngày của địa phương, nhóm lập kế hoạch đặt lên bàn cân ba phương án chiến lược:

Phương án 1 (Tổ chức chuỗi sự kiện và ngày hội tư vấn phân luồng trực tiếp): Tạo được hiệu ứng truyền thông tại thời điểm diễn ra sự kiện nhưng chi phí thuê mặt bằng, dàn dựng sân khấu cao; thời gian tiếp xúc ngắn (chỉ một đến hai buổi) nên học sinh khó tra cứu sâu thông tin và không để lại học liệu lưu giữ lâu dài.

Phương án 2 (Xây dựng phần mềm hoặc cổng thông tin hướng nghiệp trực tuyến): Có lợi thế về khả năng truy cập số nhưng đòi hỏi chi phí lập trình, thuê máy chủ và bảo trì hàng năm vượt quá ngân sách sự nghiệp ngắn hạn; đồng thời tiềm ẩn rủi ro chậm tiến độ do lỗi kỹ thuật phần mềm.

Phương án 3 (Biên soạn, thẩm định và in ấn bộ tài liệu chuẩn hóa): Tập trung ngân sách vào chất lượng nội dung sư phạm và tạo ra ấn phẩm vật lý cụ thể. Học sinh có thể trực tiếp làm bài trắc nghiệm trên giấy, mang tài liệu về nhà cùng cha mẹ đối chiếu các phương án chọn trường, trong khi cơ quan quản lý dễ dàng kiểm soát định mức chi phí, tiến độ và nghiệm thu sản phẩm đầu ra.

\begin{table}[htbp]
\centering
\caption{Ma trận đánh giá và lựa chọn chiến lược can thiệp (Thang điểm 10)}
\label{tab:lua_chon_chien_luoc}
\begin{tabularx}{\textwidth}{L{5.0cm}CY{2.2cm}Y{2.2cm}Y{2.2cm}}
\toprule
\textbf{Tiêu chí đánh giá} & \textbf{Trọng số} & \textbf{Phương án 1: Tổ chức sự kiện} & \textbf{Phương án 2: Xây phần mềm} & \textbf{Phương án 3: Biên soạn ấn phẩm} \\
\midrule
1. Tính khả thi kỹ thuật và quản lý & 20\% & 6 (Điểm: 1,20) & 5 (Điểm: 1,00) & 9 (Điểm: 1,80) \\
2. Tính kinh tế và tối ưu ngân sách & 25\% & 5 (Điểm: 1,25) & 4 (Điểm: 1,00) & 9 (Điểm: 2,25) \\
3. Mức độ tác động và tính bền vững & 25\% & 4 (Điểm: 1,00) & 7 (Điểm: 1,75) & 8 (Điểm: 2,00) \\
4. Khả năng tiếp cận của người học & 20\% & 5 (Điểm: 1,00) & 6 (Điểm: 1,20) & 8 (Điểm: 1,60) \\
5. Tính kiểm soát chất lượng & 10\% & 6 (Điểm: 0,60) & 5 (Điểm: 0,50) & 9 (Điểm: 0,90) \\
\midrule
\textbf{Tổng điểm có trọng số} & \textbf{100\%} & \textbf{5,05} & \textbf{5,45} & \textbf{8,55} \\
\bottomrule
\end{tabularx}
\end{table}

Kết quả phân tích đa tiêu chí tại Bảng \ref{tab:lua_chon_chien_luoc} cho thấy Phương án 3 đạt tổng điểm có trọng số cao nhất là \textbf{8,55/10 điểm} (so với 5,05 điểm của Phương án 1 và 5,45 điểm của Phương án 2). Do đó, chiến lược biên soạn và xuất bản bộ tài liệu in ấn được lựa chọn để triển khai.

\subsection*{1.4. Đặc trưng dự án và chu trình vòng đời bốn giai đoạn}
\addcontentsline{toc}{subsection}{1.4. Đặc trưng dự án và chu trình vòng đời bốn giai đoạn}

\subsubsection*{1.4.1. Các đặc trưng cơ bản của dự án}

Dự án thể hiện rõ bảy đặc trưng của một dự án giáo dục: có mục tiêu xác định nhằm giải quyết khoảng trống học liệu phân luồng cấp trung học cơ sở; có sản phẩm bàn giao định lượng rõ ràng (2.000 cuốn ấn phẩm và bộ tệp tin chế bản gốc); bị giới hạn chặt chẽ trong khung thời gian 90 ngày (từ ngày 02/03/2026 đến ngày 30/05/2026) và ngân sách 185.000.000 đồng; được vận hành bởi một Ban Quản lý dự án thành lập tạm thời; đòi hỏi sự phối hợp liên ngành giữa nhà khoa học giáo dục, chuyên gia tâm lý, cán bộ quản lý và doanh nghiệp in ấn; liên quan đến nhiều nhóm chủ thể từ Phòng Giáo dục và Đào tạo đến các nhà trường, phụ huynh; đồng thời chịu tác động của các yếu tố bất định về tiến độ biên soạn và kỹ thuật xuất bản.

\subsubsection*{1.4.2. Khung chu trình bốn giai đoạn vòng đời dự án}

Vòng đời 90 ngày của dự án được tổ chức thành bốn giai đoạn nối tiếp, khớp với mười công việc từ A đến J trên sơ đồ mạng tiến độ:

\begin{figure}[htbp]
\centering
\caption{Khung chu trình bốn giai đoạn vòng đời của dự án}
\label{fig:chu_trinh_vong_doi}
\resizebox{0.95\textwidth}{!}{%
\begin{tikzpicture}[
    node distance = 0.5cm,
    giai_doan/.style = {rectangle, rounded corners=4pt, draw=blue!80!black, fill=blue!6, line width=1.0pt, text width=3.3cm, align=center, inner sep=6pt, font=\footnotesize},
    arr/.style = {-{Stealth[scale=1.1]}, line width=1.2pt, draw=blue!70!black}
]
\node (g1) [giai_doan] {\textbf{GIAI ĐOẠN 1}\\XÂY DỰNG Ý TƯỞNG\\(Ngày 1 -- Ngày 8)};
\node (g2) [giai_doan, right=of g1] {\textbf{GIAI ĐOẠN 2}\\PHÁT TRIỂN KẾ HOẠCH\\(Ngày 9 -- Ngày 14)};
\node (g3) [giai_doan, right=of g2] {\textbf{GIAI ĐOẠN 3}\\THỰC HIỆN DỰ ÁN\\(Ngày 15 -- Ngày 84)};
\node (g4) [giai_doan, right=of g3] {\textbf{GIAI ĐOẠN 4}\\KẾT THÚC DỰ ÁN\\(Ngày 85 -- Ngày 90)};

\draw [arr] (g1) -- (g2);
\draw [arr] (g2) -- (g3);
\draw [arr] (g3) -- (g4);
\end{tikzpicture}%
}
\end{figure}

Giai đoạn 1 (Xây dựng ý tưởng, từ ngày 1 đến ngày 8 -- công việc A): Triển khai khảo sát nhu cầu thực tế trên 350 bộ phiếu tại năm trường trung học cơ sở đại diện, thu thập dữ liệu tuyển sinh và hoàn thành báo cáo khảo sát đầu vào.

Giai đoạn 2 (Phát triển kế hoạch, từ ngày 9 đến ngày 14 -- công việc B): Thiết kế khung cấu trúc bốn phần của tài liệu, tổ chức hội thảo thông qua đề cương chi tiết và phê duyệt kế hoạch quản lý tổng thể, cấu trúc phân rã công việc, tiến độ đường găng cùng dự toán ngân sách.

Giai đoạn 3 (Thực hiện dự án, từ ngày 15 đến ngày 84 -- các công việc từ C đến I): Tổ chức cho bốn nhóm tác giả biên soạn bản thảo (C), đồng thời lựa chọn nhà thầu in ấn (D); tiến hành thẩm định chuyên môn độc lập hai vòng (E) và hiệu đính bản thảo sạch (F); thiết kế đồ họa, dàn trang chế bản điện tử (G), in thử bản bông màu để ký duyệt mẫu (H) và in ấn offset hàng loạt 2.000 cuốn tài liệu (I).

Giai đoạn 4 (Kết thúc dự án, từ ngày 85 đến ngày 90 -- công việc J): Rút mẫu thống kê 384 cuốn để kiểm định chất lượng tại kho, họp Hội đồng nghiệm thu cấp cơ sở, vận chuyển bàn giao 2.000 cuốn tài liệu cho năm trường trung học cơ sở, thanh lý hợp đồng, đóng gói 03 bộ hồ sơ lưu trữ và giải thể Ban Quản lý dự án.


\newpage

% 5. Chương 2: Kế hoạch quản lý tổng thể và khởi động dự án
\section*{CHƯƠNG 2. KẾ HOẠCH QUẢN LÝ TỔNG THỂ VÀ KHỞI ĐỘNG DỰ ÁN}
\addcontentsline{toc}{section}{CHƯƠNG 2. KẾ HOẠCH QUẢN LÝ TỔNG THỂ VÀ KHỞI ĐỘNG DỰ ÁN}

Chương này thiết lập văn kiện pháp lý khởi động dự án, phân tích mức độ ảnh hưởng và mối quan tâm của các bên liên quan, đồng thời xác lập quy chế điều hành tích hợp và quy trình kiểm soát thay đổi tổng thể.

\subsection*{2.1. Báo cáo tóm tắt dự án và văn kiện khởi động}
\addcontentsline{toc}{subsection}{2.1. Báo cáo tóm tắt dự án và văn kiện khởi động}

Văn kiện \textit{Project Charter} (Báo cáo tóm tắt dự án) được ban hành kèm theo Quyết định số 128/QĐ-PGDĐT ngày 02 tháng 3 năm 2026 của Trưởng phòng Giáo dục và Đào tạo, chính thức xác lập tư cách pháp lý cho Dự án xây dựng bộ tài liệu tư vấn phân luồng và định hướng nghề nghiệp cho học sinh cuối cấp trung học cơ sở (mã hiệu DA-GDHN-2026) và trao quyền cho Giám đốc dự án huy động nguồn lực thực hiện.

Dự án sử dụng nguồn ngân sách sự nghiệp giáo dục địa phương năm 2026 kết hợp kinh phí chương trình giáo dục hướng nghiệp theo Quyết định số 522/QĐ-TTg với tổng mức đầu tư được duyệt là 185.000.000 đồng, thực hiện trong 90 ngày lịch (từ ngày 02/03/2026 đến ngày 30/05/2026). Phần nội dung học thuật và thiết kế ma-két chế bản điện tử do đội ngũ chuyên gia sư phạm, tâm lý học và họa sĩ thiết kế chuyên ngành đảm nhiệm; phần bình bản kẽm và in ấn công nghiệp 2.000 cuốn được thuê ngoài thông qua hình thức chào hàng cạnh tranh rút gọn. Cơ cấu nhân sự chủ chốt của dự án được xác lập tại Bảng \ref{tab:nhan_su_chu_chot}.

\begin{table}[htbp]
\centering
\caption{Cơ cấu nhân sự chủ chốt trong văn kiện khởi động dự án}
\label{tab:nhan_su_chu_chot}
\begin{tabularx}{\textwidth}{L{3.2cm}L{3.5cm}L{3.5cm}Y}
\toprule
\textbf{Vai trò trong dự án} & \textbf{Họ và tên} & \textbf{Đơn vị công tác và chức danh} & \textbf{Trách nhiệm chính xác lập} \\
\midrule
Đại diện cơ quan phê duyệt & Tiến sĩ Trần Văn Hùng & Trưởng phòng Giáo dục và Đào tạo & Đại diện Chủ đầu tư, phê duyệt chủ trương, ngân sách và ký biên bản nghiệm thu \\
Giám đốc dự án & Thạc sĩ Lê Hoàng Nam & Chuyên viên chính Phòng GD\&ĐT & Điều hành tổng thể, chịu trách nhiệm về tiến độ, chi phí, chất lượng và nhân sự \\
Chủ biên chuyên môn & PGS. TS. Nguyễn Minh Hà & Giảng viên cao cấp ĐH Sư phạm & Chủ trì thiết kế khung cấu trúc, điều phối bốn nhóm tác giả, hiệu đính nội dung \\
Chủ tịch Hội đồng thẩm định & Tiến sĩ Vũ Thị Mai & Chuyên gia Viện Khoa học Giáo dục & Chủ trì thẩm định độc lập chất lượng khoa học và tính sư phạm của bản thảo \\
Phụ trách Mua sắm và Tài chính & Thạc sĩ Phạm Quang Thắng & Chuyên viên Kế hoạch Tài chính & Tổ chức chào hàng cạnh tranh in ấn, điều phối chế bản, quản lý giải ngân hợp đồng \\
Đại diện đơn vị thụ hưởng & Thạc sĩ Đỗ Thị Lan & Hiệu trưởng trường THCS cụm & Điều phối khảo sát thực địa, chủ trì kiểm định mẫu thống kê và tiếp nhận tài liệu \\
\bottomrule
\end{tabularx}
\end{table}

\subsection*{2.2. Cơ cấu và vai trò của các bên liên quan}
\addcontentsline{toc}{subsection}{2.2. Cơ cấu và vai trò của các bên liên quan}

Nhóm lập kế hoạch tiến hành nhận diện \textit{Stakeholders} (Các bên liên quan) tham gia hoặc chịu tác động của dự án, đánh giá mức độ ảnh hưởng và mức độ quan tâm của từng nhóm chủ thể để xác lập chiến lược phối hợp phù hợp tại Bảng \ref{tab:cac_ben_lien_quan}.

\begin{table}[htbp]
\centering
\caption{Phân tích cơ cấu và chiến lược quản lý các bên liên quan}
\label{tab:cac_ben_lien_quan}
\begin{tabularx}{\textwidth}{L{2.8cm}L{3.2cm}YCC}
\toprule
\textbf{Bên tham gia} & \textbf{Vai trò} & \textbf{Mối quan tâm cốt lõi} & \textbf{Ảnh hưởng} & \textbf{Quan tâm} \\
\midrule
Lãnh đạo Phòng GD\&ĐT & Cấp phê duyệt chủ trương & Tính tuân thủ pháp luật, tiến độ giải ngân và chất lượng xuất bản phẩm & Cao & Cao \\
Ban Quản lý dự án & Cơ quan điều hành thực hiện & Kiểm soát đường găng tiến độ, cân đối ngân sách và chất lượng từng khâu & Cao & Cao \\
Bốn nhóm tác giả & Sáng tạo nội dung bản thảo & Độ tin cậy bài trắc nghiệm Holland, tính sư phạm và tiến độ nộp bài & Cao & Cao \\
Hội đồng thẩm định & Đánh giá, phản biện độc lập & Chuẩn mực học thuật, độ chính xác pháp lý và khả năng tiếp nhận của học sinh & Cao & TB \\
Nhà thầu in ấn & Cung ứng kỹ thuật xuất bản & Quy cách kỹ thuật chế bản, độ chuẩn màu sắc và tiến độ thanh toán & TB & Cao \\
Hiệu trưởng 5 trường & Đơn vị phối hợp, tiếp nhận & Tiếp nhận đủ 2.000 cuốn tài liệu đạt chuẩn trước kỳ đăng ký dự thi & TB & Cao \\
Giáo viên tư vấn & Sử dụng tài liệu tư vấn & Cấu trúc tài liệu rõ ràng, có hướng dẫn chấm điểm trắc nghiệm dễ dùng & Thấp & Cao \\
Học sinh và phụ huynh & Đối tượng thụ hưởng trực tiếp & Ấn phẩm trình bày sinh động, ngôn ngữ dễ hiểu, thông tin trường nghề rõ ràng & Thấp & Cao \\
\bottomrule
\end{tabularx}
\end{table}

\begin{figure}[htbp]
\centering
\caption{Ma trận Ảnh hưởng và Mức quan tâm của các bên liên quan}
\label{fig:ma_tran_stakeholders}
\resizebox{0.75\textwidth}{!}{%
\begin{tikzpicture}[scale=7]
    % Khung trục tọa độ
    \draw[->, line width=1.0pt] (0,0) -- (1.05,0) node[right, font=\footnotesize] {Mức độ quan tâm};
    \draw[->, line width=1.0pt] (0,0) -- (0,1.05) node[above, font=\footnotesize] {Mức độ ảnh hưởng};
    
    % Đường phân chia 4 góc
    \draw[dashed, gray] (0.5,0) -- (0.5,1.0);
    \draw[dashed, gray] (0,0.5) -- (1.0,0.5);
    
    % Nhãn 4 góc phần tư
    \node[font=\scriptsize\bfseries, color=blue!70!black] at (0.25, 0.95) {GIỮ HÀI LÒNG};
    \node[font=\scriptsize\bfseries, color=red!70!black] at (0.75, 0.95) {QUẢN LÝ CHẶT CHẼ};
    \node[font=\scriptsize\bfseries, color=gray!80!black] at (0.25, 0.05) {THEO DÕI TỐI THIỂU};
    \node[font=\scriptsize\bfseries, color=green!60!black] at (0.75, 0.05) {CUNG CẤP THÔNG TIN};
    
    % Các điểm tọa độ bên liên quan
    \filldraw[red!80!black] (0.85, 0.90) circle (0.015) node[above left, font=\scriptsize] {Lãnh đạo Phòng GD};
    \filldraw[red!80!black] (0.95, 0.95) circle (0.015) node[below left, font=\scriptsize] {Ban QLDA};
    \filldraw[red!80!black] (0.88, 0.85) circle (0.015) node[below right, font=\scriptsize] {Nhóm tác giả};
    \filldraw[blue!80!black] (0.55, 0.80) circle (0.015) node[right, font=\scriptsize] {Hội đồng thẩm định};
    \filldraw[green!60!black] (0.85, 0.50) circle (0.015) node[right, font=\scriptsize] {Nhà thầu in ấn};
    \filldraw[green!60!black] (0.75, 0.45) circle (0.015) node[left, font=\scriptsize] {Hiệu trưởng 5 trường};
    \filldraw[green!60!black] (0.80, 0.25) circle (0.015) node[right, font=\scriptsize] {Giáo viên tư vấn};
    \filldraw[green!60!black] (0.90, 0.15) circle (0.015) node[left, font=\scriptsize] {Học sinh \& Phụ huynh};
\end{tikzpicture}%
}
\end{figure}

\subsection*{2.3. Kế hoạch quản lý tổng thể và quy chế điều hành dự án}
\addcontentsline{toc}{subsection}{2.3. Kế hoạch quản lý tổng thể và quy chế điều hành dự án}

Bản kế hoạch quản lý tổng thể đóng vai trò điều phối tám kế hoạch thành phần (phạm vi, tiến độ, chi phí, chất lượng, nhân lực, rủi ro, truyền thông và mua sắm) thành một thể thống nhất. Để vận hành dự án thông suốt, Ban Quản lý dự án áp dụng cơ chế \textit{Work Authorization System} (Hệ thống ủy quyền giao việc) bằng văn bản: mỗi gói công việc trên sơ đồ mạng chỉ được phát lệnh triển khai khi gói công việc tiền nhiệm liền trước đã có biên bản xác nhận hoàn thành đạt yêu cầu chất lượng.

Chế độ hội họp điều hành được duy trì theo bốn hình thức: họp giao ban tiến độ vào sáng thứ Hai hàng tuần để rà soát các công việc trên đường găng; họp đánh giá sau các mốc kiểm soát trọng tâm (ngày thứ 14 sau khi duyệt đề cương, ngày thứ 40 sau khi thu bản thảo lần một, ngày thứ 52 sau khi họp Hội đồng thẩm định và ngày thứ 76 sau khi duyệt bản bông); cùng các cuộc họp đột xuất và họp ứng biến được triệu tập trong vòng 24 giờ khi phát sinh vướng mắc kỹ thuật hoặc cần kích hoạt phương án dự phòng.

\subsection*{2.4. Cơ chế kiểm soát thay đổi tổng thể}
\addcontentsline{toc}{subsection}{2.4. Cơ chế kiểm soát thay đổi tổng thể}

Quy trình \textit{Integrated Change Control} (Kiểm soát thay đổi tổng thể) được thiết lập nhằm thẩm định mọi đề xuất điều chỉnh phát sinh trong quá trình thực hiện, giữ vững tính ổn định của ba đường cơ sở về phạm vi, tiến độ và ngân sách.

\begin{figure}[htbp]
\centering
\caption{Quy trình bốn bước kiểm soát thay đổi tổng thể}
\label{fig:kiem_soat_thay_doi}
\resizebox{0.95\textwidth}{!}{%
\begin{tikzpicture}[
    node distance = 0.6cm,
    step_box/.style = {rectangle, rounded corners=3pt, draw=blue!80!black, fill=blue!5, line width=0.8pt, text width=12cm, align=left, inner sep=6pt, font=\footnotesize},
    arr/.style = {-{Stealth[scale=1.0]}, line width=1.0pt, draw=blue!70!black}
]
\node (b1) [step_box] {\textbf{Bước 1: Tiếp nhận yêu cầu thay đổi} -- Lập Phiếu yêu cầu thay đổi bằng văn bản, nêu rõ lý do phát sinh và nội dung đề xuất.};
\node (b2) [step_box, below=of b1] {\textbf{Bước 2: Phân tích tác động tổng thể} -- Đánh giá ảnh hưởng tới dung lượng 80 trang ruột, tiến độ đường găng 90 ngày và ngân sách 185.000.000 đồng.};
\node (b3) [step_box, below=of b2] {\textbf{Bước 3: Thẩm định và phê duyệt} -- Giám đốc dự án duyệt thay đổi kỹ thuật nội bộ; Trưởng phòng GD\&ĐT duyệt thay đổi ngân sách hoặc mốc bàn giao.};
\node (b4) [step_box, below=of b3] {\textbf{Bước 4: Thực thi và cập nhật kế hoạch} -- Điều chỉnh đường cơ sở tiến độ, chi phí; lưu trữ hồ sơ và thông báo đồng bộ các bên liên quan.};

\draw [arr] (b1) -- (b2);
\draw [arr] (b2) -- (b3);
\draw [arr] (b3) -- (b4);
\end{tikzpicture}%
}
\end{figure}

Mọi đề xuất điều chỉnh trong quá trình biên soạn và chế bản phải được lập thành Phiếu yêu cầu thay đổi bằng văn bản (Bước 1). Ngay khi tiếp nhận, Giám đốc dự án cùng bộ phận chuyên môn phân tích tác động của đề xuất đối với dung lượng 80 trang ruột, kiểm tra xem sự thay đổi có nằm trên đường găng làm ảnh hưởng đến mốc bàn giao 90 ngày hay làm vượt dự toán 185.000.000 đồng hay không (Bước 2).

Về thẩm quyền phê duyệt (Bước 3), Giám đốc dự án quyết định các điều chỉnh kỹ thuật nội bộ không làm thay đổi tiến độ đường găng và không phát sinh chi phí; trong khi Lãnh đạo Phòng Giáo dục và Đào tạo phê duyệt các thay đổi liên quan đến việc sử dụng quỹ dự phòng hoặc điều chỉnh mốc bàn giao. Khi yêu cầu thay đổi được chấp thuận, Ban Quản lý dự án cập nhật lại tài liệu kế hoạch cơ sở và thông báo bằng văn bản tới các bên liên quan để thống nhất thực hiện (Bước 4).


\newpage

% 6. Chương 3: Kế hoạch quản lý phạm vi dự án
\section*{CHƯƠNG 3. KẾ HOẠCH QUẢN LÝ PHẠM VI DỰ ÁN}
\addcontentsline{toc}{section}{CHƯƠNG 3. KẾ HOẠCH QUẢN LÝ PHẠM VI DỰ ÁN}

Chương này xác lập văn bản \textit{Scope Statement} (Tuyên bố phạm vi dự án) nhằm định rõ quy cách sản phẩm chuyển giao và ranh giới công việc, đồng thời thiết lập \textit{Work Breakdown Structure -- WBS} (Cấu trúc phân rã công việc) kèm từ điển mô tả chi tiết và quy trình kiểm soát thay đổi phạm vi.

\subsection*{3.1. Tuyên bố phạm vi dự án và đặc tả sản phẩm}
\addcontentsline{toc}{subsection}{3.1. Tuyên bố phạm vi dự án và đặc tả sản phẩm}

Văn bản Tuyên bố phạm vi dự án là cam kết chính thức giữa Ban Quản lý dự án, Phòng Giáo dục và Đào tạo cùng năm trường trung học cơ sở thụ hưởng về các sản phẩm bàn giao.

Sản phẩm chuyển giao cuối cùng là 2.000 cuốn ấn phẩm ``Bộ tài liệu tư vấn phân luồng và định hướng nghề nghiệp cho học sinh cuối cấp trung học cơ sở'' kèm theo 01 bộ tệp tin chế bản điện tử gốc. Ấn phẩm được thiết kế theo khổ tiêu chuẩn 17 $\times$ 24 cm, gồm \textbf{80 trang ruột} (chia thành 5 tay sách 16 trang, in trên giấy Fort trắng định lượng 100 gsm chống lóa) và \textbf{4 trang bìa ngoài} (Bìa 1, 2, 3, 4 in trên giấy Couche định lượng 250 gsm cán màng mờ mặt ngoài), in offset bốn màu toàn bộ và gia công đóng gáy keo nhiệt.

Nội dung 80 trang ruột được cấu trúc thành phần mở đầu và bốn phần chuyên môn (tổng cộng 20 bài chuẩn), kế thừa quy trình định hướng nghề nghiệp bốn bước của Lê Thị Quỳnh Nga, Nguyễn Thị Hương và Chu Cẩm Thơ (2026) cùng nghiên cứu chuẩn hóa công cụ mật mã Holland của Phạm Mạnh Hà (2020):

Trang đầu sách (từ trang 1 đến trang 4, gồm 4 trang): Trang phụ bìa, lời giới thiệu, hướng dẫn cách khai thác tài liệu và mục lục tổng thể.

Phần 1 -- Khám phá tiềm năng và nhận diện bản thân (từ trang 5 đến trang 22, gồm 18 trang, 5 bài chuẩn): Giới thiệu bài trắc nghiệm sở thích nghề nghiệp theo mật mã Holland rút gọn gồm 36 câu hỏi dành cho lứa tuổi 14 đến 15, hướng dẫn tự chấm điểm, vẽ biểu đồ nhận diện nhóm sở thích nổi trội và đối chiếu với năng lực học tập các môn văn hóa.

Phần 2 -- Bản đồ các con đường học tập sau tốt nghiệp trung học cơ sở (từ trang 23 đến trang 44, gồm 22 trang, 6 bài chuẩn): Hệ thống hóa bằng sơ đồ đồ họa thông tin và bảng so sánh bốn hướng đi sau lớp 9 (trường trung học phổ thông công lập, trường trung học phổ thông ngoài công lập, trung tâm Giáo dục nghề nghiệp -- Giáo dục thường xuyên và các trường trung cấp, cao đẳng đào tạo hệ chín cộng), làm rõ điều kiện xét tuyển, thời gian đào tạo và lộ trình học liên thông.

Phần 3 -- Thế giới nghề nghiệp và chính sách hỗ trợ học nghề (từ trang 45 đến trang 62, gồm 18 trang, 5 bài chuẩn): Cung cấp thông tin dự báo nhu cầu nhân lực tại địa phương và giải thích cụ thể quy định miễn 100\% học phí học nghề đối với học sinh tốt nghiệp trung học cơ sở theo Nghị định số 81/2021/NĐ-CP và Nghị định số 97/2023/NĐ-CP.

Phần 4 -- Cẩm nang ra quyết định và đồng hành cùng gia đình (từ trang 63 đến trang 80, gồm 18 trang, 4 bài chuẩn): Hướng dẫn quy trình năm bước lựa chọn nguyện vọng, cung cấp bộ phiếu bài tập đối thoại giữa cha mẹ và học sinh, kèm danh bạ thông tin tuyển sinh của 25 cơ sở giáo dục nghề nghiệp uy tín trên địa bàn.

\subsection*{3.2. Cấu trúc phân rã công việc}
\addcontentsline{toc}{subsection}{3.2. Cấu trúc phân rã công việc}

\subsubsection*{3.2.1. Phương pháp tiếp cận và nguyên lý xây dựng}

Dự án áp dụng phương pháp \textit{Top-down} (Tiếp cận từ trên xuống) kết hợp kỹ thuật \textit{Mind-mapping} (Sơ đồ tư duy) để phân rã tổng thể dự án thành 5 giai đoạn ở cấp độ 1 và 20 gói công việc cụ thể ở cấp độ 2.

Cấu trúc phân rã công việc tuân thủ bảy nguyên lý chuẩn mực: bao quát đúng 100\% phạm vi công việc đã cam kết; mỗi gói công việc chỉ xuất hiện tại một vị trí duy nhất; mỗi gói công việc ở cấp thấp nhất chỉ gắn với một sản phẩm đầu ra và do duy nhất một cá nhân chịu trách nhiệm chính; phù hợp với trình tự tổ chức thực hiện thực tế; có sự tham gia xây dựng của các tiểu ban chuyên môn và được định nghĩa rõ trong bảng từ điển công việc.

\subsubsection*{3.2.2. Sơ đồ cây phân rã công việc}

\begin{figure}[htbp]
\centering
\caption{Sơ đồ cây phân rã công việc (WBS) của dự án}
\label{fig:so_do_wbs}
\resizebox{0.95\textwidth}{!}{%
\begin{tikzpicture}[
    node distance = 0.4cm and 0.3cm,
    root_box/.style = {rectangle, rounded corners=3pt, draw=red!80!black, fill=red!10, line width=1.1pt, text width=13cm, align=center, inner sep=5pt, font=\footnotesize\bfseries},
    level1_box/.style = {rectangle, rounded corners=3pt, draw=blue!80!black, fill=blue!10, line width=0.9pt, text width=2.4cm, align=center, inner sep=4pt, font=\scriptsize\bfseries},
    level2_box/.style = {rectangle, rounded corners=2pt, draw=gray!70, fill=gray!6, line width=0.6pt, text width=2.4cm, align=left, inner sep=3pt, font=\tiny},
    arr/.style = {-{Stealth[scale=0.8]}, line width=0.7pt, draw=gray!70!black}
]

% Root
\node (root) [root_box] {DỰ ÁN XÂY DỰNG BỘ TÀI LIỆU TƯ VẤN PHÂN LUỒNG TRUNG HỌC CƠ SỞ};

% Level 1
\node (g1) [level1_box, below left=0.8cm and 3.5cm of root] {1.0 KHẢO SÁT \&\\ĐỀ CƯƠNG};
\node (g2) [level1_box, right=0.35cm of g1] {2.0 BIÊN SOẠN\\BẢN THẢO};
\node (g3) [level1_box, right=0.35cm of g2] {3.0 THẨM ĐỊNH\\\& HIỆU ĐÍNH};
\node (g4) [level1_box, right=0.35cm of g3] {4.0 CHẾ BẢN\\\& IN ẤN};
\node (g5) [level1_box, right=0.35cm of g4] {5.0 KIỂM ĐỊNH\\\& BÀN GIAO};

\draw [arr] (root.south) -- ++(0,-0.3) -| (g1.north);
\draw [arr] (root.south) -- ++(0,-0.3) -| (g2.north);
\draw [arr] (root.south) -- ++(0,-0.3) -| (g3.north);
\draw [arr] (root.south) -- ++(0,-0.3) -| (g4.north);
\draw [arr] (root.south) -- ++(0,-0.3) -| (g5.north);

% Level 2 - Branch 1
\node (g1_1) [level2_box, below=0.3cm of g1] {1.1 Khảo sát nhu cầu (350 phiếu)};
\node (g1_2) [level2_box, below=0.15cm of g1_1] {1.2 Dữ liệu và phân tích số liệu};
\node (g1_3) [level2_box, below=0.15cm of g1_2] {1.3 Xây dựng, duyệt đề cương};
\draw [arr] (g1.south) -- (g1_1.north);

% Level 2 - Branch 2
\node (g2_1) [level2_box, below=0.3cm of g2] {2.1 Viết Phần 1 (18 trang)};
\node (g2_2) [level2_box, below=0.15cm of g2_1] {2.2 Viết Phần 2 (22 trang)};
\node (g2_3) [level2_box, below=0.15cm of g2_2] {2.3 Viết Phần 3 (18 trang)};
\node (g2_4) [level2_box, below=0.15cm of g2_3] {2.4 Viết Phần 4 (18 trang)};
\node (g2_5) [level2_box, below=0.15cm of g2_4] {2.5 Biên tập bản thảo lần 1};
\draw [arr] (g2.south) -- (g2_1.north);

% Level 2 - Branch 3
\node (g3_1) [level2_box, below=0.3cm of g3] {3.1 Phản biện 3 chuyên gia};
\node (g3_2) [level2_box, below=0.15cm of g3_1] {3.2 Họp Hội đồng 2 phiên};
\node (g3_3) [level2_box, below=0.15cm of g3_2] {3.3 Lấy ý kiến đọc thử};
\node (g3_4) [level2_box, below=0.15cm of g3_3] {3.4 Hiệu đính bản sạch};
\draw [arr] (g3.south) -- (g3_1.north);

% Level 2 - Branch 4
\node (g4_1) [level2_box, below=0.3cm of g4] {4.1 Lựa chọn nhà thầu in};
\node (g4_2) [level2_box, below=0.15cm of g4_1] {4.2 Vẽ hình, dàn trang};
\node (g4_3) [level2_box, below=0.15cm of g4_2] {4.3 In thử duyệt bản bông};
\node (g4_4) [level2_box, below=0.15cm of g4_3] {4.4 In offset 2.000 cuốn};
\draw [arr] (g4.south) -- (g4_1.north);

% Level 2 - Branch 5
\node (g5_1) [level2_box, below=0.3cm of g5] {5.1 Kiểm định mẫu (384 c)};
\node (g5_2) [level2_box, below=0.15cm of g5_1] {5.2 Họp Hội đồng nghiệm thu};
\node (g5_3) [level2_box, below=0.15cm of g5_2] {5.3 Đóng thùng, bàn giao};
\node (g5_4) [level2_box, below=0.15cm of g5_3] {5.4 Thanh lý \& lưu hồ sơ};
\draw [arr] (g5.south) -- (g5_1.north);

\end{tikzpicture}%
}
\end{figure}

\subsubsection*{3.2.3. Bảng từ điển phân rã công việc chi tiết}

Tài liệu \textit{WBS Dictionary} (Từ điển phân rã công việc) tại Bảng \ref{tab:tu_dien_wbs} xác định rõ phạm vi tác nghiệp, kết quả bàn giao và cá nhân chịu trách nhiệm chính cho toàn bộ 20 gói công việc, làm cơ sở trực tiếp để lập dự toán chi phí tại Chương 5.

\begin{table}[htbp]
\centering
\caption{Từ điển phân rã công việc chi tiết của dự án}
\label{tab:tu_dien_wbs}
\begin{tabularx}{\textwidth}{cCYYL{3.4cm}}
\toprule
\textbf{Mã} & \textbf{Tên gói công việc} & \textbf{Nội dung triển khai} & \textbf{Sản phẩm đầu ra} & \textbf{Chịu trách nhiệm} \\
\midrule
\textbf{1.0} & \multicolumn{3}{l}{\textbf{Khảo sát nhu cầu và lập khung cấu trúc}} & \textbf{PGS. TS. Nguyễn Minh Hà} \\
1.1 & Khảo sát thực tế & Thu thập 350 bộ phiếu khảo sát (300 HS, 50 GV/PH) tại 5 trường & 350 bộ phiếu hợp lệ & Chuyên viên điều tra \\
1.2 & Dữ liệu tuyển sinh & Xử lý số liệu, tổng hợp chỉ tiêu 25 trường nghề và chính sách học phí & Báo cáo kết quả dữ liệu & Thư ký dự án \\
1.3 & Duyệt đề cương & Thiết kế ma trận 4 phần và hội thảo chuyên môn (12 đại biểu) & Bản đề cương được duyệt & Chủ biên chuyên môn \\
\midrule
\textbf{2.0} & \multicolumn{3}{l}{\textbf{Biên soạn bản thảo bộ tài liệu (80 trang ruột)}} & \textbf{PGS. TS. Nguyễn Minh Hà} \\
2.1 & Viết Phần 1 & Viết 5 bài và chuẩn hóa bảng hỏi Holland 36 câu (trang 5--22) & Bản thảo P1 (18 trang) & Trưởng nhóm TG1 \\
2.2 & Viết Phần 2 & Viết 6 bài và lập ma trận so sánh 4 hướng đi sau THCS (trang 23--44) & Bản thảo P2 (22 trang) & Trưởng nhóm TG2 \\
2.3 & Viết Phần 3 & Viết 5 bài xu hướng nghề và chính sách miễn 100\% học phí (trang 45--62) & Bản thảo P3 (18 trang) & Trưởng nhóm TG3 \\
2.4 & Viết Phần 4 & Viết 4 bài, phiếu đối thoại gia đình và danh bạ 25 trường (trang 63--80) & Bản thảo P4 (18 trang) & Trưởng nhóm TG4 \\
2.5 & Biên tập bản thảo 1 & Soạn 4 trang đầu sách, ráp nối 20 bài và biên tập 80 trang ruột & Bản thảo lần một (80 tr) & Chủ biên chuyên môn \\
\midrule
\textbf{3.0} & \multicolumn{3}{l}{\textbf{Thẩm định và hiệu đính bản thảo}} & \textbf{Tiến sĩ Vũ Thị Mai} \\
3.1 & Phản biện độc lập & 3 chuyên gia độc lập đọc phản biện toàn văn 80 trang bản thảo & 3 bản nhận xét phản biện & Chủ tịch HĐ thẩm định \\
3.2 & Họp Hội đồng & Tổ chức 2 phiên họp Hội đồng thẩm định (5 thành viên) & Biên bản họp Hội đồng & Chủ tịch HĐ thẩm định \\
3.3 & Lấy ý kiến đọc thử & Thảo luận nhóm tập trung đọc thử với 20 học sinh và 5 giáo viên & Báo cáo ý kiến đọc thử & Thư ký HĐ thẩm định \\
3.4 & Hiệu đính bản sạch & Bốn nhóm tác giả chỉnh sửa theo kết luận, rà soát chính tả & Bản thảo sạch hoàn chỉnh & Chủ biên chuyên môn \\
\midrule
\textbf{4.0} & \multicolumn{3}{l}{\textbf{Thiết kế chế bản và in ấn tài liệu}} & \textbf{ThS. Phạm Quang Thắng} \\
4.1 & Lựa chọn nhà thầu & Chào hàng cạnh tranh rút gọn và ký hợp đồng trọn gói & Báo cáo thầu và Hợp đồng & Phụ trách Mua sắm \\
4.2 & Thiết kế dàn trang & Thiết kế 4 trang bìa ngoài, vẽ 25 hình minh họa và dàn trang điện tử & Tệp ma-két hoàn chỉnh & Nhóm thiết kế chế bản \\
4.3 & In thử duyệt mẫu & In thử 2 bộ bản bông màu, rà soát lỗi và ký duyệt niêm phong & Bản bông có chữ ký duyệt & Chủ biên chuyên môn \\
4.4 & In offset hàng loạt & Bình bản xuất kẽm, in 2.000 cuốn, cán màng mờ và đóng keo nhiệt & 2.000 cuốn tài liệu & Giám đốc nhà thầu in \\
\midrule
\textbf{5.0} & \multicolumn{3}{l}{\textbf{Kiểm định, nghiệm thu và bàn giao}} & \textbf{Thạc sĩ Lê Hoàng Nam} \\
5.1 & Kiểm định mẫu & Rút mẫu phân tầng 384 cuốn kiểm tra quy cách kỹ thuật tại kho & Phiếu kiểm định 384 cuốn & Đại diện đơn vị thụ hưởng \\
5.2 & Họp nghiệm thu & Hội đồng nghiệm thu cấp cơ sở đánh giá sản phẩm và ký biên bản & Biên bản nghiệm thu & Giám đốc dự án \\
5.3 & Đóng thùng bàn giao & Đóng 40 thùng giấy (50 c/thùng), vận chuyển bàn giao cho 5 trường & Biên bản bàn giao 5 trường & Đại diện đơn vị thụ hưởng \\
5.4 & Thanh lý \& lưu trữ & Quyết toán hợp đồng in ấn, đóng gói 03 bộ hồ sơ lưu trữ gốc & Biên bản thanh lý \& hồ sơ & Phụ trách Mua sắm \\
\bottomrule
\end{tabularx}
\end{table}

\subsection*{3.3. Quy trình xác nhận và kiểm soát sự thay đổi phạm vi}
\addcontentsline{toc}{subsection}{3.3. Quy trình xác nhận và kiểm soát sự thay đổi phạm vi}

Xác nhận phạm vi được thực hiện tại điểm kết thúc của từng giai đoạn thông qua việc đối chiếu sản phẩm trung gian thực tế với tiêu chí đầu ra trong Bảng \ref{tab:tu_dien_wbs}. Nhằm ngăn ngừa hiện tượng \textit{Scope Creep} (Trượt phạm vi dự án), Ban Quản lý dự án áp dụng hai chốt kiểm soát định lượng: khống chế dung lượng phần ruột ở mức cố định 80 trang (5 tay sách) để không làm biến động đơn giá chế bản và giấy in; đồng thời từ chối mọi đề xuất bổ sung các hoạt động ngoài phạm vi xuất bản học liệu đã được phê duyệt tại Điều lệ dự án.


\newpage

% 7. Chương 4: Kế hoạch quản lý tiến độ dự án
\section*{CHƯƠNG 4. KẾ HOẠCH QUẢN LÝ TIẾN ĐỘ DỰ ÁN}
\addcontentsline{toc}{section}{CHƯƠNG 4. KẾ HOẠCH QUẢN LÝ TIẾN ĐỘ DỰ ÁN}

Chương này xác định danh mục mười công việc trên sơ đồ mạng từ cấu trúc phân rã công việc, ước lượng thời gian thực hiện, mô hình hóa mạng công việc theo phương pháp mũi tên và phương pháp nút, tính toán sáu bước thông số mạng theo phương pháp đường găng, lập biểu đồ tiến độ ngang Gantt và đề xuất phương án điều hòa nguồn lực trong khung thời gian 90 ngày lịch (từ ngày 02/03/2026 đến ngày 30/05/2026).

\subsection*{4.1. Danh mục công việc và ước lượng thời gian thực hiện}
\addcontentsline{toc}{subsection}{4.1. Danh mục công việc và ước lượng thời gian thực hiện}

Từ 20 gói công việc tại Chương 3, tiến trình thực thi dự án được tích hợp thành mười công việc mạng ký hiệu từ A đến J, nối giữa mười sự kiện đánh số từ 1 đến 10. Thời gian thực hiện từng công việc ($t_{ij}$) được ước lượng theo phương pháp tất định dựa trên khối lượng trang viết, quy trình thẩm định xuất bản và định mức thời gian chào hàng cạnh tranh.

\begin{table}[htbp]
\centering
\caption{Danh mục công việc, mã phân rã tương ứng và thời gian thực hiện}
\label{tab:danh_muc_cong_viec}
\begin{tabularx}{\textwidth}{cCCCYcL}
\toprule
\textbf{TT} & \textbf{Ký hiệu} & \textbf{Sự kiện} & \textbf{Mã WBS} & \textbf{Nội dung công việc chi tiết} & \textbf{$t_{ij}$ (ngày)} & \textbf{Điều kiện bắt đầu} \\
\midrule
\textbf{1} & \textbf{A} & $(1 - 2)$ & 1.1, 1.2 & Khảo sát nhu cầu (350 phiếu) \& dữ liệu & \textbf{8} & Khởi công ngày 1 (02/03/2026) \\
\textbf{2} & \textbf{B} & $(2 - 3)$ & 1.3 & Xây dựng cấu trúc \& phê duyệt đề cương & \textbf{6} & Ngay sau khi xong A \\
\textbf{3} & \textbf{C} & $(3 - 4)$ & 2.1--2.5 & Biên soạn bản thảo lần một (80 trang) & \textbf{26} & Ngay sau khi xong B \\
\textbf{4} & \textbf{D} & $(3 - 5)$ & 4.1 & Lựa chọn nhà in và ký hợp đồng & \textbf{18} & Sau B (song song với C, E) \\
\textbf{5} & \textbf{E} & $(4 - 5)$ & 3.1--3.3 & Phản biện độc lập \& họp thẩm định & \textbf{12} & Ngay sau khi xong C \\
\textbf{6} & \textbf{F} & $(5 - 6)$ & 3.4 & Tiếp thu góp ý, hiệu đính bản sạch & \textbf{8} & Sau khi xong cả D và E \\
\textbf{7} & \textbf{G} & $(6 - 7)$ & 4.2 & Thiết kế bìa, vẽ 25 hình, dàn trang & \textbf{12} & Ngay sau khi xong F \\
\textbf{8} & \textbf{H} & $(7 - 8)$ & 4.3 & In thử bản bông màu, ký duyệt mẫu & \textbf{4} & Ngay sau khi xong G \\
\textbf{9} & \textbf{I} & $(8 - 9)$ & 4.4 & In offset 2.000 cuốn, cán màng, keo gáy & \textbf{8} & Ngay sau khi xong H \\
\textbf{10} & \textbf{J} & $(9 - 10)$ & 5.1--5.4 & Kiểm định 384 cuốn, nghiệm thu bàn giao & \textbf{6} & Ngay sau khi xong I \\
\bottomrule
\end{tabularx}
\end{table}

\subsection*{4.2. Sơ đồ mạng công việc theo phương pháp mũi tên và phương pháp nút}
\addcontentsline{toc}{subsection}{4.2. Sơ đồ mạng công việc theo phương pháp mũi tên và phương pháp nút}

\subsubsection*{4.2.1. Sơ đồ mạng công việc theo phương pháp mũi tên (AOA)}

Phương pháp \textit{Activity on Arrow -- AOA} (Phương pháp biểu diễn công việc trên mũi tên) biểu thị mỗi công việc bằng một mũi tên nối giữa hai sự kiện tròn. Tại Sự kiện 3 (sau khi duyệt đề cương), mạng công việc tách thành hai nhánh chạy song song gồm nhánh biên soạn, thẩm định học thuật ($C \rightarrow E$) và nhánh lựa chọn nhà thầu in ấn ($D$). Hai nhánh này hội tụ tại Sự kiện 5, xác lập điều kiện rằng khâu hoàn thiện bản thảo sạch để chuyển giao chế bản (công việc F) chỉ bắt đầu khi cả bản thảo đã được Hội đồng thẩm định thông qua và hợp đồng in ấn đã được ký kết.

\begin{figure}[htbp]
\centering
\caption{Sơ đồ mạng công việc theo phương pháp mũi tên (AOA)}
\label{fig:so_do_aoa}
\resizebox{0.95\textwidth}{!}{%
\begin{tikzpicture}[
    node distance = 1.3cm and 1.5cm,
    event/.style = {circle, draw=red!80!black, fill=red!6, line width=1.1pt, minimum size=0.9cm, inner sep=0pt, font=\small\bfseries},
    arr_crit/.style = {-{Stealth[scale=1.0]}, line width=1.3pt, draw=red!80!black},
    arr_noncrit/.style = {-{Stealth[scale=1.0]}, line width=1.0pt, draw=blue!80!black, dashed}
]

\node (s1) [event] {1};
\node (s2) [event, right=1.6cm of s1] {2};
\node (s3) [event, right=1.6cm of s2] {3};
\node (s4) [event, above right=0.9cm and 2.2cm of s3] {4};
\node (s5) [event, below right=0.9cm and 2.2cm of s4] {5};
\node (s6) [event, right=1.6cm of s5] {6};
\node (s7) [event, right=1.6cm of s6] {7};
\node (s8) [event, right=1.6cm of s7] {8};
\node (s9) [event, right=1.6cm of s8] {9};
\node (s10) [event, right=1.6cm of s9] {10};

\draw [arr_crit] (s1) -- node[above, font=\footnotesize] {A, 8} (s2);
\draw [arr_crit] (s2) -- node[above, font=\footnotesize] {B, 6} (s3);
\draw [arr_crit] (s3) -- node[above left, font=\footnotesize] {C, 26} (s4);
\draw [arr_noncrit] (s3) -- node[below, font=\footnotesize] {D, 18} (s5);
\draw [arr_crit] (s4) -- node[above right, font=\footnotesize] {E, 12} (s5);
\draw [arr_crit] (s5) -- node[above, font=\footnotesize] {F, 8} (s6);
\draw [arr_crit] (s6) -- node[above, font=\footnotesize] {G, 12} (s7);
\draw [arr_crit] (s7) -- node[above, font=\footnotesize] {H, 4} (s8);
\draw [arr_crit] (s8) -- node[above, font=\footnotesize] {I, 8} (s9);
\draw [arr_crit] (s9) -- node[above, font=\footnotesize] {J, 6} (s10);

\end{tikzpicture}%
}
\end{figure}

\subsubsection*{4.2.2. Sơ đồ mạng công việc theo phương pháp nút (AON)}

Phương pháp \textit{Activity on Node -- AON} (Phương pháp biểu diễn công việc trên nút) đặt các công việc vào trong các ô hình chữ nhật kèm theo thông số thời điểm bắt đầu sớm ($ES$), thời điểm kết thúc sớm ($EF$) và thời gian thực hiện ($t$), còn các mũi tên thể hiện quan hệ phụ thuộc trước sau:

\begin{figure}[htbp]
\centering
\caption{Sơ đồ mạng công việc theo phương pháp nút (AON)}
\label{fig:so_do_aon}
\resizebox{0.95\textwidth}{!}{%
\begin{tikzpicture}[
    node distance = 0.5cm and 0.6cm,
    crit_box/.style = {rectangle, rounded corners=3pt, draw=red!80!black, fill=red!6, line width=1.1pt, text width=2.6cm, align=center, inner sep=4pt, font=\scriptsize},
    noncrit_box/.style = {rectangle, rounded corners=3pt, draw=blue!80!black, fill=blue!6, line width=1.0pt, text width=2.6cm, align=center, inner sep=4pt, font=\scriptsize},
    arr_c/.style = {-{Stealth[scale=0.9]}, line width=1.1pt, draw=red!80!black},
    arr_nc/.style = {-{Stealth[scale=0.9]}, line width=0.9pt, draw=blue!80!black}
]

\node (A) [crit_box] {\textbf{[A] Khảo sát nhu cầu}\\ES = 0 $\vert$ EF = 8\\t = 8 ngày};
\node (B) [crit_box, right=of A] {\textbf{[B] Duyệt đề cương}\\ES = 8 $\vert$ EF = 14\\t = 6 ngày};
\node (C) [crit_box, above right=0.3cm and 0.6cm of B] {\textbf{[C] Biên soạn bản thảo}\\ES = 14 $\vert$ EF = 40\\t = 26 ngày};
\node (D) [noncrit_box, below right=0.3cm and 0.6cm of B] {\textbf{[D] Chọn nhà in ấn}\\ES = 14 $\vert$ EF = 32\\t = 18 ngày};
\node (E) [crit_box, right=of C] {\textbf{[E] Thẩm định bài}\\ES = 40 $\vert$ EF = 52\\t = 12 ngày};
\node (F) [crit_box, below right=0.3cm and 0.6cm of E] {\textbf{[F] Hiệu đính bản sạch}\\ES = 52 $\vert$ EF = 60\\t = 8 ngày};
\node (G) [crit_box, right=of F] {\textbf{[G] Thiết kế dàn trang}\\ES = 60 $\vert$ EF = 72\\t = 12 ngày};
\node (H) [crit_box, right=of G] {\textbf{[H] In thử duyệt mẫu}\\ES = 72 $\vert$ EF = 76\\t = 4 ngày};
\node (I) [crit_box, right=of H] {\textbf{[I] In offset hàng loạt}\\ES = 76 $\vert$ EF = 84\\t = 8 ngày};
\node (J) [crit_box, right=of I] {\textbf{[J] Nghiệm thu bàn giao}\\ES = 84 $\vert$ EF = 90\\t = 6 ngày};

\draw [arr_c] (A) -- (B);
\draw [arr_c] (B) |- (C);
\draw [arr_nc] (B) |- (D);
\draw [arr_c] (C) -- (E);
\draw [arr_c] (E) -| (F);
\draw [arr_nc] (D) -| (F);
\draw [arr_c] (F) -- (G);
\draw [arr_c] (G) -- (H);
\draw [arr_c] (H) -- (I);
\draw [arr_c] (I) -- (J);

\end{tikzpicture}%
}
\end{figure}

\subsection*{4.3. Các thông số sơ đồ mạng và xác định đường găng}
\addcontentsline{toc}{subsection}{4.3. Các thông số sơ đồ mạng và xác định đường găng}

Việc giải sơ đồ mạng được tiến hành theo quy trình sáu bước của phương pháp \textit{Critical Path Method -- CPM} (Phương pháp đường găng).

\subsubsection*{4.3.1. Thời điểm sớm của các sự kiện}

Thời điểm sớm của sự kiện $j$ ($E_j$) được tính theo chiều thuận từ Sự kiện 1 đến Sự kiện 10 theo công thức $E_1 = 0$ và $E_j = \max_{i} (E_i + t_{ij})$. Cụ thể, xuất phát từ $E_1 = 0$, ta xác định được $E_2 = 0 + 8 = 8$ ngày; $E_3 = 8 + 6 = 14$ ngày; và $E_4 = 14 + 26 = 40$ ngày. Tại Sự kiện 5, do có hai công việc E và D cùng đi vào, thời điểm sớm được lấy theo giá trị cực đại giữa hai nhánh:
\[
E_5 = \max\left[E_4 + t_{45};\ E_3 + t_{35}\right] = \max\left[40 + 12;\ 14 + 18\right] = \max[52;\ 32] = 52\text{ ngày}
\]
Tiếp tục tính thuận theo chuỗi công việc nối tiếp phía sau, ta thu được $E_6 = 52 + 8 = 60$ ngày; $E_7 = 60 + 12 = 72$ ngày; $E_8 = 72 + 4 = 76$ ngày; $E_9 = 76 + 8 = 84$ ngày; và tại sự kiện kết thúc dự án là $E_{10} = 84 + 6 = 90$ ngày.

\subsubsection*{4.3.2. Thời điểm muộn của các sự kiện}

Thời điểm muộn của sự kiện $i$ ($L_i$) được tính lùi từ Sự kiện 10 về Sự kiện 1 theo công thức $L_{10} = E_{10} = 90$ ngày và $L_i = \min_{j} (L_j - t_{ij})$. Lần lượt trừ lùi thời gian thực hiện của các công việc cuối mạng, ta có $L_9 = 90 - 6 = 84$ ngày; $L_8 = 84 - 8 = 76$ ngày; $L_7 = 76 - 4 = 72$ ngày; $L_6 = 72 - 12 = 60$ ngày; $L_5 = 60 - 8 = 52$ ngày; và $L_4 = 52 - 12 = 40$ ngày. Tại Sự kiện 3, do có hai công việc C và D cùng đi ra, thời điểm muộn được xác định bằng giá trị cực tiểu:
\[
L_3 = \min\left[L_4 - t_{34};\ L_5 - t_{35}\right] = \min\left[40 - 26;\ 52 - 18\right] = \min[14;\ 34] = 14\text{ ngày}
\]
Từ kết quả $L_3 = 14$ ngày, tính lùi tiếp về đầu mạng xác định được $L_2 = 14 - 6 = 8$ ngày và $L_1 = 8 - 8 = 0$ ngày.

\subsubsection*{4.3.3. Thời gian dự trữ sự kiện và xác định đường găng}

Thời gian dự trữ của sự kiện $j$ được tính bằng hiệu số $R_j = L_j - E_j$. Các sự kiện có $R_j = 0$ là các sự kiện găng.

\begin{table}[htbp]
\centering
\caption{Thời gian dự trữ của các sự kiện và xác định sự kiện găng}
\label{tab:du_tru_su_kien}
\begin{tabularx}{0.85\textwidth}{CCCCCY}
\toprule
\textbf{Sự kiện ($j$)} & \textbf{$E_j$ (ngày)} & \textbf{$L_j$ (ngày)} & \textbf{$R_j = L_j - E_j$} & \textbf{Tính chất sự kiện} \\
\midrule
\textbf{1} & 0 & 0 & 0 & \textbf{Sự kiện găng} \\
\textbf{2} & 8 & 8 & 0 & \textbf{Sự kiện găng} \\
\textbf{3} & 14 & 14 & 0 & \textbf{Sự kiện găng} \\
\textbf{4} & 40 & 40 & 0 & \textbf{Sự kiện găng} \\
\textbf{5} & 52 & 52 & 0 & \textbf{Sự kiện găng} \\
\textbf{6} & 60 & 60 & 0 & \textbf{Sự kiện găng} \\
\textbf{7} & 72 & 72 & 0 & \textbf{Sự kiện găng} \\
\textbf{8} & 76 & 76 & 0 & \textbf{Sự kiện găng} \\
\textbf{9} & 84 & 84 & 0 & \textbf{Sự kiện găng} \\
\textbf{10} & 90 & 90 & 0 & \textbf{Sự kiện găng} \\
\bottomrule
\end{tabularx}
\end{table}

Đường găng của dự án đi qua toàn bộ mười sự kiện găng theo tuyến chín công việc nối tiếp:
\[
A \rightarrow B \rightarrow C \rightarrow E \rightarrow F \rightarrow G \rightarrow H \rightarrow I \rightarrow J
\]
Độ dài đường găng ấn định tổng thời gian hoàn thành dự án:
\[
T_{CP} = 8 + 6 + 26 + 12 + 8 + 12 + 4 + 8 + 6 = 90\text{ ngày}
\]

\subsubsection*{4.3.4. Các thông số thời gian của công việc và bảng tổng hợp thông số mạng}

Đối với mỗi công việc $(i - j)$, bốn mốc thời điểm sớm/muộn được xác định gồm: bắt đầu sớm ($ES_{ij} = E_i$), kết thúc sớm ($EF_{ij} = E_i + t_{ij}$), kết thúc muộn ($LF_{ij} = L_j$) và bắt đầu muộn ($LS_{ij} = L_j - t_{ij}$). Bốn thông số thời gian dự trữ của công việc gồm: dự trữ toàn phần ($GR_{ij} = L_j - E_i - t_{ij}$), dự trữ bắt đầu ($SR_{ij} = L_j - L_i - t_{ij}$), dự trữ kết thúc ($FR_{ij} = E_j - E_i - t_{ij}$) và dự trữ độc lập ($IR_{ij} = E_j - L_i - t_{ij}$).

\begin{table}[htbp]
\centering
\caption{Bảng tổng hợp các thông số thời gian và dự trữ của công việc (CPM)}
\label{tab:cpm_tong_hop}
\resizebox{\textwidth}{!}{%
\begin{tabular}{cccccccccccc}
\toprule
\textbf{Công việc} & \textbf{Mã $(i - j)$} & \textbf{$t_{ij}$} & \textbf{$ES_{ij}$} & \textbf{$EF_{ij}$} & \textbf{$LF_{ij}$} & \textbf{$LS_{ij}$} & \textbf{$GR_{ij}$} & \textbf{$SR_{ij}$} & \textbf{$FR_{ij}$} & \textbf{$IR_{ij}$} & \textbf{Tính chất} \\
\midrule
\textbf{A} & $(1 - 2)$ & 8 & 0 & 8 & 8 & 0 & \textbf{0} & 0 & 0 & 0 & \textbf{Găng} \\
\textbf{B} & $(2 - 3)$ & 6 & 8 & 14 & 14 & 8 & \textbf{0} & 0 & 0 & 0 & \textbf{Găng} \\
\textbf{C} & $(3 - 4)$ & 26 & 14 & 40 & 40 & 14 & \textbf{0} & 0 & 0 & 0 & \textbf{Găng} \\
\textbf{D} & $(3 - 5)$ & 18 & 14 & 32 & 52 & 34 & \textbf{20} & 20 & 20 & 20 & \textbf{Không găng} \\
\textbf{E} & $(4 - 5)$ & 12 & 40 & 52 & 52 & 40 & \textbf{0} & 0 & 0 & 0 & \textbf{Găng} \\
\textbf{F} & $(5 - 6)$ & 8 & 52 & 60 & 60 & 52 & \textbf{0} & 0 & 0 & 0 & \textbf{Găng} \\
\textbf{G} & $(6 - 7)$ & 12 & 60 & 72 & 72 & 60 & \textbf{0} & 0 & 0 & 0 & \textbf{Găng} \\
\textbf{H} & $(7 - 8)$ & 4 & 72 & 76 & 76 & 72 & \textbf{0} & 0 & 0 & 0 & \textbf{Găng} \\
\textbf{I} & $(8 - 9)$ & 8 & 76 & 84 & 84 & 76 & \textbf{0} & 0 & 0 & 0 & \textbf{Găng} \\
\textbf{J} & $(9 - 10)$ & 6 & 84 & 90 & 90 & 84 & \textbf{0} & 0 & 0 & 0 & \textbf{Găng} \\
\bottomrule
\end{tabular}%
}
\end{table}

Bảng \ref{tab:cpm_tong_hop} cho thấy chín công việc trên đường găng có mọi thông số dự trữ bằng không. Riêng công việc D (lựa chọn nhà thầu in ấn) có thời gian dự trữ ở cả bốn chỉ số bằng 20 ngày ($GR_D = SR_D = FR_D = IR_D = 20$ ngày), cho phép bộ phận mua sắm hoàn thành ký hợp đồng sớm nhất vào ngày thứ 32 hoặc muộn nhất vào ngày thứ 52 mà không làm chậm tiến độ chung.

\subsection*{4.4. Phương pháp biểu đồ tiến độ ngang Gantt}
\addcontentsline{toc}{subsection}{4.4. Phương pháp biểu đồ tiến độ ngang Gantt}

Trục thời gian 90 ngày lịch được chia thành 13 tuần (mỗi tuần 7 ngày, riêng tuần 13 gồm 6 ngày cuối từ ngày 85 đến ngày 90). Bảng \ref{tab:bieu_do_gantt_tuan} mô tả tiến trình thực hiện của toàn bộ mười công việc qua các tuần:

\begin{table}[htbp]
\centering
\caption{Phân bổ tiến độ dự án theo tuần trên biểu đồ ngang Gantt}
\label{tab:bieu_do_gantt_tuan}
\resizebox{\textwidth}{!}{%
\begin{tabular}{llcccccccccccc}
\toprule
\textbf{Mã} & \textbf{Tên công việc} & \textbf{$t$} & \textbf{T1--2} & \textbf{T3--4} & \textbf{T5--6} & \textbf{T7--8} & \textbf{T9} & \textbf{T10--11} & \textbf{T12} & \textbf{T13} \\
& & \textbf{(ngày)} & \textbf{(N1--14)} & \textbf{(N15--28)} & \textbf{(N29--42)} & \textbf{(N43--56)} & \textbf{(N57--63)} & \textbf{(N64--77)} & \textbf{(N78--84)} & \textbf{(N85--90)} \\
\midrule
\textbf{A} & Khảo sát nhu cầu \& dữ liệu & 8 & $\blacksquare\blacksquare\blacksquare\blacksquare$ & & & & & & & \\
\textbf{B} & Duyệt khung đề cương chi tiết & 6 & \phantom{$\blacksquare\blacksquare$}$\blacksquare\blacksquare$ & & & & & & & \\
\textbf{C} & Biên soạn bản thảo lần một & 26 & & $\blacksquare\blacksquare\blacksquare\blacksquare$ & $\blacksquare\blacksquare\blacksquare$ & & & & & \\
\textbf{D} & Chọn nhà in ấn \& ký hợp đồng & 18 & & $\square\square\square\square$ & $\square$ \footnotesize{[+20d]} & \footnotesize{[+20d]} & & & & \\
\textbf{E} & Thẩm định chuyên môn bản thảo & 12 & & & \phantom{$\blacksquare\blacksquare$}$\blacksquare$ & $\blacksquare\blacksquare\blacksquare$ & & & & \\
\textbf{F} & Hiệu đính, hoàn thiện bản sạch & 8 & & & & \phantom{$\blacksquare\blacksquare$}$\blacksquare$ & $\blacksquare$ & & & \\
\textbf{G} & Thiết kế đồ họa và dàn trang & 12 & & & & & $\blacksquare$ & $\blacksquare\blacksquare\blacksquare$ & & \\
\textbf{H} & In bản bông thử và duyệt mẫu & 4 & & & & & & \phantom{$\blacksquare\blacksquare$}$\blacksquare$ & & \\
\textbf{I} & In offset hàng loạt 2.000 cuốn & 8 & & & & & & \phantom{$\blacksquare\blacksquare$}$\blacksquare$ & $\blacksquare\blacksquare$ & \\
\textbf{J} & Kiểm định mẫu và bàn giao & 6 & & & & & & & & $\blacksquare\blacksquare$ \\
\bottomrule
\end{tabular}%
}
\end{table}

\subsection*{4.5. Cơ chế kiểm soát tiến độ và điều hòa nguồn lực}
\addcontentsline{toc}{subsection}{4.5. Cơ chế kiểm soát tiến độ và điều hòa nguồn lực}

\subsubsection*{4.5.1. Kiểm soát các mốc sự kiện găng trung gian}

Ban Quản lý dự án giám sát chặt chẽ bảy mốc sự kiện găng gắn với sản phẩm nghiệm thu và lịch giải ngân: Mốc 1 tại ngày thứ 14 (15/03/2026) phê duyệt khung đề cương chi tiết; Mốc 2 tại ngày thứ 40 (10/04/2026) thu hồi bản thảo lần một đủ 80 trang ruột; Mốc 3 tại ngày thứ 52 (22/04/2026) ký biên bản kết luận của Hội đồng thẩm định chuyên môn; Mốc 4 tại ngày thứ 60 (30/04/2026) ký duyệt bản thảo sạch để chuyển giao chế bản; Mốc 5 tại ngày thứ 76 (16/05/2026) ký niêm phong bản bông mẫu màu; Mốc 6 tại ngày thứ 84 (24/05/2026) tiếp nhận lô hàng 2.000 cuốn xuất xưởng tại kho; và Mốc 7 tại ngày thứ 90 (30/05/2026) hoàn tất nghiệm thu cấp cơ sở, bàn giao tài liệu cho năm trường trung học cơ sở.

\subsubsection*{4.5.2. Điều hòa nguồn lực và nén tiến độ khi phát sinh chậm trễ}

Nhờ công việc D có thời gian dự trữ độc lập 20 ngày, trong trường hợp bộ phận hành chính và mua sắm cần tập trung hỗ trợ nhóm tác giả tổng hợp dữ liệu trường nghề ở tuần thứ ba, thời điểm phát hành hồ sơ mời báo giá có thể lùi từ ngày thứ 15 sang ngày thứ 25 mà vẫn hoàn tất ký hợp đồng trước ngày thứ 52.

Đối với nhánh công việc găng có thời lượng dài nhất là biên soạn bản thảo (C, 26 ngày), nếu một nhóm tác giả chậm tiến độ tổng hợp số liệu tuyển sinh, Giám đốc dự án sẽ điều động trợ lý nghiên cứu hỗ trợ chuẩn hóa bảng biểu. Đồng thời, dự án áp dụng kỹ thuật \textit{Fast-tracking} (Kỹ thuật thực hiện gối đầu) bằng cách gửi trước các phần bản thảo đã hoàn chỉnh cho ba chuyên gia phản biện đọc độc lập ngay từ ngày thứ 36, giúp rút ngắn thời gian thẩm định (công việc E) từ 12 ngày xuống 9 ngày để bù đắp độ trễ mà không làm tăng chi phí dự toán.


\newpage

% 8. Chương 5: Kế hoạch quản lý chi phí dự án
\section{CHƯƠNG 5. KẾ HOẠCH QUẢN LÝ CHI PHÍ DỰ ÁN}

Chương này vận dụng quy trình quản lý chi phí bốn bước gồm lập kế hoạch nguồn tài nguyên, ước tính chi phí từ dưới lên theo các định mức tài chính hiện hành, thiết lập dự toán ngân sách gắn với đường mức chi phí cơ sở và xác lập cơ chế kiểm soát giải ngân trong phạm vi tổng ngân sách 185.000.000 đồng.

\subsection{5.1. Khung quản lý chi phí và kế hoạch nguồn lực tài chính}

Kế hoạch nguồn lực tài chính của dự án được xác lập dựa trên bốn yếu tố quản trị thực tế: cố định quy mô xuất bản ở mức 2.000 cuốn (80 trang ruột và 4 trang bìa ngoài) để kiểm soát trần chi phí vật tư giấy in; áp dụng định mức thù lao biên soạn, thẩm định theo quy định của Bộ Tài chính đối với nguồn vốn ngân sách nhà nước; kết hợp huy động chuyên gia trong ngành giáo dục cho khâu biên soạn nội dung với việc thuê ngoài trọn gói khâu bình bản kẽm và in ấn công nghiệp; đồng thời trích lập quỹ dự phòng rủi ro chi phí để chủ động trước biến động giá giấy in trên thị trường.

\subsection{5.2. Phương pháp ước tính chi phí theo mô hình từ dưới lên}

\subsubsection{5.2.1. Phân loại ước tính theo các giai đoạn dự án}

Quá trình xác định ngân sách dự án đi qua ba cấp độ ước tính có biên độ sai số thu hẹp dần: cấp độ \textit{Order of Magnitude} (Ước tính độ lớn thô) ở giai đoạn hình thành ý tưởng với biên độ sai số từ $-25\%$ đến $+75\%$ (khoảng 150.000.000 đồng đến 220.000.000 đồng); cấp độ \textit{Budget Estimate} (Ước tính ngân sách) khi lập Điều lệ dự án với biên độ sai số từ $-10\%$ đến $+25\%$, xác lập mức trần ngân sách là 185.000.000 đồng; và cấp độ \textit{Definitive Estimate} (Ước tính xác định) khi lập kế hoạch chi tiết với sai số từ $-5\%$ đến $+10\%$, gắn từng khoản mục chi tiêu với từng gói công việc trong cấu trúc phân rã công việc.

\subsubsection{5.2.2. Phương pháp ước tính chi phí từ dưới lên và căn cứ định mức}

Dự án sử dụng phương pháp \textit{Bottom-up Estimating} (Phương pháp ước tính chi phí từ dưới lên), bóc tách khối lượng và áp đơn giá cho từng gói công việc từ mã \texttt{1.1} đến \texttt{5.4} đã thiết lập tại Chương 3, sau đó tổng hợp thành ngân sách từng giai đoạn.

Đơn giá dự toán được xây dựng dựa trên ba căn cứ: Thông tư số 51/2019/TT-BTC ngày 15 tháng 8 năm 2019 của Bộ Tài chính hướng dẫn xây dựng dự toán, thanh quyết toán kinh phí biên soạn, chỉnh lý và thẩm định tài liệu giáo dục của địa phương; Thông tư số 76/2018/TT-BTC ngày 17 tháng 8 năm 2018 của Bộ Tài chính hướng dẫn mức chi xây dựng chương trình đào tạo, biên soạn giáo trình tài liệu; và báo giá khảo sát thị trường của các doanh nghiệp in xuất bản phẩm đối với gói dịch vụ in offset bốn màu.

\subsection{5.3. Dự toán ngân sách chi tiết theo cấu trúc phân rã công việc}

Bảng 5.1 trình bày dự toán chi tiết cho toàn bộ 20 gói công việc (từ \texttt{1.1} đến \texttt{5.4}) và Quỹ dự phòng rủi ro chi phí (\texttt{6.0}), khớp trực tiếp với Từ điển phân rã công việc tại Bảng 3.1.

\begin{center}
{\small
\setstretch{1.15}
\begin{longtable}{|c|p{4.5cm}|p{5.5cm}|r|c|}
\caption{Bảng dự toán chi tiết ngân sách dự án theo mã phân rã công việc} \label{tab:du_toan_chi_phi} \\
\hline
\textbf{Mã CV} & \textbf{Hạng mục công việc và nội dung chi tiêu} & \textbf{Căn cứ định mức và khối lượng tính toán} & \textbf{Thành tiền (đ)} & \textbf{Tỷ trọng} \\
\hline
\endfirsthead

\multicolumn{5}{c}{\textit{(Tiếp theo Bảng 5.1: Dự toán chi tiết ngân sách dự án)}} \\
\hline
\textbf{Mã CV} & \textbf{Hạng mục công việc và nội dung chi tiêu} & \textbf{Căn cứ định mức và khối lượng tính toán} & \textbf{Thành tiền (đ)} & \textbf{Tỷ trọng} \\
\hline
\endhead

\hline
\multicolumn{5}{r}{\textit{Còn tiếp sang trang sau...}} \\
\endfoot

\hline
\endlastfoot

\textbf{1.0} & \multicolumn{2}{l|}{\textbf{Khảo sát nhu cầu và lập khung cấu trúc}} & \textbf{12.500.000} & \textbf{6,76\%} \\
\hline
1.1a & In ấn phiếu khảo sát nhu cầu thực tế & 350 bộ phiếu $\times$ 3.000 đồng/bộ phiếu & 1.050.000 & -- \\
\hline
1.1b & Thù lao điều tra viên phát và thu phiếu & 5 trường $\times$ 500.000 đồng/trường & 2.500.000 & -- \\
\hline
1.2 & Xử lý số liệu khảo sát và tổng hợp dữ liệu & 2 chuyên viên $\times$ 1.500.000 đồng/người & 3.000.000 & -- \\
\hline
1.3a & Thù lao xây dựng khung đề cương chi tiết 4 phần & Định mức trọn gói theo TT số 51/2019/TT-BTC & 3.500.000 & -- \\
\hline
1.3b & Tổ chức hội thảo chuyên môn thông qua đề cương & Bồi dưỡng 12 đại biểu $\times$ 150.000 đ (1.800.000 đ) và tài liệu, nước uống (650.000 đ) & 2.450.000 & -- \\
\hline
\textbf{2.0} & \multicolumn{2}{l|}{\textbf{Biên soạn bản thảo bộ tài liệu (80 trang ruột)}} & \textbf{32.800.000} & \textbf{17,73\%} \\
\hline
2.1 & Biên soạn Phần 1: Khám phá bản thân (18 trang) & 5 bài chuẩn $\times$ 450.000 đ/bài (2.250.000 đ, TT 51) và chuẩn hóa trắc nghiệm Holland (4.750.000 đ) & 7.000.000 & -- \\
\hline
2.2 & Biên soạn Phần 2: Bản đồ phân luồng (22 trang) & 6 bài chuẩn $\times$ 450.000 đ/bài (2.700.000 đ, TT 51) và lập ma trận so sánh 4 hướng đi (3.300.000 đ) & 6.000.000 & -- \\
\hline
2.3 & Biên soạn Phần 3: Thị trường và chính sách (18 trang) & 5 bài chuẩn $\times$ 450.000 đ/bài (2.250.000 đ, TT 51) và dữ liệu chính sách học phí (3.550.000 đ) & 5.800.000 & -- \\
\hline
2.4 & Biên soạn Phần 4: Cẩm nang quyết định (18 trang) & 4 bài chuẩn $\times$ 450.000 đ/bài (1.800.000 đ, TT 51) và lập phiếu gia đình, danh bạ trường (5.200.000 đ) & 7.000.000 & -- \\
\hline
2.5 & Tổng hợp và biên tập sơ bộ bản thảo lần một & Thù lao Chủ biên 20 bài $\times$ 50.000 đ (1.000.000 đ); biên tập 80 trang $\times$ 75.000 đ (6.000.000 đ) & 7.000.000 & -- \\
\hline
\textbf{3.0} & \multicolumn{2}{l|}{\textbf{Thẩm định chuyên môn và hiệu đính bản thảo}} & \textbf{17.200.000} & \textbf{9,30\%} \\
\hline
3.1 & Thù lao đọc phản biện độc lập (3 chuyên gia) & 80 trang $\times$ 10.000 đồng/trang $\times$ 3 người (Thông tư số 51/2019/TT-BTC) & 2.400.000 & -- \\
\hline
3.2a & Thù lao Chủ tịch, Ủy viên Hội đồng thẩm định & Chủ tịch: 2 buổi $\times$ 200.000 đ (400.000 đ); 4 Ủy viên, Thư ký: 4 người $\times$ 150.000 đ $\times$ 2 buổi (1.200.000 đ) & 1.600.000 & -- \\
\hline
3.2b & Chi phí phục vụ 2 phiên họp Hội đồng thẩm định & In sao bản thảo thẩm định và hậu cần 2 phiên họp $\times$ 1.000.000 đồng/phiên & 2.000.000 & -- \\
\hline
3.3 & Tổ chức lấy ý kiến đọc thử thực tế & 2 nhóm thảo luận tập trung (20 học sinh lớp 9 và 5 giáo viên tư vấn) & 3.200.000 & -- \\
\hline
3.4 & Chỉnh sửa sau thẩm định và hiệu đính bản thảo & Bốn nhóm tác giả chỉnh sửa nội dung và ban biên tập hiệu đính 80 trang bản sạch & 8.000.000 & -- \\
\hline
\textbf{4.0} & \multicolumn{2}{l|}{\textbf{Thiết kế chế bản và in ấn tài liệu}} & \textbf{104.500.000} & \textbf{56,49\%} \\
\hline
4.1 & Tổ chức lựa chọn nhà thầu in ấn công nghiệp & Chào hàng cạnh tranh rút gọn (chi phí hành chính lồng ghép trong mục 6.2) & 0 & -- \\
\hline
4.2a & Thiết kế mỹ thuật 4 trang bìa ngoài và phụ bìa & 4 trang bìa $\times$ 350.000 đ (1.400.000 đ) và trang phụ bìa (100.000 đ) theo TT 51/2019/TT-BTC & 1.500.000 & -- \\
\hline
4.2b & Vẽ minh họa chuyên môn và sơ đồ đồ họa & 25 hình minh họa $\times$ 160.000 đồng/hình (Thông tư số 51/2019/TT-BTC) & 4.000.000 & -- \\
\hline
4.2c & Dàn trang chế bản điện tử chuyên dụng & 80 trang ruột $\times$ 50.000 đồng/trang kỹ thuật & 4.000.000 & -- \\
\hline
4.3 & In thử nghiệm 2 bộ bản bông màu và rà soát lỗi & Chi phí in thử màu kỹ thuật số và kiểm tra đối soát mẫu kỹ thuật & 1.000.000 & -- \\
\hline
4.4 & Thuê ngoài in offset hàng loạt 2.000 cuốn & 2.000 cuốn $\times$ 47.000 đồng/cuốn (giá trúng thầu trọn gói, dự toán gói thầu là 95.000.000 đ) & 94.000.000 & -- \\
\hline
\textbf{5.0} & \multicolumn{2}{l|}{\textbf{Kiểm định, nghiệm thu và bàn giao}} & \textbf{8.500.000} & \textbf{4,59\%} \\
\hline
5.1 & Thù lao kiểm định mẫu thống kê tại kho & Đo đạc và kiểm tra quy cách 384 cuốn mẫu ngẫu nhiên theo tiêu chuẩn & 2.000.000 & -- \\
\hline
5.2 & Tổ chức họp Hội đồng nghiệm thu cấp cơ sở & Bồi dưỡng hội đồng nghiệm thu và hậu cần phiên họp đánh giá kết quả & 2.500.000 & -- \\
\hline
5.3a & Mua thùng giấy tiêu chuẩn và đóng gói & 40 thùng (50 cuốn/thùng) $\times$ 25.000 đồng/thùng & 1.000.000 & -- \\
\hline
5.3b & Chi phí vận chuyển và bàn giao đến 5 trường & Thuê phương tiện vận chuyển bàn giao tận kho 5 trường trung học cơ sở & 3.000.000 & -- \\
\hline
5.4 & Thanh lý hợp đồng và đóng gói 03 bộ hồ sơ lưu trữ & Thực hiện nghiệp vụ thanh quyết toán và lưu trữ hành chính & 0 & -- \\
\hline
\textbf{6.0} & \multicolumn{2}{l|}{\textbf{Quỹ dự phòng rủi ro chi phí}} & \textbf{9.500.000} & \textbf{5,14\%} \\
\hline
6.1 & Dự phòng rủi ro biến động giá vật tư giấy in & Dự phòng trượt giá giấy và vật tư xuất bản công nghiệp & 6.000.000 & -- \\
\hline
6.2 & Dự phòng chi phí hành chính và chỉnh lý kỹ thuật & Chi phí in ấn hồ sơ mời báo giá, lưu trữ và chỉnh lý phát sinh & 3.500.000 & -- \\
\hline
\multicolumn{3}{|l|}{\textbf{Tổng cộng toàn bộ ngân sách dự án (Cân đối Điều lệ)}} & \textbf{185.000.000} & \textbf{100,0\%} \\
\hline
\end{longtable}
}
\end{center}

Trong tổng ngân sách 185.000.000 đồng, khối lượng công việc học thuật từ khảo sát đến hoàn thiện bản thảo sạch (Mục 1.0, 2.0 và 3.0) chiếm 33,78\% (62.500.000 đồng); khối lượng thiết kế chế bản và in ấn (Mục 4.0) chiếm 56,49\% (104.500.000 đồng, gồm 10.500.000 đồng thiết kế dàn trang, in thử và 94.000.000 đồng hợp đồng in công nghiệp); công tác kiểm định, nghiệm thu bàn giao (Mục 5.0) chiếm 4,59\% (8.500.000 đồng); và khoản \textit{Contingency Fund} (Quỹ dự phòng rủi ro chi phí -- Mục 6.0) chiếm 5,14\% (9.500.000 đồng).

\subsection{5.4. Đường mức chi phí cơ sở và quy trình kiểm soát giải ngân}

\subsubsection{5.4.1. Đường mức chi phí cơ sở gắn với năm đợt giải ngân theo mốc nghiệm thu}

Đường \textit{Cost Baseline} (Đường mức chi phí cơ sở) của dự án được thiết lập dựa trên nguyên tắc chỉ giải ngân khi gói công việc tại mốc sự kiện găng tương ứng đã có biên bản nghiệm thu kỹ thuật. Tiến trình giải ngân lũy kế từ ngày khởi công đến ngày kết thúc dự án được thể hiện trên biểu đồ Hình~\ref{fig:cost_baseline} và phân kỳ qua năm đợt thanh toán:

\begin{figure}[htbp]
\centering
\begin{tikzpicture}[x=2.2cm,y=0.032cm]
    % Trục tọa độ
    \draw[->, thick, >=stealth] (0,0) -- (6.5,0) node[right, font=\small\bfseries] {Mốc thời gian};
    \draw[->, thick, >=stealth] (0,0) -- (0,210) node[above, font=\small\bfseries] {Chi phí lũy kế (triệu đồng)};
    
    % Lưới ngang
    \foreach \y in {20, 40, 60, 80, 100, 120, 140, 160, 180} {
        \draw[dotted, gray!50] (0,\y) -- (6.2,\y);
        \node[left, font=\scriptsize] at (0,\y) {\y};
    }
    \node[left, font=\scriptsize] at (0,0) {0};

    % Các cột biểu diễn chi phí lũy kế
    \fill[blue!20, draw=blue!60, line width=0.8pt] (0.8,0) rectangle (1.2,12.5);
    \fill[blue!20, draw=blue!60, line width=0.8pt] (1.8,0) rectangle (2.2,45.3);
    \fill[blue!20, draw=blue!60, line width=0.8pt] (2.8,0) rectangle (3.2,62.5);
    \fill[blue!20, draw=blue!60, line width=0.8pt] (3.8,0) rectangle (4.2,101.2);
    \fill[blue!20, draw=blue!60, line width=0.8pt] (4.8,0) rectangle (5.2,167.0);
    \fill[blue!20, draw=blue!60, line width=0.8pt] (5.8,0) rectangle (6.2,185.0);

    % Đường nối Cost Baseline
    \draw[red!80!black, line width=1.5pt] 
        (1,12.5) -- (2,45.3) -- (3,62.5) -- (4,101.2) -- (5,167.0) -- (6,185.0);

    % Các điểm nút và nhãn giá trị
    \foreach \x/\y/\txt in {
        1/12.5/{12,5 tr}, 
        2/45.3/{45,3 tr}, 
        3/62.5/{62,5 tr}, 
        4/101.2/{101,2 tr}, 
        5/167.0/{167,0 tr}, 
        6/185.0/{185,0 tr}
    } {
        \fill[red!80!black] (\x,\y) circle (2.5pt);
        \node[above, font=\footnotesize\bfseries, text=red!80!black] at (\x,\y+3) {\txt};
    }

    % Nhãn trục hoành
    \node[below=3pt, align=center, font=\scriptsize] at (1,0) {Đợt 1\\(Ngày 14)};
    \node[below=3pt, align=center, font=\scriptsize] at (2,0) {Đợt 2\\(Ngày 40)};
    \node[below=3pt, align=center, font=\scriptsize] at (3,0) {Đợt 3\\(Ngày 60)};
    \node[below=3pt, align=center, font=\scriptsize] at (4,0) {Đợt 4\\(Ngày 76)};
    \node[below=3pt, align=center, font=\scriptsize] at (5,0) {Mốc in\\(Ngày 84)};
    \node[below=3pt, align=center, font=\scriptsize] at (6,0) {Đợt 5\\(Ngày 90)};
\end{tikzpicture}
\caption{Đường mức chi phí tích lũy theo mốc nghiệm thu kỹ thuật của dự án}
\label{fig:cost_baseline}
\end{figure}

Đợt giải ngân thứ nhất tại Ngày thứ 14 (Tuần 2): Giải ngân \textbf{12.500.000 đồng} (tương ứng 6,76\% ngân sách, lũy kế đạt \textbf{12,5 triệu đồng}) để thanh toán chi phí khảo sát 350 bộ phiếu, xử lý dữ liệu và hội thảo phê duyệt khung đề cương (Mục 1.0) sau khi đề cương chi tiết được thông qua.

Đợt giải ngân thứ hai tại Ngày thứ 40 (Tuần 6): Giải ngân \textbf{32.800.000 đồng} (tương ứng 17,73\% ngân sách, nâng lũy kế lên \textbf{45,3 triệu đồng}, đạt 24,49\%) để chi trả thù lao viết 20 bài chuẩn, xây dựng công cụ trắc nghiệm và biên tập sơ bộ (Mục 2.0) khi Chủ biên bàn giao đủ 80 trang bản thảo lần một.

Đợt giải ngân thứ ba tại Ngày thứ 60 (Tuần 9): Giải ngân \textbf{17.200.000 đồng} (tương ứng 9,30\% ngân sách, nâng lũy kế lên \textbf{62,5 triệu đồng}, đạt 33,78\%) để thanh toán thù lao phản biện độc lập, chi phí họp Hội đồng thẩm định, nhóm đọc thử và thù lao hiệu đính (Mục 3.0) sau khi Giám đốc dự án ký duyệt bản thảo sạch.

Đợt giải ngân thứ tư tại Ngày thứ 76 (Tuần 11): Giải ngân \textbf{38.700.000 đồng} (tương ứng 20,92\% ngân sách, nâng lũy kế lên \textbf{101,2 triệu đồng}, đạt 54,70\%) ngay sau khi ký niêm phong mẫu in thử bản bông, bao gồm thanh toán 10.500.000 đồng chi phí thiết kế bìa, vẽ minh họa, dàn trang và in thử (\texttt{4.2a} đến \texttt{4.3}) cùng khoản tạm ứng 30\% giá trị hợp đồng in ấn cho nhà thầu ($30\% \times 94.000.000 = 28.200.000$ đồng) để nhập vật tư giấy in.

Đợt thanh quyết toán thứ năm từ Ngày thứ 84 đến Ngày thứ 90 (Tuần 12 và Tuần 13): Tại ngày thứ 84 khi nhà thầu bàn giao 2.000 cuốn tại kho, giá trị khối lượng hoàn thành cộng thêm 70\% hợp đồng in còn lại (65.800.000 đồng) đưa chi phí tích lũy bốn giai đoạn đầu lên \textbf{167,0 triệu đồng} (90,27\%). Sau khi kiểm định mẫu 384 cuốn đạt yêu cầu, dự án chuyển khoản thanh toán nốt 65.800.000 đồng cho nhà thầu in vào ngày thứ 88, đồng thời giải ngân \textbf{8.500.000 đồng} chi phí kiểm định, họp nghiệm thu cấp cơ sở, đóng thùng vận chuyển bàn giao (Mục 5.0) và quyết toán \textbf{9.500.000 đồng} quỹ dự phòng (Mục 6.0) vào ngày thứ 90, khép lại tổng mức ngân sách \textbf{185.000.000 đồng} (100\%).

\subsection{5.4.2. Cơ chế quản lý và sử dụng quỹ dự phòng chi phí}

Khoản quỹ dự phòng rủi ro trị giá 9.500.000 đồng (Mục 6.0) thuộc thẩm quyền phê duyệt của Trưởng phòng Giáo dục và Đào tạo (Chủ đầu tư). Quỹ này chỉ được phép giải ngân khi phát sinh biến động giá giấy in công nghiệp trên thị trường vượt quá 10\% tại thời điểm ký hợp đồng hoặc khi Hội đồng thẩm định yêu cầu bổ sung thêm các bảng biểu đồ họa ngoài dự toán ban đầu. Khi cần sử dụng, Giám đốc dự án phải lập tờ trình thuyết minh nguyên nhân và dự toán chi tiết trình Chủ đầu tư phê duyệt bằng văn bản trước khi chi tiêu; phần kinh phí dự phòng không sử dụng đến được hoàn trả ngân sách nhà nước khi quyết toán đóng dự án.


\newpage

% 9. Chương 6: Kế hoạch quản lý chất lượng dự án
\section{CHƯƠNG 6. KẾ HOẠCH QUẢN LÝ CHẤT LƯỢNG DỰ ÁN}

Chương này thiết lập bộ tiêu chuẩn chất lượng kép cho sản phẩm xuất bản, xác định quy trình đảm bảo chất lượng ở giai đoạn biên soạn và chế bản, áp dụng công thức tính kích cỡ mẫu thống kê trong kiểm tra chất lượng nghiệm thu và bóc tách cơ cấu năm nhóm chi phí chất lượng của dự án.

\subsection{6.1. Khung khái niệm và nguyên tắc quản lý chất lượng}

Quản lý chất lượng trong dự án phát triển học liệu hướng nghiệp được vận hành theo hai nguyên tắc nền tảng: nguyên tắc \textit{Fitness for requirements} (Nguyên tắc yêu cầu phù hợp), đòi hỏi ấn phẩm tuân thủ đúng yêu cầu cần đạt của môn Hoạt động trải nghiệm, hướng nghiệp lớp 9 thuộc Chương trình Giáo dục phổ thông 2018 (Thông tư số 32/2018/TT-BGDĐT) và các quy định của Luật Xuất bản năm 2012; cùng nguyên tắc \textit{Fitness for use} (Nguyên tắc tiện lợi cho sử dụng), yêu cầu ngôn ngữ trình bày, cấu trúc bảng hỏi trắc nghiệm và hệ thống sơ đồ phải trực quan, vừa sức tiếp nhận của học sinh độ tuổi 14 đến 15 và phụ huynh.

Chu trình quản lý chất lượng được thực hiện liên hoàn qua ba bước: lập kế hoạch chất lượng, thực hiện \textit{Quality Assurance -- QA} (Đảm bảo chất lượng) trong quá trình tạo lập sản phẩm và tiến hành \textit{Quality Control -- QC} (Kiểm tra chất lượng) khi nghiệm thu lô hàng xuất bản (Hình~\ref{fig:quy_trinh_chat_luong}).

\begin{figure}[htbp]
\centering
\begin{tikzpicture}[
    box/.style={rectangle, draw=blue!80, fill=blue!5, thick, rounded corners=3pt, align=center, inner sep=6pt, font=\small},
    arr/.style={->, thick, >=stealth, draw=blue!80}
]
    \node (step1) [box] {\textbf{6.1. LẬP KẾ HOẠCH CHẤT LƯỢNG}\\\textit{Bộ tiêu chuẩn kép: Học thuật sư phạm}\\\textit{và Kỹ thuật in ấn công nghiệp}};
    \node (step2) [box, right=0.8cm of step1] {\textbf{6.2. ĐẢM BẢO CHẤT LƯỢNG}\\\textit{Thẩm định 2 vòng kín, đọc thử nhóm}\\\textit{và ký duyệt niêm phong bản bông}};
    \node (step3) [box, right=0.8cm of step2] {\textbf{6.3. KIỂM TRA CHẤT LƯỢNG}\\\textit{Rút mẫu phân tầng $n = 384$ cuốn,}\\\textit{đo độ lệch chuẩn và nghiệm thu}};

    \draw [arr] (step1) -- (step2);
    \draw [arr] (step2) -- (step3);
\end{tikzpicture}
\caption{Quy trình ba bước quản lý chất lượng toàn diện của dự án}
\label{fig:quy_trinh_chat_luong}
\end{figure}

\subsection{6.2. Bộ tiêu chuẩn chất lượng sản phẩm xuất bản}

\subsubsection{6.2.1. Tiêu chuẩn về nội dung học thuật và sư phạm}

Nội dung bốn phần của bộ tài liệu phải bảo đảm độ chính xác về pháp lý và dữ liệu tuyển sinh: phản ánh đúng thông tin bốn hướng đi sau trung học cơ sở, danh bạ 25 cơ sở giáo dục nghề nghiệp trên địa bàn và chính sách miễn 100\% học phí học nghề theo Nghị định số 81/2021/NĐ-CP và Nghị định số 97/2023/NĐ-CP.

Đối với công cụ trắc nghiệm sở thích nghề nghiệp theo mật mã Holland rút gọn (36 chỉ báo hành vi tại Phần 1), thang đo phải đạt hệ số tin cậy nội tại từ 0,80 trở lên theo chuẩn thích ứng tâm lý học tại Việt Nam, đi kèm hướng dẫn giải mã kết quả mang tính định hướng mở, không dán nhãn hay áp đặt năng lực của học sinh.

\subsubsection{6.2.2. Tiêu chuẩn về kỹ thuật chế bản và in ấn công nghiệp}

Trên cơ sở quy cách sản phẩm đã xác lập tại Chương 3 (khổ 17 $\times$ 24 cm, gồm 80 trang ruột giấy Fort 100 gsm và 4 trang bìa ngoài giấy Couche 250 gsm cán màng mờ), bộ tiêu chuẩn kỹ thuật nghiệm thu ấn định các ngưỡng dung sai cụ thể: sai số kích thước thành phẩm sau cắt xén không vượt quá $\pm 1\text{ mm}$; độ trắng của giấy ruột đạt từ 84\% ISO trở lên để chống thấu mực và chống lóa; độ phân giải của toàn bộ 25 hình minh họa và đồ họa thông tin đạt tối thiểu 300 DPI; bản in offset bốn màu sắc nét, chồng màu chuẩn xác, không lem mực hoặc nhảy trang; gáy sách phay rãnh và vào keo nhiệt chuyên dụng bảo đảm góc mở phẳng $180^\circ$ không bong tróc khi lật mở nhiều lần.

\subsection{6.3. Quy trình đảm bảo chất lượng học thuật và kỹ thuật}

Hoạt động đảm bảo chất lượng tập trung vào bốn chốt phòng ngừa sai lỗi trước khi đưa vào in ấn hàng loạt:

Trong giai đoạn viết bản thảo, bốn nhóm tác giả thực hiện cơ chế đọc chéo giữa chuyên gia tâm lý và chuyên gia giáo dục học, sau đó Chủ biên biên tập đồng nhất thuật ngữ trên toàn bộ 80 trang ruột (gói công việc \texttt{2.5}).

Khi hoàn thành bản thảo lần một, dự án thực hiện thẩm định chuyên môn qua hai vòng độc lập: gửi toàn văn bản thảo cho ba chuyên gia ngoài dự án đọc phản biện kín theo phiếu tiêu chí (gói công việc \texttt{3.1}), tiếp theo là tổ chức hai phiên họp Hội đồng thẩm định gồm năm thành viên để đánh giá báo cáo giải trình tiếp thu (gói công việc \texttt{3.2}); bản thảo chỉ được thông qua khi đạt tối thiểu 4/5 phiếu đồng thuận.

Song song với thẩm định chuyên gia, dự án tổ chức khảo nghiệm thực tế theo phương pháp \textit{Focus Group} (Thảo luận nhóm tập trung) với 20 học sinh lớp 9 và 5 giáo viên tư vấn (gói công việc \texttt{3.3}) để rà soát các từ ngữ khó hiểu hoặc cách tính điểm trắc nghiệm còn gây nhầm lẫn, làm cơ sở hiệu đính bản thảo sạch (gói công việc \texttt{3.4}).

Ở khâu chế bản, sau khi hoàn tất dàn trang điện tử, bộ phận kỹ thuật và nhà thầu phải in thử nghiệm hai bộ bản bông màu (gói công việc \texttt{4.3}). Giám đốc dự án và Chủ biên trực tiếp kiểm tra đối soát từng trang và ký niêm phong lên bản bông làm mẫu đối chứng pháp lý duy nhất cho khâu in offset hàng loạt.

\subsection{6.4. Quy trình kiểm tra chất lượng theo phương pháp lấy mẫu thống kê}

\subsubsection{6.4.1. Công thức và bài toán tính kích cỡ mẫu thống kê}

Để kiểm định khách quan lô hàng 2.000 cuốn tài liệu khi xuất xưởng mà không cần kiểm tra toàn bộ 100\% sản phẩm, dự án áp dụng công thức xác định kích cỡ mẫu thống kê:

\begin{equation}
n = 0,25 \times \left(\frac{Z}{e}\right)^2
\end{equation}

Trong đó, $Z$ là hệ số tra bảng tương ứng với mức độ tin cậy của phép lấy mẫu (theo bảng chuẩn trong quản lý chất lượng: độ tin cậy 80\% ứng với $Z = 1,281$; độ tin cậy 90\% ứng với $Z = 1,645$; độ tin cậy 95\% ứng với $Z = 1,960$). Dự án lựa chọn mức độ tin cậy 95\% ($Z = 1,960$) và biên độ sai số chấp nhận được là $e = 5\% = 0,05$. Thay số vào công thức:

\begin{equation}
n = 0,25 \times \left(\frac{1,960}{0,05}\right)^2 = 0,25 \times (39,2)^2 = 0,25 \times 1.536,64 = 384,16
\end{equation}

Làm tròn theo quy chuẩn thống kê, kích cỡ mẫu kiểm tra chất lượng là \textbf{$n = 384\text{ cuốn}$}.

\subsubsection{6.4.2. Quy trình rút mẫu phân tầng và ngưỡng chấp nhận nghiệm thu}

Khi nhà thầu giao 2.000 cuốn đóng trong 40 thùng giấy tiêu chuẩn (50 cuốn/thùng), Tiểu ban Kiểm định tiến hành rút mẫu phân tầng ngẫu nhiên bằng cách lấy từ mỗi thùng từ 9 đến 10 cuốn ở cả ba lớp (trên, giữa, đáy thùng) cho đủ 384 cuốn mẫu.

Trên 384 cuốn mẫu, nhóm kiểm định đo kích thước cạnh xén (yêu cầu độ lệch chuẩn $\sigma \le 0,5\text{ mm}$, nằm trong giới hạn kiểm soát $\pm 3\sigma$), kiểm tra thứ tự trang từ trang 1 đến trang 80, thử nghiệm mở phẳng $180^\circ$ kiểm tra keo gáy và đối chiếu màu sắc với bản bông đã niêm phong.

Quy tắc nghiệm thu lô hàng được ấn định theo ba mức cụ thể:
\begin{itemize}
    \item Nếu số cuốn lỗi trong mẫu 384 cuốn từ 1 cuốn trở xuống (tỷ lệ lỗi $\le 0,26\%$), lô hàng 2.000 cuốn đạt yêu cầu và được ký biên bản nghiệm thu;
    \item Nếu phát hiện từ 2 đến 4 cuốn lỗi, nhà thầu phải rà soát lại toàn bộ 40 thùng hàng và in bù đổi mới các cuốn lỗi trong vòng 48 giờ;
    \item Nếu phát hiện từ 5 cuốn lỗi trở lên (tỷ lệ lỗi $> 1,30\%$), Hội đồng nghiệm thu từ chối tiếp nhận và yêu cầu nhà thầu tái sản xuất toàn bộ lô hàng bằng kinh phí của nhà thầu.
\end{itemize}

\subsection{6.5. Cơ cấu chi phí chất lượng của dự án}

Từ bảng dự toán ngân sách tại Chương 5 (Bảng 5.1), các khoản chi phục vụ trực tiếp cho công tác quản lý chất lượng được phân nhóm theo năm cấu phần chi phí chất lượng tại Bảng 6.1 (đảm bảo truy vết 1-1 theo mã công việc và không tính trùng giữa các nhóm):

\begin{table}[htbp]
\centering
\caption{Phân tích năm nhóm chi phí chất lượng của dự án}
\label{tab:chi_phi_chat_luong}
\small
\setstretch{1.15}
\begin{tabularx}{\textwidth}{|p{3.2cm}|Y|Y|r|c|}
\hline
\textbf{Nhóm chi phí} & \textbf{Nội hàm quản lý} & \textbf{Các gói công việc cấu thành (Bảng 5.1)} & \textbf{Chi phí (đ)} & \textbf{Tỷ trọng} \\
\hline
\textbf{1. Chi phí ngăn ngừa} & Hoạch định và thiết kế chuẩn đầu vào nhằm ngăn ngừa sai lỗi ngay từ đầu. & Khảo sát thực tế, phân tích dữ liệu, xây dựng và phê duyệt khung đề cương (\texttt{1.1a, 1.1b, 1.2, 1.3a, 1.3b}). & \textbf{12.500.000} & 6,76\% \\
\hline
\textbf{2. Chi phí đánh giá} & Phản biện, thẩm định chuyên môn và duyệt mẫu kỹ thuật đánh giá sự phù hợp. & Phản biện độc lập (\texttt{3.1}: 2,4 tr); họp Hội đồng thẩm định (\texttt{3.2a, 3.2b}: 3,6 tr); duyệt bản bông (\texttt{4.3}: 1,0 tr); họp nghiệm thu (\texttt{5.2}: 2,5 tr). & \textbf{9.500.000} & 5,14\% \\
\hline
\textbf{3. Chi phí sai sót bên trong} & Rà soát, chỉnh sửa và khắc phục các khiếm khuyết nội dung trước xuất bản. & Biên tập đồng nhất văn phong 80 trang bản thảo 1 (\texttt{2.5} phần biên tập: 6,0 tr); chỉnh sửa sau thẩm định (\texttt{3.4}: 8,0 tr). & \textbf{14.000.000} & 7,57\% \\
\hline
\textbf{4. Chi phí sai sót bên ngoài} & Thu hồi, in đền bù nếu sản phẩm đã bàn giao phát sinh lỗi nghiêm trọng. & Kiểm soát tuyệt đối qua chốt duyệt bản bông và kiểm định mẫu thống kê tại kho trước khi nhận. & \textbf{0} & 0,00\% \\
\hline
\textbf{5. Chi phí thử nghiệm, đo lường} & Khảo nghiệm thực tế trên người dùng và đo đạc mẫu kiểm định thống kê. & Đọc thử nghiệm với học sinh, giáo viên (\texttt{3.3}: 3,2 tr); đo đạc kiểm định 384 cuốn tại kho (\texttt{5.1}: 2,0 tr). & \textbf{5.200.000} & 2,81\% \\
\hline
\multicolumn{3}{|l|}{\textbf{Tổng cộng chi phí chất lượng toàn diện của dự án}} & \textbf{41.200.000} & \textbf{22,27\%} \\
\hline
\end{tabularx}
\end{table}

Tổng chi phí chất lượng của dự án là \textbf{41.200.000 đồng}, chiếm \textbf{22,27\%} tổng ngân sách 185.000.000 đồng. Việc phân bổ nguồn lực thỏa đáng cho khâu ngăn ngừa, đánh giá và thử nghiệm đo lường (tổng cộng 27.200.000 đồng) cùng khâu biên tập, hiệu đính nội bộ (14.000.000 đồng) giúp dự án loại bỏ hoàn toàn rủi ro phát sinh chi phí sai sót bên ngoài sau khi bàn giao ấn phẩm cho các nhà trường.


\newpage

% 10. Chương 7: Kế hoạch quản lý nguồn nhân lực dự án
\section{CHƯƠNG 7. KẾ HOẠCH QUẢN LÝ NGUỒN NHÂN LỰC DỰ ÁN}

Chương này thiết lập sơ đồ cơ cấu tổ chức dự án theo mô hình ma trận cân bằng, phân định trách nhiệm của từng vị trí bằng ma trận phân công trách nhiệm, đồng thời vận dụng thuyết thứ bậc nhu cầu, mười phẩm chất tâm lý của người lãnh đạo và bảy thói quen quản lý hiệu quả vào điều hành đội ngũ nhân sự dự án.

\subsection{7.1. Cơ cấu tổ chức và bộ máy quản lý dự án}

Ban Quản lý dự án được tổ chức theo mô hình ma trận chức năng cân bằng nhằm huy động đội ngũ chuyên gia tâm lý học, giáo dục học từ các trường đại học, viện nghiên cứu phối hợp cùng đội ngũ cán bộ quản lý của Phòng Giáo dục và Đào tạo và các trường trung học cơ sở mà không làm xáo trộn biên chế thường xuyên của các đơn vị.

\subsubsection{7.1.1. Sơ đồ cơ cấu tổ chức Ban Quản lý dự án}

Bộ máy dự án gồm bốn tiểu ban chuyên trách đặt dưới sự điều phối trực tiếp của Giám đốc dự án, trong đó Tiểu ban Chuyên môn nội dung quản lý bốn nhóm tác giả phụ trách bốn phần của bộ tài liệu (Hình~\ref{fig:co_cau_to_chuc}):

\begin{figure}[htbp]
\centering
\begin{tikzpicture}[
    level 1/.style={sibling distance=3.6cm, level distance=1.4cm},
    level 2/.style={sibling distance=1.8cm, level distance=1.8cm},
    box/.style={rectangle, draw=blue!80, fill=blue!5, thick, rounded corners=2pt, align=center, inner sep=4pt, font=\scriptsize},
    root/.style={rectangle, draw=red!80, fill=red!10, thick, rounded corners=3pt, align=center, inner sep=5pt, font=\footnotesize\bfseries},
    manager/.style={rectangle, draw=blue!80, fill=blue!15, thick, rounded corners=3pt, align=center, inner sep=5pt, font=\scriptsize\bfseries},
    subbox/.style={rectangle, draw=gray!80, fill=gray!10, thin, rounded corners=2pt, align=center, inner sep=3pt, font=\tiny}
]
    \node (CDT) [root] {LÃNH ĐẠO PHÒNG GIÁO DỤC VÀ ĐÀO TẠO\\(Cơ quan phê duyệt và Chủ đầu tư)}
        child {
            node (GD) [manager] {GIÁM ĐỐC DỰ ÁN\\ThS. Lê Hoàng Nam (Phòng GD\&ĐT)}
            child {
                node (CB) [box] {\textbf{Tiểu ban Chuyên môn}\\\textit{Chủ biên: PGS.TS. Nguyễn Minh Hà}}
                child { node [subbox] {4 Nhóm tác giả\\(Phần 1 đến Phần 4)} }
            }
            child {
                node (HD) [box] {\textbf{Hội đồng Thẩm định}\\\textit{Chủ tịch: TS. Vũ Thị Mai}}
                child { node [subbox] {3 Chuyên gia phản biện\\và Thư ký thẩm định} }
            }
            child {
                node (MS) [box] {\textbf{Tiểu ban Mua sắm \& TC}\\\textit{Phụ trách: ThS. Phạm Quang Thắng}}
                child { node [subbox] {Kế toán, Thiết kế\\và Nhà thầu in} }
            }
            child {
                node (TH) [box] {\textbf{Tiểu ban Kiểm định}\\\textit{Đại diện: ThS. Đỗ Thị Lan}}
                child { node [subbox] {Tổ giáo viên 5 trường\\và Nhóm đo đạc mẫu} }
            }
        };
\end{tikzpicture}
\caption{Sơ đồ cơ cấu tổ chức Ban Quản lý dự án theo mô hình ma trận cân bằng}
\label{fig:co_cau_to_chuc}
\end{figure}

\subsubsection{7.1.2. Chức năng và nhiệm vụ của các tiểu ban chuyên trách}

Trên cơ sở danh sách nhân sự chủ chốt đã xác lập tại Bảng 2.1 (Chương 2), phạm vi điều hành của bốn tiểu ban được quy định cụ thể theo từng mảng nghiệp vụ:

Tiểu ban Chuyên môn nội dung (do Chủ biên Phó Giáo sư, Tiến sĩ Nguyễn Minh Hà phụ trách) chịu trách nhiệm thiết kế khung đề cương 20 bài chuẩn, điều phối bốn nhóm tác giả biên soạn 80 trang ruột, tổ chức đọc chéo nội bộ, biên tập đồng nhất văn phong sư phạm và trực tiếp chỉnh sửa, hiệu đính bản thảo sạch theo kết luận của Hội đồng thẩm định.

Hội đồng thẩm định độc lập (do Tiến sĩ Vũ Thị Mai làm Chủ tịch) chịu trách nhiệm điều phối ba chuyên gia phản biện độc lập đánh giá toàn văn bản thảo, chủ trì hai phiên họp thẩm định chính thức và giám sát mức độ tiếp thu chỉnh sửa trước khi cấp phép chế bản.

Tiểu ban Mua sắm và Tài chính (do Thạc sĩ Phạm Quang Thắng phụ trách) điều phối nhóm họa sĩ thiết kế bìa, vẽ 25 hình minh họa và dàn trang điện tử; đồng thời tổ chức quy trình chào hàng cạnh tranh rút gọn để lựa chọn nhà thầu in ấn, kiểm soát chứng từ giải ngân năm đợt theo mốc nghiệm thu và thực hiện thủ tục thanh quyết toán, lưu trữ hồ sơ.

Tiểu ban Kiểm định và Thụ hưởng (do Thạc sĩ Đỗ Thị Lan đại diện) điều phối mạng lưới giáo viên tư vấn tại năm trường trung học cơ sở để thu thập 350 bộ phiếu khảo sát đầu vào, tổ chức hai nhóm đọc thử bản thảo (20 học sinh và 5 giáo viên), chủ trì rút mẫu phân tầng 384 cuốn để kiểm định chất lượng kỹ thuật tại kho và tổ chức tiếp nhận, phân phối 2.000 cuốn ấn phẩm về các trường.

\subsection{7.2. Ma trận phân bổ trách nhiệm}

Dự án áp dụng mô hình \textit{RACI} (\textit{Responsible} -- Người trực tiếp thực hiện, \textit{Accountable} -- Người chịu trách nhiệm phê duyệt, \textit{Consulted} -- Người được tham vấn, \textit{Informed} -- Người được thông báo) đối với mười công việc từ A đến J trên sơ đồ mạng. Nguyên tắc cốt lõi được duy trì xuyên suốt là mỗi công việc chỉ có duy nhất một vị trí giữ vai trò \texttt{A} để tránh chồng chéo thẩm quyền hoặc đùn đẩy trách nhiệm khi nghiệm thu.

\begin{table}[htbp]
\centering
\caption{Ma trận phân bổ trách nhiệm theo mô hình RACI}
\label{tab:raci}
\small
\setstretch{1.15}
\begin{tabularx}{\textwidth}{|C{0.6cm}|Y|C{0.9cm}|C{0.9cm}|C{0.9cm}|C{0.9cm}|C{0.9cm}|C{0.9cm}|C{0.9cm}|C{0.9cm}|}
\hline
\textbf{Mã} & \textbf{Tên công việc trên sơ đồ mạng} & \textbf{LĐ Phòng} & \textbf{GĐ DA} & \textbf{Chủ biên} & \textbf{Tác giả} & \textbf{HĐ Thẩm} & \textbf{Mua sắm} & \textbf{In \& TK} & \textbf{Trường cơ sở} \\
\hline
\textbf{A} & Khảo sát nhu cầu và thu thập dữ liệu & I & \textbf{A} & C & \textbf{R} & I & I & I & C \\
\hline
\textbf{B} & Lập khung cấu trúc và duyệt đề cương & \textbf{A} & C & \textbf{R} & C & C & I & I & I \\
\hline
\textbf{C} & Biên soạn bản thảo lần một (80 trang) & I & \textbf{A} & \textbf{R} & \textbf{R} & I & I & I & I \\
\hline
\textbf{D} & Chọn nhà thầu in và ký hợp đồng & C & \textbf{A} & I & I & I & \textbf{R} & C & I \\
\hline
\textbf{E} & Đọc phản biện và họp thẩm định & I & \textbf{A} & C & I & \textbf{R} & I & I & C \\
\hline
\textbf{F} & Tiếp thu ý kiến, hiệu đính bản sạch & I & \textbf{A} & \textbf{R} & C & C & I & I & I \\
\hline
\textbf{G} & Thiết kế đồ họa và dàn trang chế bản & I & \textbf{A} & C & I & I & C & \textbf{R} & I \\
\hline
\textbf{H} & In bản bông và ký duyệt mẫu in thử & I & \textbf{A} & \textbf{R} & I & I & C & C & I \\
\hline
\textbf{I} & In offset hàng loạt 2.000 cuốn & I & \textbf{A} & I & I & I & C & \textbf{R} & I \\
\hline
\textbf{J} & Kiểm định mẫu, nghiệm thu, bàn giao & \textbf{A} & \textbf{R} & C & I & C & C & C & C \\
\hline
\end{tabularx}
\end{table}

\subsection{7.3. Động lực làm việc và nghệ thuật lãnh đạo nhóm dự án}

\subsubsection{7.3.1. Ứng dụng thuyết thứ bậc nhu cầu trong tạo động lực nhân sự}

Do đội ngũ tham gia dự án gồm cả các giảng viên đại học, nghiên cứu viên và cán bộ quản lý cơ sở có đặc thù công việc kiêm nhiệm, Giám đốc dự án vận dụng thuyết thứ bậc nhu cầu của Abraham Maslow để thiết kế cơ chế khuyến khích phù hợp với từng nhóm nhân sự:

Ở tầng nhu cầu cơ bản và nhu cầu an toàn, dự án bảo đảm chi trả thù lao viết bài, phản biện và dàn trang đúng định mức tại Thông tư số 51/2019/TT-BTC, giải ngân kịp thời ngay sau từng mốc nghiệm thu để các thành viên yên tâm công tác. Ở tầng nhu cầu xã hội và sự gắn kết, Ban Quản lý dự án duy trì không khí trao đổi học thuật cởi mở trong các buổi hội thảo đề cương và thảo luận nhóm, kết nối góc nhìn lý luận của nhà nghiên cứu với kinh nghiệm thực tiễn của giáo viên trung học cơ sở.

Ở tầng nhu cầu được tôn trọng và tự khẳng định bản thân, học hàm, học vị và họ tên của Chủ biên, từng tác giả biên soạn bốn phần và các ủy viên Hội đồng thẩm định được ghi nhận trang trọng tại trang thông tin xuất bản của cuốn sách. Đồng thời, giá trị thực tiễn của bộ tài liệu khi đến tay 2.000 học sinh lớp 9 và các gia đình trên địa bàn chính là động lực nghề nghiệp bền vững đối với đội ngũ những người làm giáo dục.

\subsubsection{7.3.2. Mười phẩm chất tâm lý của người lãnh đạo tài năng}

Trong điều hành tác nghiệp, Giám đốc dự án rèn luyện phong cách lãnh đạo đồng hành thông qua mười chuẩn mực hành vi đối lập với lối quản lý mệnh lệnh hành chính đơn thuần (Bảng~\ref{tab:lanh_dao}):

\begin{table}[htbp]
\centering
\caption{Bảng đối chiếu mười phẩm chất tâm lý của người lãnh đạo nhóm dự án}
\label{tab:lanh_dao}
\small
\setstretch{1.15}
\begin{tabularx}{\textwidth}{|C{0.8cm}|Y|Y|}
\hline
\textbf{TT} & \textbf{Phẩm chất tâm lý của người lãnh đạo tài năng} & \textbf{Hành vi quản lý tiêu cực cần loại bỏ} \\
\hline
\textbf{1} & Hướng dẫn, chia sẻ chuyên môn; cung cấp rõ quy chuẩn biên soạn ngay từ đầu. & Ra lệnh cứng nhắc, áp đặt ý kiến chủ quan từ trên xuống. \\
\hline
\textbf{2} & Đặt niềm tin vào năng lực học thuật và trách nhiệm nghề nghiệp của tác giả. & Sử dụng quyền lực hành chính để gây sức ép nặng nề. \\
\hline
\textbf{3} & Động viên, khích lệ khi xử lý các bộ dữ liệu tuyển sinh và trường nghề phức tạp. & Tạo tâm lý căng thẳng, đe dọa cắt giảm kinh phí hoặc kỷ luật. \\
\hline
\textbf{4} & Dùng đại từ tập thể ``chúng ta'', ``các bạn'' để đề cao sự đóng góp của nhóm. & Thường xuyên nhấn mạnh chữ ``tôi'', đề cao cái tôi cá nhân. \\
\hline
\textbf{5} & Chủ động tháo gỡ vướng mắc thủ tục và nhận trách nhiệm khi trễ tiến độ. & Đổ lỗi cho tác giả hoặc đối tác khi xảy ra sự cố tiến độ. \\
\hline
\textbf{6} & Lắng nghe cầu thị ý kiến phản biện độc lập nâng cao chất lượng bản thảo. & Độc đoán, tự quyết định chuyên môn không qua tham vấn. \\
\hline
\textbf{7} & Khơi dậy niềm hứng khởi sáng tạo học liệu phục vụ học sinh cuối cấp. & Biến việc viết tài liệu thành nghĩa vụ hành chính khô khan. \\
\hline
\textbf{8} & Khuyến khích giải pháp sơ đồ hóa thông tin và bảng hỏi trực quan sinh động. & Bảo thủ, định kiến, chỉ tin dùng một nhóm hẹp quen thuộc. \\
\hline
\textbf{9} & Xây dựng quan hệ hợp tác tôn trọng, minh bạch với nhà thầu in ấn. & Đặt ra các yêu cầu phi thực tế rồi dồn ép đối tác thực hiện. \\
\hline
\textbf{10} & Trực tiếp sát cánh cùng các tiểu ban tại cơ sở theo tinh thần cùng làm việc. & Đứng ngoài chỉ đạo hình thức và khoán trắng cho cấp dưới. \\
\hline
\end{tabularx}
\end{table}

\subsubsection{7.3.3. Hệ thống bảy thói quen quản lý hiệu quả trong điều hành dự án}

Cùng với mười phẩm chất tâm lý trên, bảy thói quen quản lý hiệu quả được cụ thể hóa vào quy trình điều hành 90 ngày của dự án:
\begin{enumerate}
    \item \textit{Thói quen 1 (Chủ động):} Giữ thái độ chủ động trong việc cập nhật sớm chỉ tiêu tuyển sinh lớp 10 và khảo sát giá vật tư giấy in;
    \item \textit{Thói quen 2 (Bắt đầu từ mục tiêu):} Luôn bắt đầu từ hình ảnh cuốn cẩm nang hoàn chỉnh đến tay học sinh trước kỳ đăng ký nguyện vọng;
    \item \textit{Thói quen 3 (Ưu tiên điều quan trọng nhất):} Ưu tiên nguồn lực cao nhất cho chín công việc nằm trên đường găng 90 ngày;
    \item \textit{Thói quen 4 (Tư duy cùng thắng):} Duy trì tinh thần cùng thắng trong thương thảo hợp đồng in ấn và thỏa thuận thù lao tác giả;
    \item \textit{Thói quen 5 (Lắng nghe để thấu hiểu):} Lắng nghe thấu đáo phản hồi của học sinh, giáo viên trong buổi đọc thử nghiệm trước khi chốt bản thảo;
    \item \textit{Thói quen 6 (Hợp lực cộng hưởng):} Kết hợp sức mạnh cộng hưởng giữa chuyên gia tâm lý, nhà giáo dục và họa sĩ thiết kế đồ họa;
    \item \textit{Thói quen 7 (Rèn giũa bản thân):} Bố trí nhịp độ làm việc hợp lý giữa các chặng biên soạn và thẩm định để tái tạo sức sáng tạo cho toàn đội ngũ.
\end{enumerate}


\newpage

% 11. Chương 8: Kế hoạch phân tích và quản lý rủi ro dự án
\section{CHƯƠNG 8. KẾ HOẠCH PHÂN TÍCH VÀ QUẢN LÝ RỦI RO DỰ ÁN}

Chương này vận dụng quy trình quản lý rủi ro sáu bước nhằm nhận diện các yếu tố bất định qua ma trận bốn chiều, phân tích định tính danh mục rủi ro trọng yếu theo tháng, phân tích định lượng phương án hợp đồng in ấn bằng mô hình cây quyết định và giá trị tiền tệ kỳ vọng, từ đó xác lập các chiến lược ứng phó và kế hoạch dự phòng.

\subsection{8.1. Khung quản lý rủi ro dự án}

Trong dự án phát triển học liệu hướng nghiệp, rủi ro bao gồm mọi biến cố không chắc chắn có thể làm chậm tiến độ đường găng 90 ngày, làm biến động dự toán 185.000.000 đồng hoặc ảnh hưởng đến chất lượng của 2.000 cuốn tài liệu. Quy trình quản lý rủi ro được thiết lập khép kín qua sáu bước (Hình~\ref{fig:quy_trinh_rui_ro}):

\begin{figure}[htbp]
\centering
\begin{tikzpicture}[
    box/.style={rectangle, draw=red!80, fill=red!5, thick, rounded corners=2pt, align=center, inner sep=4pt, font=\scriptsize},
    arr/.style={->, thick, >=stealth, draw=red!80}
]
    \node (s1) [box] {\textbf{8.1. Lập kế hoạch}\\\textit{Quản lý rủi ro}};
    \node (s2) [box, right=0.35cm of s1] {\textbf{8.2. Nhận diện}\\\textit{Mô hình SWOT}};
    \node (s3) [box, right=0.35cm of s2] {\textbf{8.3. Định tính}\\\textit{Ưu tiên theo tháng}};
    \node (s4) [box, right=0.35cm of s3] {\textbf{8.4. Định lượng}\\\textit{Cây quyết định}};
    \node (s5) [box, right=0.35cm of s4] {\textbf{8.5. Ứng phó}\\\textit{4 Chiến lược}};
    \node (s6) [box, right=0.35cm of s5] {\textbf{8.6. Kiểm soát}\\\textit{Giám sát định kỳ}};

    \draw [arr] (s1) -- (s2);
    \draw [arr] (s2) -- (s3);
    \draw [arr] (s3) -- (s4);
    \draw [arr] (s4) -- (s5);
    \draw [arr] (s5) -- (s6);
\end{tikzpicture}
\caption{Quy trình sáu bước quản lý và kiểm soát rủi ro dự án}
\label{fig:quy_trinh_rui_ro}
\end{figure}

\subsection{8.2. Ma trận nhận diện rủi ro theo mô hình điểm mạnh, điểm yếu, thời cơ và nguy cơ}

Nhóm lập kế hoạch sử dụng mô hình \textit{SWOT} (\textit{Strengths} -- Điểm mạnh, \textit{Weaknesses} -- Điểm yếu, \textit{Opportunities} -- Thời cơ, \textit{Threats} -- Nguy cơ) để rà soát các yếu tố nội tại của Ban Quản lý dự án và các tác động từ môi trường bên ngoài, làm cơ sở nhận diện các rủi ro tiềm ẩn:

\begin{table}[htbp]
\centering
\caption{Ma trận nhận diện rủi ro dự án theo mô hình SWOT}
\label{tab:swot}
\small
\setstretch{1.15}
\begin{tabularx}{\textwidth}{|p{2.8cm}|Y|Y|}
\hline
\textbf{Chiều phân tích} & \textbf{Đặc điểm thực tế của dự án} & \textbf{Rủi ro tiềm ẩn được nhận diện} \\
\hline
\textbf{Điểm mạnh (S)} & Đội ngũ Chủ biên và tác giả có chuyên môn sâu về tâm lý học và giáo dục học; Ban QLDA trực thuộc Phòng GD\&ĐT thuận lợi điều phối 5 trường. & Nguy cơ tác giả trình bày thiên về lý luận hàn lâm, thiếu tính trực quan đối với học sinh lớp 9. \\
\hline
\textbf{Điểm yếu (W)} & Nhân sự làm việc kiêm nhiệm; ít kinh nghiệm giám sát in công nghiệp; 9/10 công việc nằm trên đường găng không có thời gian dự trữ. & Rủi ro chậm hạn nộp bản thảo lần một (ngày 40); rủi ro lỗi màu sắc hoặc bung keo gáy nếu thiếu khâu duyệt bản bông. \\
\hline
\textbf{Thời cơ (O)} & Quyết định số 522/QĐ-TTg và chính sách miễn 100\% học phí học nghề được ưu tiên; các trường THCS đang cần tài liệu chuẩn hóa. & Thuận lợi trong khâu khảo sát đầu vào, đọc thử bản thảo và tiếp nhận ấn phẩm tại các nhà trường. \\
\hline
\textbf{Nguy cơ (T)} & Tâm lý chuộng bằng cấp đại học của phụ huynh; giá giấy in có xu hướng biến động; chỉ tiêu tuyển sinh lớp 10 có thể điều chỉnh. & Phụ huynh ít quan tâm lộ trình học nghề hệ 9+; rủi ro trượt giá vật tư in ấn; dữ liệu tuyển sinh cần cập nhật gấp. \\
\hline
\end{tabularx}
\end{table}

\subsection{8.3. Danh mục các rủi ro trọng yếu theo chu kỳ tháng}

Từ các rủi ro đã nhận diện, dự án lập bảng theo dõi thứ hạng ưu tiên của năm rủi ro trọng yếu nhất qua ba tháng thực hiện (Tháng 1: từ ngày 1 đến ngày 30; Tháng 2: từ ngày 31 đến ngày 60; Tháng 3: từ ngày 61 đến ngày 90):

\begin{table}[htbp]
\centering
\caption{Bảng theo dõi danh mục năm rủi ro trọng yếu theo chu kỳ tháng}
\label{tab:rui_ro_thang}
\small
\setstretch{1.15}
\begin{tabularx}{\textwidth}{|Y|C{1.1cm}|C{1.1cm}|C{1.1cm}|p{2.2cm}|Y|}
\hline
\textbf{Danh mục rủi ro trọng yếu} & \textbf{Hạng T1} & \textbf{Hạng T2} & \textbf{Hạng T3} & \textbf{Xu hướng} & \textbf{Biện pháp kiểm soát theo giai đoạn} \\
\hline
\textbf{1. Tác giả chậm nộp bản thảo (CV C)} & \textbf{1} & \textbf{2} & Đã qua & Giảm dần sau ngày 40 & Khoán theo 20 bài chuẩn; giao ban tiến độ sáng thứ Hai; bố trí trợ lý hỗ trợ dữ liệu. \\
\hline
\textbf{2. Bản thảo phải sửa nhiều sau thẩm định (CV E, F)} & \textbf{3} & \textbf{1} & Đã qua & Đạt đỉnh ở Tháng 2 & Duyệt đề cương kỹ tại ngày 14; đọc chéo nội bộ và đọc thử 20 học sinh trước họp Hội đồng. \\
\hline
\textbf{3. Nhà in trễ hạn hoặc lỗi kỹ thuật (CV I)} & \textbf{4} & \textbf{3} & \textbf{1} & Tăng cao ở Tháng 3 & Duyệt bản bông tại ngày 76; phạt trễ hạn 1\%/ngày; kiểm định mẫu 384 cuốn tại kho. \\
\hline
\textbf{4. Biến động tăng giá giấy in (CV D, I)} & \textbf{2} & \textbf{4} & Đã qua & Giảm sau khi ký HĐ & Chào hàng cạnh tranh sớm và ký hợp đồng trọn gói cố định đơn giá từ ngày thứ 32. \\
\hline
\textbf{5. Phụ huynh e ngại thông tin phân luồng nghề} & \textbf{5} & \textbf{5} & \textbf{2} & Tăng khi bàn giao sách & Thiết kế bộ phiếu đối thoại gia đình; làm rõ quyền lợi miễn 100\% học phí và liên thông. \\
\hline
\end{tabularx}
\end{table}

\subsection{8.4. Mô hình cây quyết định và giá trị tiền tệ kỳ vọng}

Để định lượng rủi ro khi lựa chọn phương thức thực hiện gói thầu in ấn 2.000 cuốn tài liệu (mục \texttt{4.4}), Ban Quản lý dự án sử dụng \textit{Decision Tree} (Mô hình cây quyết định) kết hợp chỉ số \textit{Expected Monetary Value -- EMV} (Giá trị tiền tệ kỳ vọng).

\subsubsection{8.4.1. Bài toán lựa chọn phương án hợp đồng in ấn công nghiệp}

Tại điểm quyết định lựa chọn nhà thầu in ấn, dự án cân nhắc hai phương án có cấu trúc chi phí và phân bố xác suất rủi ro khác nhau:

\textbf{Phương án 1 (Ký hợp đồng in ấn trọn gói với nhà in công nghiệp uy tín A):} Chi phí hợp đồng dự toán là \textbf{95.000.000 đồng}. Nếu quá trình bình bản kẽm và chạy máy offset diễn ra thuận lợi (xác suất $P_1 = 0,80$), sản phẩm đạt chuẩn ngay lần đầu và đúng tiến độ, mang lại giá trị quy đổi cho dự án là 120.000.000 đồng; lợi ích ròng đạt $120.000.000 - 95.000.000 = +25.000.000$ đồng. Nếu xảy ra lỗi lệch màu nhẹ ở một số tay sách phải căn chỉnh chạy lại máy làm chậm tiến độ 2 ngày (xác suất $P_2 = 0,20$), giá trị thu về giảm xuống còn 75.000.000 đồng do phát sinh chi phí giám sát hiện trường và áp lực tiến độ; lợi ích ròng trong trường hợp này là $75.000.000 - 95.000.000 = -20.000.000$ đồng.

\textbf{Phương án 2 (Thuê cơ sở in gia công giá rẻ B sử dụng thiết bị cũ):} Chi phí hợp đồng là \textbf{80.000.000 đồng} (thấp hơn 15.000.000 đồng so với Phương án 1). Nếu quá trình in thuận lợi (xác suất $P_1 = 0,40$), giá trị mang lại ước tính đạt 110.000.000 đồng; lợi ích ròng đạt $110.000.000 - 80.000.000 = +30.000.000$ đồng. Nếu gặp sự cố lệch tram khi xuất kẽm làm kéo dài thời gian in thêm 5 ngày (xác suất $P_2 = 0,40$), giá trị thu về giảm xuống còn 55.000.000 đồng; lợi ích ròng là $55.000.000 - 80.000.000 = -25.000.000$ đồng. Nếu xảy ra lỗi nghiêm trọng (lem mực, bung keo gáy diện rộng không vượt qua kỳ kiểm định mẫu 384 cuốn với xác suất $P_3 = 0,20$), giá trị thu về chỉ còn 20.000.000 đồng; lợi ích ròng âm ở mức $20.000.000 - 80.000.000 = -60.000.000$ đồng.

\subsubsection{8.4.2. Sơ đồ cây quyết định và tính toán giá trị tiền tệ kỳ vọng}

Trên sơ đồ cây quyết định Hình~\ref{fig:decision_tree}, điểm quyết định được biểu diễn bằng ô hình vuông và các điểm cơ hội gắn với xác suất biến cố được biểu diễn bằng nút hình tròn:

\begin{figure}[htbp]
\centering
\begin{tikzpicture}[
    decision/.style={rectangle, draw=blue!90, fill=blue!10, thick, minimum size=8mm, align=center, font=\small\bfseries},
    chance/.style={circle, draw=orange!90, fill=orange!15, thick, minimum size=7mm, font=\small\bfseries},
    outcome/.style={rectangle, draw=gray!70, fill=gray!5, rounded corners=2pt, font=\scriptsize, align=left, inner sep=3pt},
    edge from parent/.style={draw=black!70, thick, -{Stealth[length=2mm]}}
]
    % Điểm quyết định
    \node (dec) at (0,0) [decision] {Quyết định\\chọn gói\\in ấn};

    % Nút cơ hội 1
    \node (c1) at (3.5, 2.0) [chance] {$O_1$};
    % Nút cơ hội 2
    \node (c2) at (3.5, -2.0) [chance] {$O_2$};

    \draw [thick, -{Stealth[length=2mm]}] (dec) -- node[above=2pt, sloped, font=\scriptsize] {PA 1: Nhà in uy tín A (95 tr)} (c1);
    \draw [thick, -{Stealth[length=2mm]}] (dec) -- node[below=2pt, sloped, font=\scriptsize] {PA 2: Xưởng in giá rẻ B (80 tr)} (c2);

    % Nhánh của C1
    \node (o11) at (8.5, 2.7) [outcome] {Thuận lợi ($P=0,80$): \textbf{+25.000.000 đ}};
    \node (o12) at (8.5, 1.3) [outcome] {Lỗi nhẹ ($P=0,20$): \textbf{-20.000.000 đ}};
    \draw [thick, -{Stealth[length=2mm]}] (c1) -- (o11);
    \draw [thick, -{Stealth[length=2mm]}] (c1) -- (o12);

    % Nhánh của C2
    \node (o21) at (8.5, -1.0) [outcome] {Thuận lợi ($P=0,40$): \textbf{+30.000.000 đ}};
    \node (o22) at (8.5, -2.0) [outcome] {Lỗi kẽm trễ 5 ngày ($P=0,40$): \textbf{-25.000.000 đ}};
    \node (o23) at (8.5, -3.0) [outcome] {Hỏng lô hàng ($P=0,20$): \textbf{-60.000.000 đ}};
    \draw [thick, -{Stealth[length=2mm]}] (c2) -- (o21);
    \draw [thick, -{Stealth[length=2mm]}] (c2) -- (o22);
    \draw [thick, -{Stealth[length=2mm]}] (c2) -- (o23);
\end{tikzpicture}
\caption{Sơ đồ cây quyết định lựa chọn phương án hợp đồng in ấn công nghiệp}
\label{fig:decision_tree}
\end{figure}

Tính toán giá trị tiền tệ kỳ vọng từ phải sang trái tại hai nút cơ hội:

\begin{equation}
\begin{split}
\text{EMV}_1 &= (0,80 \times 25.000.000) + (0,20 \times (-20.000.000)) \\
&= 20.000.000 - 4.000.000 = \mathbf{+16.000.000\text{ đồng}}
\end{split}
\end{equation}

\begin{equation}
\begin{split}
\text{EMV}_2 &= (0,40 \times 30.000.000) + (0,40 \times (-25.000.000)) + (0,20 \times (-60.000.000)) \\
&= 12.000.000 - 10.000.000 - 12.000.000 = \mathbf{-10.000.000\text{ đồng}}
\end{split}
\end{equation}

Do $\text{EMV}(\text{Phương án 1}) = +16.000.000\text{ đồng} > \text{EMV}(\text{Phương án 2}) = -10.000.000\text{ đồng}$, Ban Quản lý dự án quyết định lựa chọn Phương án 1 (mời thầu các doanh nghiệp in công nghiệp uy tín có máy in offset bốn màu hiện đại với mức trần dự toán gói thầu 95.000.000 đồng) và loại bỏ hoàn toàn phương án thuê cơ sở in gia công giá rẻ.

\subsection{8.5. Chiến lược đối phó rủi ro và kế hoạch dự phòng}

\subsubsection{8.5.1. Bốn chiến lược đối phó rủi ro}

Dự án thiết lập hệ thống giải pháp ứng phó theo bốn nhóm chiến lược gồm \textit{Risk Avoidance} (Tránh rủi ro), \textit{Risk Mitigation} (Giảm nhẹ rủi ro), \textit{Risk Transference} (Thuyên chuyển rủi ro) và \textit{Risk Acceptance} (Chấp nhận rủi ro) tại Bảng~\ref{tab:chien_luoc_rui_ro}:

\begin{table}[htbp]
\centering
\caption{Phân nhóm chiến lược đối phó rủi ro của dự án}
\label{tab:chien_luoc_rui_ro}
\small
\setstretch{1.15}
\begin{tabularx}{\textwidth}{|p{2.6cm}|p{3.2cm}|Y|Y|}
\hline
\textbf{Chiến lược} & \textbf{Nguyên tắc tác động} & \textbf{Rủi ro áp dụng} & \textbf{Giải pháp thực thi cụ thể} \\
\hline
\textbf{Tránh rủi ro} & Thay đổi phương thức thực hiện loại bỏ nguyên nhân. & Hỏng hàng và trễ tiến độ do thuê cơ sở in gia công lạc hậu. & Đưa tiêu chí bắt buộc máy offset 4 màu và mẫu giấy đạt chuẩn vào vòng kỹ thuật để loại cơ sở yếu kém. \\
\hline
\textbf{Giảm nhẹ rủi ro} & Giảm xác suất xảy ra hoặc thu hẹp mức độ thiệt hại. & Chậm nộp bản thảo (C) và bản thảo không qua thẩm định (E). & Áp dụng kỹ thuật gối đầu gửi phản biện từng phần và đọc thử với học sinh để chỉnh sửa sớm. \\
\hline
\textbf{Thuyên chuyển rủi ro} & Chuyển giao gánh nặng tài chính sang bên thứ ba qua hợp đồng. & Rủi ro trượt giá giấy và sai hỏng kỹ thuật khi in hàng loạt. & Ký hợp đồng trọn gói cố định 47.000 đ/cuốn; ràng buộc nhà thầu chịu 100\% phí in bù nếu mẫu vượt ngưỡng lỗi. \\
\hline
\textbf{Chấp nhận rủi ro} & Ghi nhận rủi ro khách quan và chuẩn bị phương án dự phòng. & Địa phương điều chỉnh chỉ tiêu tuyển sinh sau khi sách đã in. & Chuẩn bị sẵn mẫu tờ rời phụ lục cập nhật chỉ tiêu kẹp bổ sung vào cuốn sách khi bàn giao về trường. \\
\hline
\end{tabularx}
\end{table}

\subsubsection{8.5.2. Kế hoạch dự phòng và kế hoạch rút lui}

Đối với nhân sự biên soạn, dự án thiết lập \textit{Contingency Plan} (Kế hoạch dự phòng): nếu một tác giả chính gặp sự cố sức khỏe đột xuất trong giai đoạn viết bản thảo (từ ngày 15 đến ngày 40), Chủ biên chuyên môn cùng thành viên dự bị trong nhóm tác giả sẽ trực tiếp tiếp quản phần bài viết đó để giữ vững mốc bàn giao bản thảo lần một vào ngày thứ 40.

Đối với nhà thầu in ấn, dự án xác lập \textit{Fallback Plan} (Kế hoạch rút lui): trong trường hợp dây chuyền in của nhà thầu trúng thầu gặp sự cố kỹ thuật không thể khắc phục trước ngày thứ 77, Ban Quản lý dự án sẽ kích hoạt biên bản bảo lưu kết quả chào hàng để ký hợp đồng thay thế với nhà thầu xếp hạng hai (Công ty Cổ phần In và Truyền thông Giáo dục Thủ đô). Các phát sinh tài chính ngoài dự toán (nếu có) được xử lý thông qua Quỹ dự phòng rủi ro chi phí 9.500.000 đồng theo đúng quy trình phê duyệt đã xác lập tại Mục 5.4.2 (Chương 5).


\newpage

% 12. Chương 9: Kế hoạch quản lý truyền thông dự án
\section{CHƯƠNG 9. KẾ HOẠCH QUẢN LÝ TRUYỀN THÔNG DỰ ÁN}

Chương này xác lập kế hoạch truyền thông với các bên liên quan, quy định phương thức phân phối thông tin hành chính và tác nghiệp, thiết kế hệ thống báo cáo kết quả công việc (gồm báo cáo hiện trạng, báo cáo tiến độ và báo cáo dự báo) và chuẩn hóa thủ tục kết thúc hành chính, lưu trữ hồ sơ dự án.

\subsection{9.1. Kế hoạch truyền thông với các bên liên quan}

Để bảo đảm thông tin điều hành được chuyển tải đúng đối tượng và đúng thời điểm, Bảng 9.1 quy định cụ thể nội dung, định dạng, tần suất và đầu mối chịu trách nhiệm truyền thông đối với bảy nhóm chủ thể tham gia dự án:

\begin{center}
{\small
\setstretch{1.15}
\begin{longtable}{|p{2.6cm}|p{3.2cm}|p{2.6cm}|p{2.0cm}|p{2.2cm}|Y|}
\caption{Kế hoạch truyền thông chi tiết với các bên liên quan} \label{tab:truyen_thong} \\
\hline
\textbf{Nhóm chủ thể} & \textbf{Nội dung truyền thông} & \textbf{Hình thức, kênh} & \textbf{Đầu mối} & \textbf{Tần suất, thời điểm} & \textbf{Mục đích quản trị} \\
\hline
\endfirsthead

\multicolumn{6}{c}{\textit{(Tiếp theo Bảng 9.1: Kế hoạch truyền thông chi tiết)}} \\
\hline
\textbf{Nhóm chủ thể} & \textbf{Nội dung truyền thông} & \textbf{Hình thức, kênh} & \textbf{Đầu mối} & \textbf{Tần suất, thời điểm} & \textbf{Mục đích quản trị} \\
\hline
\endhead

\hline
\multicolumn{6}{r}{\textit{Còn tiếp sang trang sau...}} \\
\endfoot

\hline
\endlastfoot

\textbf{Lãnh đạo Phòng GD\&ĐT (Chủ đầu tư)} & Tờ trình phê duyệt; báo cáo định kỳ mốc kiểm soát; hồ sơ nghiệm thu tổng kết. & Văn bản hành chính có chữ ký, dấu pháp nhân kèm email công vụ. & Giám đốc dự án & Khởi công; cuối mỗi tháng (ngày 30, 60); kết thúc (ngày 90). & Báo cáo giải trình, xin ý kiến chỉ đạo và phê chuẩn nghiệm thu, giải ngân. \\
\hline
\textbf{Ban Quản lý dự án (Điều hành)} & Kế hoạch tác nghiệp tuần; biên bản giao ban; bảng theo dõi đường găng và rủi ro. & Họp giao ban trực tiếp kết hợp trao đổi nhóm trực tuyến nội bộ. & Giám đốc DA và Thư ký & Sáng thứ Hai hàng tuần và khi có sự cố kỹ thuật. & Điều phối tiến độ hàng tuần, tháo gỡ kịp thời vướng mắc đường găng. \\
\hline
\textbf{Chủ biên và 4 nhóm tác giả} & Khung đề cương chi tiết; bản thảo các phần; nhận xét phản biện và thẩm định. & Tệp tin lưu trữ đám mây dùng chung và họp chuyên môn trực tiếp. & Chủ biên chuyên môn & Hàng tuần giai đoạn viết (ngày 15--40) và hiệu đính (53--60). & Đồng nhất cấu trúc, văn phong sư phạm và hướng dẫn chỉnh sửa bản thảo. \\
\hline
\textbf{Hội đồng thẩm định} & Giấy mời họp; bản thảo đóng cuốn lần một; mẫu phiếu đánh giá phản biện. & Hồ sơ văn bản giấy trực tiếp kèm tệp tin bảo mật qua email công vụ. & Chủ tịch Hội đồng & Gửi trước phiên họp chính thức tối thiểu 07 ngày. & Cung cấp tài liệu để chuyên gia đọc phản biện độc lập trước họp tập trung. \\
\hline
\textbf{Thiết kế chế bản và Nhà thầu in} & Hồ sơ yêu cầu báo giá; hợp đồng kinh tế; ma-két chế bản; bản bông niêm phong. & Hợp đồng giấy tờ pháp lý; dữ liệu qua máy chủ; biên bản giao nhận. & Phụ trách Mua sắm & Giai đoạn mời thầu (15--32) và chế bản, in ấn (61--84). & Ràng buộc pháp lý, kiểm soát quy cách chế bản và giám sát xuất xưởng. \\
\hline
\textbf{Hiệu trưởng 5 trường THCS} & Công văn phối hợp khảo sát; thông báo lịch giao sách và giấy mời nghiệm thu. & Công văn điện tử qua hệ thống quản lý văn bản của ngành giáo dục. & Đại diện Tiểu ban Kiểm định & Tuần 1 (khảo sát) và Tuần 13 (bàn giao 2.000 cuốn). & Phối hợp thu thập dữ liệu và chuẩn bị kho tiếp nhận, phân phối sách. \\
\hline
\textbf{Giáo viên tư vấn, học sinh, phụ huynh} & Bản đọc thử nghiệm; ấn phẩm chính thức sau khi bàn giao về nhà trường. & Tài liệu in thử trong thảo luận nhóm và sách hoàn chỉnh tại thư viện. & Tổ GV tư vấn 5 trường cơ sở & Tuần 7 (đọc thử) và sau ngày 90 (khai thác sử dụng). & Ghi nhận phản hồi thực tế và đưa tài liệu vào tư vấn phân luồng hiệu quả. \\
\hline
\end{longtable}
}
\end{center}

\subsection{9.2. Cơ chế phân phối thông tin và quản lý phiên bản tài liệu}

Hệ thống thông tin của dự án được phân luồng thành hai kênh rõ rệt. Kênh văn bản hành chính pháp lý được áp dụng bắt buộc đối với các tài liệu làm căn cứ thanh quyết toán và nghiệm thu nhà nước (Quyết định phê duyệt, Hợp đồng kinh tế, Biên bản họp Hội đồng thẩm định, Biên bản kiểm định mẫu thống kê và Biên bản bàn giao tài liệu), yêu cầu lập thành văn bản giấy có đầy đủ chữ ký tươi và dấu pháp nhân.

Kênh tác nghiệp số hóa được sử dụng cho hoạt động biên soạn và chế bản hàng ngày thông qua thư điện tử công vụ và thư mục lưu trữ đám mây phân quyền theo từng tiểu ban. Để tránh nhầm lẫn giữa các lần chỉnh sửa bản thảo của bốn nhóm tác giả, mọi tệp tin đưa lên hệ thống đều được đặt tên theo quy tắc mã hóa thống nhất gồm: \texttt{[Mã gói công việc]\_[Tên phần]\_[Ngày cập nhật]\_[Phiên bản V1.0 cho bản thảo lần một / V2.0 cho bản thảo sạch sau thẩm định]}.

\subsection{9.3. Hệ thống báo cáo kết quả công việc dự án}

\subsubsection{9.3.1. Cấu trúc ba chiều của hệ thống báo cáo dự án}

Theo lý thuyết quản lý truyền thông dự án, báo cáo kết quả công việc được cấu trúc thành ba phần bổ trợ nhau (Hình~\ref{fig:he_thong_bao_cao}): \textit{Status Report} (Báo cáo hiện trạng) phản ánh vị trí hiện tại của dự án về tỷ lệ hoàn thành cấu trúc phân rã công việc và số tiền đã giải ngân so với đường mức chi phí cơ sở; \textit{Progress Report} (Báo cáo tiến độ) tổng hợp các sản phẩm đầu ra cụ thể đã hoàn thành nghiệm thu trong kỳ và đối chiếu với các mốc sự kiện trên đường găng; và \textit{Forecast Report} (Báo cáo dự báo) đưa ra dự kiến về khả năng hoàn thành các công việc tiếp theo cùng cảnh báo sớm về tiến độ hoặc chi phí.

\begin{figure}[htbp]
\centering
\begin{tikzpicture}[
    box/.style={rectangle, draw=blue!80, fill=blue!5, thick, rounded corners=3pt, align=center, inner sep=6pt, font=\small},
    arr/.style={->, thick, >=stealth, draw=blue!80}
]
    \node (b1) [box] {\textbf{1. BÁO CÁO HIỆN TRẠNG}\\\textit{Vị trí hiện tại của dự án:}\\\textit{Khối lượng WBS và giải ngân lũy kế}};
    \node (b2) [box, right=0.8cm of b1] {\textbf{2. BÁO CÁO TIẾN ĐỘ}\\\textit{Kết quả đạt được trong kỳ:}\\\textit{Sản phẩm nghiệm thu so với đường găng}};
    \node (b3) [box, right=0.8cm of b2] {\textbf{3. BÁO CÁO DỰ BÁO}\\\textit{Xu hướng kỳ tiếp theo:}\\\textit{Dự kiến ngày hoàn thành và rủi ro chi phí}};

    \draw [arr] (b1) -- (b2);
    \draw [arr] (b2) -- (b3);
\end{tikzpicture}
\caption{Cấu trúc ba chiều bổ trợ của hệ thống báo cáo kết quả công việc dự án}
\label{fig:he_thong_bao_cao}
\end{figure}

\subsubsection{9.3.2. Báo cáo tích hợp tại mốc kiểm soát ngày thứ 60 (Ngày 30/04/2026)}

Tại thời điểm kết thúc Giai đoạn 3.0 trên cấu trúc phân rã công việc (ngày thứ 60 của dự án, tức ngày 30 tháng 4 năm 2026), Giám đốc dự án trình Lãnh đạo Phòng Giáo dục và Đào tạo Báo cáo tích hợp số 02/BC-GDHN:

\textbf{Về hiện trạng dự án tại ngày thứ 60:} Dự án đã đi qua $60/90$ ngày (tương ứng 66,7\% tổng quỹ thời gian), hoàn thành trọn vẹn 6 trên 10 công việc của sơ đồ mạng (gồm các công việc A, B, C, D, E, F, đạt 60\% số công việc mạng) và 13 trên 20 gói công việc trong cấu trúc phân rã công việc (toàn bộ các mục \texttt{1.0}, \texttt{2.0}, \texttt{3.0} và gói \texttt{4.1}, đạt 65\% số gói công việc). Về tài chính, bộ phận kế toán đã hoàn tất giải ngân 3 đợt đầu tiên với tổng số tiền lũy kế là \textbf{62.500.000 đồng} (gồm 12.500.000 đồng Đợt 1, 32.800.000 đồng Đợt 2 và 17.200.000 đồng Đợt 3), đạt \textbf{33,78\%} tổng ngân sách, khớp với đường mức chi phí cơ sở tại Chương 5.

\textbf{Về kết quả tiến độ đạt được trong kỳ (từ ngày thứ 31 đến ngày thứ 60):} Bốn nhóm tác giả đã bàn giao bản thảo lần một đúng ngày thứ 40. Hội đồng thẩm định chuyên môn đã họp hai phiên và thông qua bản thảo với tỷ lệ 5/5 phiếu tán thành vào ngày thứ 52. Sau khi tiếp thu ý kiến chuyên gia và kết quả đọc thử của 20 học sinh, Chủ biên đã hoàn thành bản thảo sạch 80 trang ruột và được Giám đốc dự án ký duyệt xuất bản đúng ngày thứ 60 (đảm bảo mốc sự kiện găng thứ 6). Ở nhánh mua sắm song song, Hợp đồng in ấn trọn gói số 08/HĐ-IN đã được ký kết từ ngày thứ 32 với giá trúng thầu 94.000.000 đồng (tiết kiệm 1.000.000 đồng so với dự toán gói thầu 95.000.000 đồng).

\textbf{Về dự báo kỳ tiếp theo (từ ngày thứ 61 đến ngày thứ 90):} Nhóm thiết kế chế bản bắt đầu thực hiện công việc G đúng ngày thứ 61 (01/05/2026) và dự kiến bàn giao tệp tin ma-két vào ngày thứ 72 để in thử bản bông (ký duyệt ngày thứ 76). Nhờ nhà thầu in đã chuẩn bị sẵn nguồn giấy Fort 100 gsm và Couche 250 gsm theo hợp đồng, lô hàng 2.000 cuốn dự báo sẽ hoàn thành xuất xưởng đúng ngày thứ 84 và bàn giao cho 5 trường vào ngày thứ 90 mà không cần sử dụng đến quỹ dự phòng 9.500.000 đồng.

\subsection{9.4. Thủ tục kết thúc hành chính và lưu trữ hồ sơ}

\subsubsection{9.4.1. Thủ tục xác nhận và bàn giao chính thức sản phẩm}

Tại ngày thứ 90 (30/05/2026), thủ tục kết thúc hành chính và chấp nhận sản phẩm được thực hiện khi hội đủ ba căn cứ: Phiếu xác nhận kết quả kiểm định mẫu thống kê 384 cuốn tại kho đạt yêu cầu kỹ thuật; Giấy phép xuất bản và xác nhận nộp lưu chiểu theo Luật Xuất bản; cùng Biên bản bàn giao ba bên (giữa đại diện Phòng Giáo dục và Đào tạo, Ban Quản lý dự án và Hiệu trưởng năm trường trung học cơ sở) xác nhận đã giao nhận đủ 2.000 cuốn ấn phẩm (40 thùng) kèm 01 bộ tệp tin chế bản điện tử gốc.

\subsubsection{9.4.2. Đóng gói lưu trữ hồ sơ và đúc rút bài học kinh nghiệm}

Theo đúng định mức tại gói công việc \texttt{5.4} (Chương 3), toàn bộ chứng từ và tài liệu phát sinh trong 90 ngày được đóng thành \textbf{03 bộ hồ sơ lưu trữ gốc} (01 bộ lưu tại Văn phòng Phòng Giáo dục và Đào tạo, 01 bộ lưu tại bộ phận Kế hoạch -- Tài chính phục vụ kiểm toán nhà nước, và 01 bộ lưu tại trường trung học cơ sở đại diện cụm). Mỗi bộ hồ sơ lưu trữ bao gồm năm tập tài liệu thành phần:
\begin{enumerate}
    \item \textit{Tập 1 -- Hồ sơ pháp lý và kế hoạch cơ sở:} Quyết định số 128/QĐ-PGDĐT, Điều lệ dự án, Tuyên bố phạm vi, Từ điển phân rã công việc, Sơ đồ mạng đường găng và Dự toán ngân sách.
    \item \textit{Tập 2 -- Hồ sơ khảo sát và học thuật:} 350 phiếu khảo sát gốc, Báo cáo dữ liệu đầu vào, Khung đề cương, Bản thảo lần một, 03 phiếu nhận xét phản biện độc lập, Biên bản họp Hội đồng thẩm định và Bản thảo sạch.
    \item \textit{Tập 3 -- Hồ sơ chế bản và mẫu đối chứng:} Thiết bị lưu trữ chứa tệp tin chế bản gốc và cuốn bản bông mẫu có chữ ký niêm phong duyệt màu.
    \item \textit{Tập 4 -- Hồ sơ mua sắm và thanh quyết toán:} Hồ sơ yêu cầu báo giá, Báo cáo đánh giá của Tổ chuyên gia, Hợp đồng in ấn số 08/HĐ-IN, Hóa đơn giá trị gia tăng, Biên bản thanh lý hợp đồng và chứng từ giải ngân 5 đợt.
    \item \textit{Tập 5 -- Hồ sơ kiểm định và nghiệm thu bàn giao:} Phiếu đo đạc mẫu thống kê 384 cuốn, Biên bản nghiệm thu cấp cơ sở, Biên bản bàn giao cho 5 trường và Báo cáo tổng kết đóng dự án.
\end{enumerate}

Báo cáo tổng kết đóng dự án đồng thời ghi nhận ba bài học kinh nghiệm cốt lõi cho cơ quan quản lý giáo dục địa phương: việc tách nhánh mua sắm in ấn (công việc D) chạy song song với nhánh biên soạn (công việc C) trên sơ đồ mạng giúp tiết kiệm 18 ngày trên trục tiến độ; việc kết hợp thẩm định của nhà khoa học với buổi đọc thử của chính học sinh lớp 9 giúp dung hòa tính chuẩn mực pháp lý với ngôn ngữ gần gũi lứa tuổi; và việc ràng buộc thanh toán qua chốt ký duyệt bản bông cùng kiểm định mẫu thống kê $n = 384$ cuốn giúp bảo vệ an toàn chất lượng cho nguồn vốn ngân sách nhà nước.


\newpage

% 13. Chương 10: Kế hoạch quản lý mua sắm dự án
\section{CHƯƠNG 10. KẾ HOẠCH QUẢN LÝ MUA SẮM DỰ ÁN}

Chương này vận dụng quy trình quản lý mua sắm sáu bước và các quy định hiện hành của Luật Đấu thầu số 22/2023/QH15 (được sửa đổi, bổ sung bởi Luật số 57/2024/QH15 và Luật số 90/2025/QH15) cùng Nghị định số 24/2024/NĐ-CP để lập kế hoạch mua sắm, xây dựng hồ sơ yêu cầu báo giá, đánh giá lựa chọn nhà thầu in ấn, quản lý hợp đồng trọn gói và thực hiện nghiệm thu, thanh lý hợp đồng.

\subsection{10.1. Kế hoạch mua sắm và quyết định thuê ngoài}

Kế hoạch mua sắm của dự án được xác lập thông qua việc giải quyết năm câu hỏi định hướng chiến lược, khởi đầu từ kỹ thuật \textit{Make or Buy Analysis} (Phân tích quyết định tự thực hiện hay thuê ngoài) và lập \textit{Request for Quotation -- RFQ} (Hồ sơ yêu cầu báo giá) tại Hình~\ref{fig:nam_cau_hoi_mua_sam}:

\begin{figure}[htbp]
\centering
\begin{tikzpicture}[
    box/.style={rectangle, draw=blue!80, fill=blue!5, thick, rounded corners=3pt, align=left, inner sep=5pt, font=\scriptsize, text width=12.5cm},
    arr/.style={->, thick, >=stealth, draw=blue!80}
]
    \node (q1) [box] {\textbf{1. Có mua ngoài hay tự thực hiện?} Tự thực hiện nghiên cứu, biên soạn và thiết kế dàn trang; thuê ngoài trọn gói bình bản kẽm và in ấn offset công nghiệp.};
    \node (q2) [box, below=0.35cm of q1] {\textbf{2. Mua sản phẩm hoặc dịch vụ gì?} Gói dịch vụ phi tư vấn: Xuất kẽm công nghiệp, in offset 4 màu, cán màng bìa và đóng keo nhiệt 2.000 cuốn tài liệu.};
    \node (q3) [box, below=0.35cm of q2] {\textbf{3. Mua khối lượng bao nhiêu?} Đúng 2.000 cuốn ấn phẩm hoàn chỉnh đóng trong 40 thùng tiêu chuẩn kèm bộ dữ liệu bình bản xuất kẽm gốc.};
    \node (q4) [box, below=0.35cm of q3] {\textbf{4. Mua bằng phương thức nào?} Chào hàng cạnh tranh theo quy trình rút gọn (Luật Đấu thầu số 22/2023/QH15 và Nghị định số 24/2024/NĐ-CP).};
    \node (q5) [box, below=0.35cm of q4] {\textbf{5. Khi nào tiến hành mua sắm?} Phát hành hồ sơ ngày 15, mở báo giá ngày 28, ký hợp đồng ngày 32 (thực hiện trong 18 ngày của công việc D).};

    \draw [arr] (q1) -- (q2);
    \draw [arr] (q2) -- (q3);
    \draw [arr] (q3) -- (q4);
    \draw [arr] (q4) -- (q5);
\end{tikzpicture}
\caption{Mô hình năm câu hỏi định hướng lập kế hoạch mua sắm dự án}
\label{fig:nam_cau_hoi_mua_sam}
\end{figure}

Phân tích tự thực hiện hay thuê ngoài chỉ rõ: Phòng Giáo dục và Đào tạo cùng đội ngũ chuyên gia sư phạm, tâm lý học và nhóm họa sĩ đồ họa chuyên môn có đủ năng lực thực hiện các gói công việc nghiên cứu khảo sát (\texttt{1.0}), biên soạn (\texttt{2.0}), thẩm định (\texttt{3.0}) cũng như thiết kế bìa, vẽ 25 hình minh họa, dàn trang điện tử 80 trang ruột và in thử kỹ thuật số (\texttt{4.2a} đến \texttt{4.3} với kinh phí 10.500.000 đồng). Ngược lại, cơ quan giáo dục không có chức năng in xuất bản phẩm công nghiệp và không sở hữu dây chuyền máy ghi kẽm, máy in offset bốn màu khổ lớn hay máy vào keo nhiệt tự động. Do đó, dự án quyết định thuê ngoài trọn gói khâu bình bản kẽm, in ấn offset hàng loạt 2.000 cuốn và gia công thành phẩm (gói công việc \texttt{4.4}).

Về hình thức lựa chọn nhà thầu, gói thầu in ấn là dịch vụ phi tư vấn thông dụng, có tiêu chuẩn kỹ thuật rõ ràng và có giá dự toán gói thầu được duyệt là \textbf{95.000.000 đồng} (nằm trong hạn mức áp dụng chào hàng cạnh tranh). Căn cứ Điều 24 Luật Đấu thầu số 22/2023/QH15 và Điều 78, Điều 79 Nghị định số 24/2024/NĐ-CP, dự án áp dụng hình thức \textbf{Chào hàng cạnh tranh theo quy trình rút gọn} để tiết kiệm thời gian và chi phí tổ chức đấu thầu. Toàn bộ thủ tục lựa chọn nhà thầu được thực hiện trong 18 ngày của công việc không găng D (từ ngày thứ 15 đến ngày thứ 32).

\subsection{10.2. Hồ sơ yêu cầu báo giá và tiếp cận nhà thầu}

\subsubsection{10.2.1. Các thông số pháp lý và kỹ thuật trong hồ sơ yêu cầu báo giá}

Hồ sơ yêu cầu báo giá do Ban Quản lý dự án phát hành vào ngày thứ 15 (16/03/2026) xác định rõ tên gói thầu là \textit{``Dịch vụ bình bản kẽm công nghiệp, in ấn offset và gia công thành phẩm 2.000 cuốn Bộ tài liệu tư vấn phân luồng và định hướng nghề nghiệp cho học sinh cuối cấp trung học cơ sở''}, với mức trần giá gói thầu là \textbf{95.000.000 đồng} (đã bao gồm thuế giá trị gia tăng, vật tư giấy, kẽm, mực in, gia công và chi phí giao hàng tại kho).

Hồ sơ yêu cầu nhà thầu cam kết đáp ứng đầy đủ các tiêu chuẩn kỹ thuật xuất bản đã quy định tại Mục 3.1 (Chương 3) và Mục 6.2.2 (Chương 6), đồng thời gửi kèm mẫu giấy ruột Fort trắng 100 gsm và mẫu giấy bìa Couche 250 gsm trong hồ sơ chào giá để Tổ chuyên gia kiểm tra thực tế.

Thời gian thực hiện hợp đồng tại xưởng in được ấn định là \textbf{12 ngày} (từ ngày thứ 73 đến ngày thứ 84 của dự án), tính từ thời điểm nhà thầu nhận bộ tệp tin ma-két chế bản điện tử vào cuối ngày thứ 72: bao gồm \textbf{4 ngày} bình bản và phối hợp in thử, rà soát ký duyệt mẫu bản bông (công việc H, từ ngày 73 đến ngày 76) và \textbf{8 ngày} chạy máy in offset hàng loạt 2.000 cuốn, cán màng mờ bìa, vào keo nhiệt và đóng thùng xuất xưởng (công việc I, từ ngày 77 đến ngày 84).

\subsubsection{10.2.2. Danh sách các nhà thầu được gửi yêu cầu báo giá}

Ban Quản lý dự án gửi trực tiếp Hồ sơ yêu cầu báo giá tới ba doanh nghiệp có giấy phép hoạt động in xuất bản phẩm độc lập trên địa bàn:
\begin{enumerate}
    \item Nhà thầu 1: Công ty Trách nhiệm hữu hạn In ấn và Thương mại Quân đội.
    \item Nhà thầu 2: Công ty Cổ phần In và Truyền thông Giáo dục Thủ đô.
    \item Nhà thầu 3: Công ty Trách nhiệm hữu hạn Sản xuất Bao bì và Xuất bản phẩm Bách Khoa.
\end{enumerate}

\subsection{10.3. Kết quả đánh giá báo giá và lựa chọn nhà thầu}

Đúng ngày thứ 28 (29/03/2026), Tổ chuyên gia mở và đánh giá ba hồ sơ báo giá theo quy trình hai bước: kiểm tra tư cách pháp lý, năng lực kỹ thuật theo tiêu chí \textit{``Đạt / Không đạt''}, sau đó xếp hạng giá dự thầu thấp nhất trong số các nhà thầu vượt qua vòng kỹ thuật (Bảng~\ref{tab:danh_gia_thau}):

\begin{table}[htbp]
\centering
\caption{Bảng đánh giá hồ sơ chào hàng cạnh tranh gói thầu in ấn}
\label{tab:danh_gia_thau}
\small
\setstretch{1.15}
\begin{tabularx}{\textwidth}{|Y|p{2.6cm}|p{2.6cm}|p{2.6cm}|p{2.6cm}|}
\hline
\textbf{Tiêu chí đánh giá} & \textbf{Mức yêu cầu} & \textbf{Công ty In Quân đội} & \textbf{Công ty In Giáo dục} & \textbf{Công ty In Bách Khoa} \\
\hline
\multicolumn{5}{|l|}{\textbf{Bước 1: Đánh giá năng lực kỹ thuật}} \\
\hline
1. Tư cách pháp nhân, giấy phép & Đạt / Không đạt & Đạt & Đạt & Đạt \\
\hline
2. Dây chuyền offset 4 màu & Đạt / Không đạt & Đạt (Máy Heidelberg) & Đạt (Máy Komori) & Đạt (Máy Ryobi) \\
\hline
3. Mẫu giấy gửi kèm & Fort 100, Couche 250 & Đạt chuẩn định lượng & Đạt chuẩn định lượng & Không đạt (Fort 80) \\
\hline
4. Cam kết tiến độ 12 ngày & Đạt / Không đạt & Đạt (Đúng 12 ngày) & Đạt (Đúng 12 ngày) & Đạt (Đúng 12 ngày) \\
\hline
\textit{Kết luận vòng kỹ thuật} & & \textbf{Đạt vào xét giá} & \textbf{Đạt vào xét giá} & \textbf{Bị loại kỹ thuật} \\
\hline
\multicolumn{5}{|l|}{\textbf{Bước 2: Đánh giá tài chính và xếp hạng}} \\
\hline
Giá chào thầu trọn gói & Trần: 95.000.000 đ & \textbf{94.000.000 đ} & \textbf{96.800.000 đ} & Không xét giá \\
\hline
Đơn giá / cuốn ấn phẩm & Trần: 47.500 đ/cuốn & \textbf{47.000 đ/cuốn} & \textbf{48.400 đ/cuốn} & Không áp dụng \\
\hline
Chênh lệch so với dự toán & & \textbf{Tiết kiệm 1.000.000 đ} & Vượt 1.800.000 đ & Không áp dụng \\
\hline
\textbf{Kết luận xếp hạng} & & \textbf{Hạng 1 (Trúng thầu)} & \textbf{Hạng 2 (Dự bị)} & \textbf{Không đạt} \\
\hline
\end{tabularx}
\end{table}

Căn cứ kết quả tại Bảng~\ref{tab:danh_gia_thau}, vào ngày thứ 32 (02/04/2026), Giám đốc dự án phê duyệt kết quả lựa chọn nhà thầu và ký Hợp đồng kinh tế số 08/HĐ-IN với \textbf{Công ty Trách nhiệm hữu hạn In ấn và Thương mại Quân đội} với giá hợp đồng trọn gói là \textbf{94.000.000 đồng} (đơn giá 47.000 đồng/cuốn, tiết kiệm 1.000.000 đồng cho ngân sách dự án).

\subsection{10.4. Khung hợp đồng kinh tế và cơ chế kiểm soát thay đổi}

\subsubsection{10.4.1. Hình thức hợp đồng và các điều khoản ràng buộc}

Dự án áp dụng hình thức \textit{Lump-sum Fixed-price Contract} (Hợp đồng trọn gói) quy định tại Điều 64 Luật Đấu thầu số 22/2023/QH15, trong đó mức giá 94.000.000 đồng được cố định trong suốt thời gian thực hiện hợp đồng.

Nội dung Hợp đồng số 08/HĐ-IN ràng buộc chặt chẽ bốn nghĩa vụ pháp lý giữa hai bên:
\begin{enumerate}
    \item Nhà thầu có nghĩa vụ bảo mật bản thảo và bản quyền, không được tự ý sao chụp hoặc in thêm ngoài số lượng 2.000 cuốn;
    \item Nhà thầu chỉ được phép xuất kẽm chạy máy in hàng loạt sau khi có chữ ký duyệt niêm phong mẫu bản bông của Giám đốc dự án và Chủ biên tại ngày thứ 76;
    \item Nếu giao hàng trễ sau ngày thứ 84 mà không do sự kiện bất khả kháng, nhà thầu chịu mức phạt vi phạm tiến độ bằng 1\% giá trị hợp đồng cho mỗi ngày chậm trễ (khấu trừ trực tiếp vào đợt thanh toán cuối);
    \item Nhà thầu phải thực hiện nghĩa vụ in bù đổi mới trong vòng 48 giờ hoặc tái sản xuất lô hàng theo đúng các ngưỡng nghiệm thu mẫu thống kê 384 cuốn đã quy định tại Mục 6.4.2 (Chương 6).
\end{enumerate}

Về lộ trình thanh toán hợp đồng, khớp với lịch giải ngân tại Chương 5, Ban Quản lý dự án chuyển khoản tạm ứng 30\% giá trị hợp đồng (\textbf{28.200.000 đồng}) vào ngày thứ 76 ngay sau khi ký duyệt bản bông, và tất toán 70\% giá trị hợp đồng còn lại (\textbf{65.800.000 đồng}) sau khi hoàn thành kiểm định nghiệm thu lô hàng.

\subsubsection{10.4.2. Kiểm soát thay đổi hợp đồng}

Trong quá trình thực hiện hợp đồng, mọi điều chỉnh về thông số kỹ thuật hoặc lịch giao nhận (nếu phát sinh) đều phải được lập thành Phụ lục hợp đồng bằng văn bản có chữ ký của đại diện pháp luật hai bên; nhân viên giám sát kỹ thuật và thợ in tại xưởng không được tự ý thỏa thuận thay đổi quy cách giấy, mực hoặc bỏ qua bước ký duyệt bản bông.

\subsection{10.5. Xác minh sản phẩm và thanh lý hợp đồng}

Vào ngày thứ 84 (24/05/2026), nhà thầu bàn giao đủ 2.000 cuốn tài liệu trong 40 thùng giấy niêm phong tại kho tiếp nhận. Tiểu ban Kiểm định chất lượng thực hiện rút mẫu phân tầng 384 cuốn theo quy trình tại Chương 6; kết quả đo đạc thực tế ghi nhận toàn bộ 384/384 cuốn mẫu đều đạt chuẩn kích thước 17 $\times$ 24 cm, giấy đúng định lượng, màu sắc khớp với bản bông đối chứng và keo gáy chắc chắn (0 cuốn lỗi), đủ điều kiện ký Biên bản xác minh và giao nhận hàng hóa.

Vào ngày thứ 88 (28/05/2026), sau khi nhà thầu bàn giao bộ dữ liệu bình bản xuất kẽm gốc và Hóa đơn giá trị gia tăng trị giá 94.000.000 đồng, hai bên ký Biên bản thanh lý Hợp đồng kinh tế số 08/TLHĐ-IN. Bộ phận kế toán dự án thực hiện lệnh chuyển khoản thanh toán nốt \textbf{65.800.000 đồng} (70\% giá trị hợp đồng còn lại) cho nhà thầu và chuyển toàn bộ chứng từ vào 03 bộ hồ sơ lưu trữ dự án trước khi kết thúc dự án vào ngày thứ 90.


\newpage

% 14. Phần Kết luận
\section{PHẦN KẾT LUẬN}

\subsection{1. Đánh giá tổng quan về tính khả thi và sự đồng bộ của kế hoạch dự án}

Bản Kế hoạch Quản lý Dự án đối với đề tài \textit{``Dự án xây dựng bộ tài liệu tư vấn phân luồng và định hướng nghề nghiệp cho học sinh cuối cấp trung học cơ sở''} đã thiết lập một phương án can thiệp đồng bộ về phạm vi, tiến độ, ngân sách, chất lượng, nhân sự, rủi ro, truyền thông và mua sắm nhằm giải quyết trực tiếp nhu cầu học liệu hướng nghiệp cho học sinh lớp 9.

Mười chương nội dung của bản kế hoạch được liên kết chặt chẽ theo một mạch thông số thống nhất giữa các hợp phần quản trị:

Từ kết quả phân tích Cây vấn đề và Cây mục tiêu (Chương 1), ma trận lựa chọn chiến lược đã luận chứng định lượng lý do lựa chọn phương án biên soạn và xuất bản ấn phẩm chuẩn hóa (đạt \textbf{8,55/10 điểm}). Trên cơ sở đó, Điều lệ dự án và cơ chế kiểm soát thay đổi tổng thể (Chương 2) xác lập khung pháp lý cho Ban Quản lý dự án triển khai đề án trong 90 ngày lịch (từ ngày 02/03/2026 đến ngày 30/05/2026) với ngân sách 185.000.000 đồng.

Tuyên bố phạm vi và Cấu trúc phân rã công việc (Chương 3) chia nhỏ dự án thành 5 giai đoạn và 20 gói công việc cụ thể (\texttt{1.1} đến \texttt{5.4}) nhằm tạo ra 2.000 cuốn tài liệu khổ 17 $\times$ 24 cm (gồm 80 trang ruột và 4 trang bìa ngoài) kèm bộ tệp tin chế bản điện tử gốc. Hai mươi gói công việc này được tích hợp thành 10 công việc mạng (A đến J) tại Chương 4, qua đó xác định chính xác đường găng gồm 9 công việc có tổng độ dài 90 ngày và phát hiện khoản dự trữ 20 ngày của công việc mua sắm in ấn (D) để chủ động điều hòa nguồn lực.

Hệ thống 20 gói công việc trong Từ điển phân rã công việc tiếp tục được ánh xạ trực tiếp sang Bảng dự toán chi phí từ dưới lên tại Chương 5 (áp dụng định mức của Thông tư số 51/2019/TT-BTC) và gắn với đường mức chi phí cơ sở qua 5 đợt giải ngân theo mốc nghiệm thu. Tại Chương 6, công tác kiểm soát chất lượng được định lượng hóa bằng bài toán lấy mẫu thống kê $n = 384$ cuốn (độ tin cậy 95\%, sai số chấp nhận được 5\%) cùng bảng phân tích 5 nhóm chi phí chất lượng trị giá \textbf{41.200.000 đồng} (chiếm \textbf{22,27\%} tổng ngân sách).

Ở góc độ tổ chức thực thi, Chương 7 thiết lập cơ cấu ma trận cân bằng gồm 4 tiểu ban (với 4 nhóm tác giả chuyên môn) và ma trận phân bổ trách nhiệm bảo đảm mỗi công việc chỉ có duy nhất một cá nhân giữ quyền phê duyệt (\texttt{A}). Chương 8 định lượng hóa bài toán lựa chọn nhà thầu in bằng mô hình Cây quyết định ($\text{EMV}_1 = +16.000.000\text{ đồng} > \text{EMV}_2 = -10.000.000\text{ đồng}$), làm tiền đề trực tiếp cho Chương 10 tổ chức chào hàng cạnh tranh rút gọn gói thầu in ấn (dự toán 95.000.000 đồng, trúng thầu 94.000.000 đồng theo Luật Đấu thầu số 22/2023/QH15). Toàn bộ luồng thông tin phối hợp, báo cáo tích hợp tại mốc ngày thứ 60 và thủ tục đóng gói \textbf{03 bộ hồ sơ lưu trữ gốc} được chuẩn hóa trọn vẹn tại Chương 9.

\subsection{2. Khuyến nghị khai thác sản phẩm sau bàn giao}

Để 2.000 cuốn tài liệu và bộ tệp tin chế bản điện tử gốc phát huy hiệu quả sau khi dự án kết thúc vào ngày 30 tháng 5 năm 2026, Ban Quản lý dự án đề xuất ba nhóm khuyến nghị đối với các đơn vị tiếp nhận:

\textbf{Đối với Phòng Giáo dục và Đào tạo:} Ban hành văn bản hướng dẫn các trường trung học cơ sở sử dụng bộ tài liệu trong các tiết chủ đề hướng nghiệp thuộc môn Hoạt động trải nghiệm, hướng nghiệp lớp 9; đồng thời đăng tải phiên bản tài liệu điện tử (từ bộ tệp tin chế bản gốc đã nghiệm thu) lên Cổng thông tin điện tử của Phòng để học sinh và phụ huynh toàn địa bàn có thể tải về tra cứu miễn phí.

\textbf{Đối với Ban Giám hiệu và giáo viên năm trường trung học cơ sở thụ hưởng:} Phân bổ ngay 2.000 cuốn tài liệu về các lớp 9 và thư viện trường; hướng dẫn học sinh thực hiện bài trắc nghiệm Holland 36 câu hỏi tại Phần 1 và khuyến khích các em mang cuốn tài liệu về nhà để cùng cha mẹ điền phiếu đối thoại chọn trường tại Phần 4 trước thời điểm đăng ký nguyện vọng thi tuyển sinh vào lớp 10.

\textbf{Đối với các cơ sở giáo dục nghề nghiệp trên địa bàn:} Phối hợp định kỳ hàng năm với Phòng Giáo dục và Đào tạo để cung cấp số liệu cập nhật về chỉ tiêu tuyển sinh hệ chín cộng, ngành nghề đào tạo mới và cam kết việc làm sau tốt nghiệp, làm cơ sở cho việc cập nhật tờ phụ lục thông tin hướng nghiệp trong các năm học tiếp theo.


\newpage

% 15. Danh mục Tài liệu tham khảo
\renewcommand{\refname}{DANH MỤC TÀI LIỆU THAM KHẢO}
\addcontentsline{toc}{section}{DANH MỤC TÀI LIỆU THAM KHẢO}

\setlength{\bibhang}{1.27cm}
\setlength{\bibsep}{6pt}

\begin{thebibliography}{99}

\bibitem[Bộ Chính trị(2024)]{bo_chinh_tri_2024}
Bộ Chính trị. (2024). \textit{Kết luận số 91-KL/TW ngày 12 tháng 8 năm 2024 về tiếp tục thực hiện Nghị quyết số 29-NQ/TW của Ban Chấp hành Trung ương Đảng khóa XI về đổi mới căn bản, toàn diện giáo dục và đào tạo}.

\bibitem[Bộ Giáo dục và Đào tạo(2017)]{bgd_tt33_2017}
Bộ Giáo dục và Đào tạo. (2017). \textit{Thông tư số 33/2017/TT-BGDĐT ngày 18 tháng 12 năm 2017 hướng dẫn công tác tư vấn tâm lý cho học sinh trong trường phổ thông}.

\bibitem[Bộ Giáo dục và Đào tạo(2018)]{bgd_tt32_2018}
Bộ Giáo dục và Đào tạo. (2018). \textit{Thông tư số 32/2018/TT-BGDĐT ngày 26 tháng 12 năm 2018 ban hành Chương trình Giáo dục phổ thông} (được sửa đổi, bổ sung bởi Thông tư số 13/2022/TT-BGDĐT ngày 03 tháng 8 năm 2022 của Bộ Giáo dục và Đào tạo).

\bibitem[Bộ Giáo dục và Đào tạo(2023)]{bgd_tt20_2023}
Bộ Giáo dục và Đào tạo. (2023). \textit{Thông tư số 20/2023/TT-BGDĐT ngày 30 tháng 10 năm 2023 hướng dẫn về vị trí việc làm, cơ cấu viên chức theo chức danh nghề nghiệp và định mức số lượng người làm việc trong các cơ sở giáo dục phổ thông công lập}.

\bibitem[Bộ Tài chính(2018)]{btc_tt76_2018}
Bộ Tài chính. (2018). \textit{Thông tư số 76/2018/TT-BTC ngày 17 tháng 8 năm 2018 hướng dẫn nội dung, mức chi xây dựng chương trình đào tạo, biên soạn giáo trình môn học đối với giáo dục đại học, giáo dục nghề nghiệp}.

\bibitem[Bộ Tài chính(2019)]{btc_tt51_2019}
Bộ Tài chính. (2019). \textit{Thông tư số 51/2019/TT-BTC ngày 15 tháng 8 năm 2019 hướng dẫn xây dựng dự toán, thanh quyết toán kinh phí biên soạn, chỉnh lý và thẩm định tài liệu giáo dục của địa phương}.

\bibitem[Bùi Thị Thúy Hằng(2020)]{bui_thi_thuy_hang_2020}
Bùi Thị Thúy Hằng. (2020). Ứng dụng phương pháp dạy học theo dự án và các kỹ thuật quản lý dự án trong đổi mới giáo dục kỹ thuật -- nghề nghiệp. \textit{Tạp chí Khoa học và Công nghệ}, Đại học Bách khoa Hà Nội, (142), 55--61.

\bibitem[Bùi Thị Thúy Hằng(2026)]{bui_thi_thuy_hang_2026}
Bùi Thị Thúy Hằng. (2026). \textit{Hệ thống bài giảng học phần Quản lý dự án giáo dục (11 chương)}. Chương trình Cao học Lý luận và Phương pháp giảng dạy, Đại học Bách khoa Hà Nội.

\bibitem[Chính phủ(2013)]{cp_nd195_2013}
Chính phủ. (2013). \textit{Nghị định số 195/2013/NĐ-CP ngày 21 tháng 11 năm 2013 quy định chi tiết một số điều và biện pháp thi hành Luật Xuất bản} (được sửa đổi, bổ sung bởi Nghị định số 150/2018/NĐ-CP ngày 07 tháng 11 năm 2018 của Chính phủ).

\bibitem[Chính phủ(2021)]{cp_nd81_2021}
Chính phủ. (2021). \textit{Nghị định số 81/2021/NĐ-CP ngày 27 tháng 8 năm 2021 quy định về cơ chế thu, quản lý học phí đối với cơ sở giáo dục thuộc hệ thống giáo dục quốc dân và chính sách miễn, giảm học phí, hỗ trợ chi phí học tập} (được sửa đổi, bổ sung bởi Nghị định số 97/2023/NĐ-CP ngày 31 tháng 12 năm 2023 của Chính phủ).

\bibitem[Chính phủ(2024)]{cp_nd24_2024}
Chính phủ. (2024). \textit{Nghị định số 24/2024/NĐ-CP ngày 27 tháng 2 năm 2024 quy định chi tiết một số điều và biện pháp thi hành Luật Đấu thầu về lựa chọn nhà thầu}.

\bibitem[Đặng Bá Lãm(2019)]{dang_ba_lam_2019}
Đặng Bá Lãm. (2019). Phân luồng học sinh sau trung học cơ sở ở Việt Nam: Thực trạng, nguyên nhân và giải pháp chính sách. \textit{Tạp chí Khoa học Giáo dục Việt Nam}, Viện Khoa học Giáo dục Việt Nam, (163), 1--8.

\bibitem[Lê Thị Quỳnh Nga et~al.(2026)]{le_thi_quynh_nga_2026}
Lê Thị Quỳnh Nga, Nguyễn Thị Hương, \& Chu Cẩm Thơ. (2026). Quy trình định hướng nghề nghiệp cho học sinh trung học cơ sở đáp ứng yêu cầu Chương trình giáo dục phổ thông 2018. \textit{Tạp chí Giáo dục}, Bộ Giáo dục và Đào tạo, số tháng 2 năm 2026.

\bibitem[Nguyễn Đức Sơn et~al.(2026)]{nguyen_duc_son_2026}
Nguyễn Đức Sơn, Lê Minh Nguyệt, Đàm Thị Vân Anh, \& Dương Minh Huệ. (2026). Hoạt động hướng nghiệp và phân luồng học sinh: Nghiên cứu thực trạng tại các trường trung học cơ sở trên địa bàn Thành phố Hà Nội. \textit{Tạp chí Giáo dục}, 26(16), 93--99.

\bibitem[Nguyễn Thị Mai Hương(2022)]{nguyen_thi_mai_huong_2022}
Nguyễn Thị Mai Hương. (2022). Tư vấn tâm lý và định hướng phân luồng cho học sinh cuối cấp trung học cơ sở trước kỳ thi chuyển cấp. \textit{Tạp chí Quản lý Giáo dục}, 14(6), 72--79.

\bibitem[OECD(2021)]{oecd_2021}
Organisation for Economic Co-operation and Development. (2021). \textit{Career readiness in the pandemic and post-pandemic era: How schools can help youth transition into work and upper secondary education} (OECD Education Working Papers No. 241). OECD Publishing, Paris.

\bibitem[Phạm Mạnh Hà(2020)]{pham_manh_ha_2020}
Phạm Mạnh Hà. (2020). Ứng dụng lý thuyết mật mã Holland trong công tác tư vấn hướng nghiệp cho học sinh phổ thông Việt Nam: Khả năng thích ứng và chuẩn hóa công cụ. \textit{Tạp chí Tâm lý học}, 8(257), 45--56.

\bibitem[Project Management Institute(2021)]{pmi_2021}
Project Management Institute. (2021). \textit{A guide to the project management body of knowledge (PMBOK guide)} (7th ed.). Project Management Institute, Newtown Square, Pennsylvania, United States of America.

\bibitem[Quốc hội(2012)]{qh_luat_xuat_ban_2012}
Quốc hội nước Cộng hòa xã hội chủ nghĩa Việt Nam. (2012). \textit{Luật Xuất bản số 19/2012/QH13}, ban hành ngày 20 tháng 11 năm 2012 (được sửa đổi, bổ sung bởi Luật số 35/2018/QH14).

\bibitem[Quốc hội(2019)]{qh_luat_giao_duc_2019}
Quốc hội nước Cộng hòa xã hội chủ nghĩa Việt Nam. (2019). \textit{Luật Giáo dục số 43/2019/QH14}, ban hành ngày 14 tháng 6 năm 2019.

\bibitem[Quốc hội(2023)]{qh_luat_dau_thau_2023}
Quốc hội nước Cộng hòa xã hội chủ nghĩa Việt Nam. (2023). \textit{Luật Đấu thầu số 22/2023/QH15}, ban hành ngày 23 tháng 6 năm 2023.

\bibitem[Quốc hội(2024, 2025)]{qh_luat_dau_thau_sua_doi}
Quốc hội nước Cộng hòa xã hội chủ nghĩa Việt Nam. (2024, 2025). \textit{Luật số 57/2024/QH15} và \textit{Luật số 90/2025/QH15 sửa đổi, bổ sung một số điều của Luật Đấu thầu}.

\bibitem[Quốc hội(2025)]{qh_luat_giao_duc_sua_doi_2025}
Quốc hội nước Cộng hòa xã hội chủ nghĩa Việt Nam. (2025). \textit{Luật số 123/2025/QH15 sửa đổi, bổ sung một số điều của Luật Giáo dục}, ban hành ngày 10 tháng 12 năm 2025.

\bibitem[Thủ tướng Chính phủ(2018)]{ttcp_qd522_2018}
Thủ tướng Chính phủ. (2018). \textit{Quyết định số 522/QĐ-TTg ngày 14 tháng 5 năm 2018 phê duyệt Đề án Giáo dục hướng nghiệp và định hướng phân luồng học sinh trong giáo dục phổ thông giai đoạn 2018 -- 2025}.

\bibitem[Tổ chức Lao động Quốc tế(2020)]{ilo_2020}
Tổ chức Lao động Quốc tế tại Việt Nam. (2020). \textit{Chuyển tiếp từ trường học đến việc làm của thanh thiếu niên Việt Nam}. Nhà xuất bản Lao động, Hà Nội.

\bibitem[Trần Thị Hương \& Hoàng Mai Khanh(2021)]{tran_thi_huong_2021}
Trần Thị Hương, \& Hoàng Mai Khanh. (2021). Giáo dục hướng nghiệp cho học sinh trung học cơ sở trong bối cảnh thực hiện Chương trình giáo dục phổ thông mới. \textit{Tạp chí Khoa học Đại học Sư phạm Thành phố Hồ Chí Minh}, 18(5), 815--826.

\bibitem[UNESCO-UNEVOC(2021)]{unesco_unevoc_2021}
UNESCO International Centre for Technical and Vocational Education and Training. (2021). \textit{Career guidance in technical and vocational education and training and secondary education: A global perspective on transitions to work}. Bonn, Germany.

\end{thebibliography}


\newpage

\end{document}

